\chapter{Cohomology of sheaves of modules}
\label{chap:Cohomology}

The site-theoretic conventions and standard cohomological facts are presented following
the Stacks Project \cite{stacks-project}; the projective-space calculation is cross-checked
against Hartshorne \cite{hartshorne-ag}.

Fix a commutative ring \(R\); for the genus of a curve, \(R\) will be the base field \(k\).
This chapter defines the cohomology of sheaves of \(R\)-modules on a small site as
\(\Ext\)-groups from the constant sheaf, and more generally the cohomology of an object of
the site as \(\Ext\)-groups from a free sheaf of \(R\)-modules; it identifies degree-zero
cohomology with (global) sections and proves the vanishing of degree-one cohomology of the
structure sheaf on an affine open of a scheme. It then develops the Mayer--Vietoris
\((0,1)\)-slice for a sheaf of \(R\)-modules on a site and specializes it to a scheme
covered by two affine opens, producing an explicit description of \(H^1(X,\struct{X/k})\)
as the cokernel of a map of section modules. It then runs the
two-lattice ladder of Chapter~\ref{chap:Algebra} on that cokernel along a finite morphism
to the projective line, proving the finiteness of \(H^1(C, \struct C)\) --- so that the
genus of Definition~\ref{def:genus} is the dimension of a finite-dimensional vector space.
The two closing sections glue a twisted sheaf out of a unit cocycle on a two-chart cover
and prove that the two-cover \(H^1\) of the relative curve --- for the structure sheaf,
and termwise for the twisted sheaf --- commutes with arbitrary change of the affine test
ring.

\section{Cohomology on a site}

Let \((\mathcal C, J)\) be a small site: a small category \(\mathcal C\) with a Grothendieck
topology \(J\). The category of sheaves of \(R\)-modules on \((\mathcal C, J)\) is a
Grothendieck abelian category; in particular it has enough injectives, and all
\(\Ext\)-groups exist and carry \(R\)-module structures.

\begin{definition}[The constant sheaf]
  \label{def:constModuleSheaf}
  \dcref{custom}
  \lean{CategoryTheory.Sheaf.constModuleSheaf}
  \leanok
  The constant sheaf \(R_J\) of \(R\)-modules on \((\mathcal C, J)\) with value \(R\): the
  sheafification of the constant presheaf with value \(R\).
\end{definition}

\begin{definition}[Sheaf cohomology]
  \label{def:HModule}
  \dcref{custom}
  \lean{CategoryTheory.Sheaf.HModule}
  \uses{def:constModuleSheaf}
  \leanok
  For a sheaf \(F\) of \(R\)-modules on \((\mathcal C, J)\) and \(n \ge 0\), the degree-\(n\)
  cohomology of \(F\) is the \(R\)-module
  \[
    H^n(F) \;=\; \Ext^n\bigl(R_J,\, F\bigr),
  \]
  the \(\Ext\)-group computed in the category of sheaves of \(R\)-modules.
\end{definition}

\begin{definition}[Functoriality]
  \label{def:HModule_map}
  \dcref{custom}
  \lean{CategoryTheory.Sheaf.HModule.map}
  \uses{def:HModule}
  \leanok
  A morphism of sheaves \(f : F \to G\) induces an \(R\)-linear map
  \(H^n(F) \to H^n(G)\) for every \(n\), by covariant functoriality of \(\Ext\) in the second
  argument. The identity induces the identity, and a composite induces the composite.
\end{definition}
\begin{proof}
  \leanok
  Functoriality of \(\Ext^n(R_J, -)\); the \(R\)-linearity comes from the \(R\)-linear
  structure of the sheaf category.
\end{proof}

\begin{lemma}[No higher cohomology on injectives]
  \label{lem:HModule_injective_vanishing}
  \dcref{custom}
  \lean{CategoryTheory.Sheaf.HModule.instSubsingletonHAddNatOfNat}
  \uses{def:HModule}
  \leanok
  If \(F\) is an injective object of the category of sheaves of \(R\)-modules, then
  \(H^{n+1}(F) = 0\) for every \(n \ge 0\).
\end{lemma}
\begin{proof}
  \leanok
  Positive-degree \(\Ext\)-groups into an injective object vanish.
\end{proof}

\begin{definition}[The adjunction, linearly]
  \label{def:constAdjHomEquiv}
  \dcref{custom}
  \lean{CategoryTheory.Sheaf.constantSheafAdjHomLinearEquiv}
  \leanok
  Assume the category \(\mathcal C\) has a terminal object \(\top\). For every \(R\)-module
  \(M\) and every sheaf \(F\) of \(R\)-modules there is an \(R\)-linear equivalence
  \[
    \Hom\bigl(\underline{M},\, F\bigr) \;\cong\; \Hom_R\bigl(M,\, F(\top)\bigr),
  \]
  where \(\underline{M}\) denotes the constant sheaf with value \(M\); it is the hom-bijection
  of the adjunction between the constant-sheaf functor and evaluation at \(\top\).
\end{definition}
\begin{proof}
  \leanok
  The constant-sheaf functor is left adjoint to evaluation at the terminal object. The
  bijection sends \(\varphi\) to the composite of the unit \(M \to \underline{M}(\top)\) with
  \(\varphi_\top\); both this assignment and its inverse are additive and commute with the
  \(R\)-action.
\end{proof}

\begin{definition}[Morphisms out of the constant sheaf are sections]
  \label{def:constHomEquiv}
  \dcref{custom}
  \lean{CategoryTheory.Sheaf.constModuleSheafHomEquiv}
  \uses{def:constModuleSheaf, def:constAdjHomEquiv}
  \leanok
  The \(R\)-linear equivalence
  \(\Hom(R_J, F) \cong F(\top)\): specialize \ref{def:constAdjHomEquiv} to \(M = R\) and
  compose with the canonical evaluation \(\Hom_R(R, N) \cong N\), \(u \mapsto u(1)\).
\end{definition}

\begin{lemma}[Naturality]
  \label{lem:constHomEquiv_naturality}
  \dcref{custom}
  \lean{CategoryTheory.Sheaf.constModuleSheafHomEquiv_naturality}
  \uses{def:constHomEquiv}
  \leanok
  The equivalence of \ref{def:constHomEquiv} is natural in the sheaf: for
  \(\varphi : R_J \to F\) and \(f : F \to G\), the section corresponding to
  \(f \circ \varphi\) is \(f_\top\) applied to the section corresponding to \(\varphi\).
\end{lemma}
\begin{proof}
  \leanok
  Naturality of the adjunction hom-bijection in the second variable, and naturality of the
  evaluation \(\Hom_R(R, N) \cong N\).
\end{proof}

\begin{definition}[Degree-zero cohomology is global sections]
  \label{def:HModule_linearEquiv0}
  \dcref{custom}
  \lean{CategoryTheory.Sheaf.HModule.linearEquiv₀}
  \uses{def:HModule, def:constHomEquiv}
  \leanok
  The \(R\)-linear equivalence \(H^0(F) \cong F(\top)\): the canonical identification
  \(\Ext^0(R_J, F) \cong \Hom(R_J, F)\) followed by \ref{def:constHomEquiv}.
\end{definition}

\begin{lemma}[Naturality in degree zero]
  \label{lem:HModule_linearEquiv0_naturality}
  \dcref{custom}
  \lean{CategoryTheory.Sheaf.HModule.linearEquiv₀_naturality}
  \uses{def:HModule_linearEquiv0, def:HModule_map}
  \leanok
  The equivalence \(H^0(F) \cong F(\top)\) of \ref{def:HModule_linearEquiv0} is natural: for
  \(f : F \to G\), applying \(f_\top\) to the section corresponding to \(x \in H^0(F)\)
  yields the section corresponding to the image of \(x\) under \(H^0(f)\).
\end{lemma}
\begin{proof}
  \leanok
  \uses{lem:constHomEquiv_naturality}
  The identification \(\Ext^0 \cong \Hom\) turns composition with the degree-zero class of
  \(f\) into postcomposition with \(f\); conclude by \ref{lem:constHomEquiv_naturality}.
\end{proof}

\section{Cohomology of an object of the site}

For an object \(U\) of \(\mathcal C\), this section defines the cohomology of \(U\) with
coefficients in a sheaf \(F\) of \(R\)-modules as \(\Ext\)-groups from a \emph{free} sheaf
of \(R\)-modules on \(U\), generalizing the previous section (where the implicit object was
the terminal object \(\top\)). This is the coefficient system needed to run the
Mayer--Vietoris sequence of Section~\ref{sec:mv-in-text} \(R\)-linearly, and it recovers the
cohomology of the site when \(U\) is terminal (Definition~\ref{def:HModule_linearEquivHModule'}
below).

\begin{definition}[The free sheaf of \(R\)-modules on an object of the site]
  \label{def:freeModuleSheaf}
  \dcref{custom}
  \lean{CategoryTheory.Sheaf.freeModuleSheaf}
  \leanok
  For an object \(U\) of \(\mathcal C\), the free sheaf of \(R\)-modules on \(U\), written
  \(R[U]\): the sheafification of the presheaf \(V \mapsto\) the free \(R\)-module on the
  set \(\Hom(V, U)\), with restriction maps induced by precomposition.
\end{definition}

\begin{definition}[Functoriality in the object of the site]
  \label{def:freeModuleSheafMap}
  \dcref{custom}
  \lean{CategoryTheory.Sheaf.freeModuleSheafMap}
  \uses{def:freeModuleSheaf}
  \leanok
  A morphism \(i : U \to V\) of \(\mathcal C\) induces a morphism of sheaves of
  \(R\)-modules \(R[i] : R[U] \to R[V]\): apply the free-\(R\)-module-on-a-set functor and
  sheafification to the induced map of representable presheaves
  \(\Hom(-,U) \to \Hom(-,V)\), \(g \mapsto i \circ g\).
\end{definition}

\begin{definition}[Free sheaves of modules represent sections]
  \label{def:freeModuleSheafHomEquiv}
  \dcref{custom}
  \lean{CategoryTheory.Sheaf.freeModuleSheafHomEquiv}
  \uses{def:freeModuleSheaf}
  \leanok
  For every object \(U\) of \(\mathcal C\) and every sheaf \(F\) of \(R\)-modules, an
  \(R\)-linear equivalence
  \[
    \Hom\bigl(R[U],\, F\bigr) \;\cong\; F(U).
  \]
\end{definition}
\begin{proof}
  \leanok
  Write \(y_U = \Hom(-,U)\) for the presheaf represented by \(U\), so that \(R[U]\) is the
  sheafification of the presheaf of \(R\)-modules obtained by applying the free-\(R\)-module
  functor to \(y_U\) objectwise. Compose three natural bijections:
  \begin{itemize}
    \item sheafification is left adjoint to the inclusion of sheaves into presheaves, giving
      \(\Hom(R[U],F) \cong \Hom_{\mathrm{PSh}}(y_U \otimes R,\, F)\), naturally in \(F\);
    \item on each object \(V\), the free-\(R\)-module functor is left adjoint to the
      forgetful functor from \(R\)-modules to sets, giving
      \(\Hom_{\mathrm{PSh}}(y_U \otimes R,\, F) \cong \Hom_{\mathrm{PSh}(\mathrm{Set})}(y_U,\, F)\)
      (\(F\) viewed as a presheaf of sets), naturally in \(F\);
    \item the Yoneda lemma identifies \(\Hom_{\mathrm{PSh}(\mathrm{Set})}(y_U, F) \cong F(U)\).
  \end{itemize}
  Each bijection is additive and \(R\)-linear in \(F\) (the first two are hom-sets of
  \(R\)-linear/additive adjunctions; the last carries the \(R\)-module structure of \(F(U)\)
  tautologically, evaluating a natural transformation at the identity of \(U\)), so the
  composite is an \(R\)-linear equivalence.
\end{proof}

\begin{lemma}[Naturality of the section equivalence]
  \label{lem:freeModuleSheafHomEquiv_naturality}
  \dcref{custom}
  \lean{CategoryTheory.Sheaf.freeModuleSheafHomEquiv_comp,
    CategoryTheory.Sheaf.freeModuleSheafHomEquiv_map}
  \uses{def:freeModuleSheafHomEquiv, def:freeModuleSheafMap}
  \leanok
  The equivalence of \ref{def:freeModuleSheafHomEquiv} is natural in both variables:
  \begin{enumerate}
    \item for \(\varphi : R[U] \to F\) and \(f : F \to G\), the section of \(G(U)\)
      corresponding to \(f \circ \varphi\) is \(f_U\) applied to the section of \(F(U)\)
      corresponding to \(\varphi\);
    \item for \(i : U \to V\) and \(\varphi : R[V] \to F\), the section of \(F(U)\)
      corresponding to the composite \(R[U] \xrightarrow{R[i]} R[V] \xrightarrow{\varphi} F\)
      is the restriction along \(i\) of the section of \(F(V)\) corresponding to \(\varphi\).
  \end{enumerate}
\end{lemma}
\begin{proof}
  \leanok
  Both are naturality of the composite of adjunction bijections used in
  \ref{def:freeModuleSheafHomEquiv}: (1) is naturality of each of the three bijections in
  the target sheaf \(F\); (2) is naturality of the Yoneda bijection
  \(\Hom_{\mathrm{PSh}(\mathrm{Set})}(y_U,-) \cong \mathrm{ev}_U\) in \(U\) (precomposition
  with \(y_i : y_U \to y_V\) corresponds to restriction along \(i\)), transported through the
  other two (fixed) bijections.
\end{proof}

\begin{definition}[Cohomology of an object of the site]
  \label{def:HModule'}
  \dcref{custom}
  \lean{CategoryTheory.Sheaf.HModule'}
  \uses{def:freeModuleSheaf}
  \leanok
  For an object \(U\) of \(\mathcal C\), a sheaf \(F\) of \(R\)-modules and \(n \ge 0\), the
  degree-\(n\) cohomology of \(U\) with coefficients in \(F\) is the \(R\)-module
  \[
    H'^n(U, F) \;=\; \Ext^n\bigl(R[U],\, F\bigr).
  \]
\end{definition}

\begin{definition}[Degree zero is sections]
  \label{def:HModule'_linearEquiv0}
  \dcref{custom}
  \lean{CategoryTheory.Sheaf.HModule'.linearEquiv₀}
  \uses{def:HModule', def:freeModuleSheafHomEquiv}
  \leanok
  The \(R\)-linear equivalence \(H'^0(U, F) \cong F(U)\): the canonical identification
  \(\Ext^0(R[U], F) \cong \Hom(R[U], F)\) followed by \ref{def:freeModuleSheafHomEquiv}.
\end{definition}

\begin{definition}[Restriction]
  \label{def:HModule'_res}
  \dcref{custom}
  \lean{CategoryTheory.Sheaf.HModule'.res}
  \uses{def:HModule', def:freeModuleSheafMap}
  \leanok
  For \(i : U \to V\) in \(\mathcal C\), restriction along \(i\) is the \(R\)-linear map
  \(H'^n(V, F) \to H'^n(U, F)\), \(x \mapsto x \circ R[i]\): precomposition with
  \(R[i] : R[U] \to R[V]\) (contravariant functoriality of \(\Ext^n(-, F)\)).
\end{definition}

\begin{definition}[Functoriality in the coefficient sheaf]
  \label{def:HModule'_mapCoeff}
  \dcref{custom}
  \lean{CategoryTheory.Sheaf.HModule'.mapCoeff,
    CategoryTheory.Sheaf.HModule'.mapCoeff_id_apply,
    CategoryTheory.Sheaf.HModule'.mapCoeff_comp_apply}
  \uses{def:HModule'}
  \leanok
  A morphism of sheaves of \(R\)-modules \(f : F \to G\) induces, for every object \(U\) of
  \(\mathcal C\) and every \(n\), an \(R\)-linear map \(H'^n(U, F) \to H'^n(U, G)\),
  \(x \mapsto x \cdot f\): composition (Yoneda product) with the degree-zero class of \(f\), i.e.\
  covariant functoriality of \(\Ext^n(R[U], -)\) in the coefficient variable. This is the
  coefficient-variable analogue of the restriction \ref{def:HModule'_res} (functoriality in the
  object of the site). The identity induces the identity, and a composite \(g \circ f\) induces the
  composite of the induced maps.
\end{definition}
\begin{proof}
  \leanok
  Postcomposition with the degree-zero \(\Ext\)-class of \(f\) is \(R\)-linear; it sends the
  degree-zero class of the identity to the identity operation, and by associativity of the Yoneda
  product it turns the class of a composite into the composite of the two postcompositions.
\end{proof}

\begin{definition}[Coefficient-isomorphism transport for cohomology of an object]
  \label{def:HModule'_congrCoeff}
  \dcref{custom}
  \lean{CategoryTheory.Sheaf.HModule'.congrCoeff}
  \uses{def:HModule'_mapCoeff}
  \leanok
  An isomorphism \(\alpha : F \cong G\) of coefficient sheaves of \(R\)-modules induces, for every
  object \(U\) and every \(n\), an \(R\)-linear isomorphism \(H'^n(U, F) \cong H'^n(U, G)\).
\end{definition}
\begin{proof}
  \leanok
  Take the induced maps of \ref{def:HModule'_mapCoeff} for \(\alpha\) and for \(\alpha^{-1}\). By
  functoriality their two composites are the maps induced by \(\alpha^{-1}\circ\alpha = \id_F\) and
  \(\alpha\circ\alpha^{-1} = \id_G\), which are the identity operations; hence the two induced maps
  are mutually inverse \(R\)-linear isomorphisms.
\end{proof}

\begin{lemma}[Transport of vanishing along a coefficient isomorphism]
  \label{lem:HModule'_subsingleton_of_congrCoeff}
  \dcref{custom}
  \lean{CategoryTheory.Sheaf.HModule'.subsingleton_of_congrCoeff}
  \uses{def:HModule'_congrCoeff}
  \leanok
  If \(F \cong G\) as coefficient sheaves of \(R\)-modules and \(H'^n(U, G) = 0\), then
  \(H'^n(U, F) = 0\).
\end{lemma}
\begin{proof}
  \leanok
  The \(R\)-linear isomorphism \(H'^n(U, F) \cong H'^n(U, G)\) of \ref{def:HModule'_congrCoeff}
  identifies \(H'^n(U, F)\) with the zero module \(H'^n(U, G)\).
\end{proof}

\begin{definition}[The free sheaf on a terminal object is the constant sheaf]
  \label{def:freeModuleSheafIsoConstModuleSheaf}
  \dcref{custom}
  \lean{CategoryTheory.Sheaf.freeModuleSheafIsoConstModuleSheaf}
  \uses{def:freeModuleSheaf, def:constModuleSheaf, def:freeModuleSheafHomEquiv,
    def:constHomEquiv}
  \leanok
  If \(T\) is a terminal object of \(\mathcal C\), a canonical isomorphism of sheaves of
  \(R\)-modules \(R[T] \cong R_J\).
\end{definition}
\begin{proof}
  \leanok
  Both sides corepresent the same functor: \(\Hom(R[T], -) \cong \mathrm{ev}_T\) by
  \ref{def:freeModuleSheafHomEquiv}, and \(\Hom(R_J, -) \cong \mathrm{ev}_T\) by
  \ref{def:constHomEquiv} (valid since \(T\) is terminal), both naturally in the sheaf
  argument. Two objects corepresenting the same functor via natural equivalences are
  canonically isomorphic: the isomorphism corresponds, under each equivalence, to the section
  of the other object at \(T\) that corresponds to its own identity morphism, and the two
  composites are the identity by naturality applied to the identity morphism of each side.
\end{proof}

\begin{definition}[Cohomology of the site is cohomology of a terminal object]
  \label{def:HModule_linearEquivHModule'}
  \dcref{custom}
  \lean{CategoryTheory.Sheaf.HModule.linearEquivHModule'}
  \uses{def:HModule, def:HModule', def:freeModuleSheafIsoConstModuleSheaf}
  \leanok
  For \(T\) terminal in \(\mathcal C\) and every \(n\), an \(R\)-linear equivalence
  \(H^n(F) \cong H'^n(T, F)\).
\end{definition}
\begin{proof}
  \leanok
  Apply the (contravariant, \(R\)-linear) functoriality of \(\Ext^n(-, F)\) to the
  isomorphism \(R[T] \cong R_J\) of \ref{def:freeModuleSheafIsoConstModuleSheaf}: an
  isomorphism of the first \(\Ext\)-argument induces, by composition (Yoneda product) with
  its degree-zero class, a natural \(R\)-linear isomorphism on \(\Ext\)-groups in every
  degree, with inverse given by composition with the inverse isomorphism.
\end{proof}

\section{The structure sheaf as a sheaf of modules}

Let \(k\) be a commutative ring and let \(X\) be a scheme over \(\Spec k\). The relevant site
is the \emph{small Zariski site} of \(X\): the partially ordered set of open subsets of
\(X\), with coverings the open coverings. The whole space is a terminal object of this site.

\begin{definition}[The algebra map on sections]
  \label{def:overAlgebraMap}
  \dcref{custom}
  \lean{AlgebraicGeometry.Scheme.overAlgebraMap}
  \leanok
  For every open \(U \subseteq X\), the canonical ring homomorphism \(k \to \Gamma(X, U)\):
  the map induced on global sections by the structure morphism, followed by restriction to
  \(U\). It commutes with the restriction maps of the structure sheaf; in particular each
  \(\Gamma(X, U)\) is a \(k\)-algebra --- hence a \(k\)-module --- and all restriction maps
  are \(k\)-linear.
\end{definition}
\begin{proof}
  \leanok
  Restriction maps are ring homomorphisms and commute with restriction from the global
  sections; the maps from \(k\) are fixed composites of such maps, so they commute with
  further restriction.
\end{proof}

\begin{definition}[The structure sheaf as a sheaf of \(k\)-modules]
  \label{def:moduleKSheaf}
  \dcref{custom}
  \lean{AlgebraicGeometry.Scheme.moduleKSheaf}
  \uses{def:overAlgebraMap}
  \leanok
  The structure sheaf of \(X\), viewed as a sheaf \(\struct{X/k}\) of \(k\)-modules on the
  small Zariski site of \(X\): its sections over an open \(U\) are \(\Gamma(X, U)\) with the
  \(k\)-module structure of \ref{def:overAlgebraMap}, and its restriction maps are those of
  the structure sheaf.
\end{definition}
\begin{proof}
  \leanok
  The restriction maps are \(k\)-linear by \ref{def:overAlgebraMap}, so this is a presheaf of
  \(k\)-modules. Its underlying presheaf of sets is that of the structure sheaf, which
  satisfies the sheaf condition; since limits of \(k\)-modules are computed on underlying
  sets, the presheaf of \(k\)-modules is a sheaf.
\end{proof}

\begin{definition}[Degree-zero cohomology of the structure sheaf]
  \label{def:moduleKSheafHZero}
  \dcref{custom}
  \lean{AlgebraicGeometry.Scheme.moduleKSheafHZero}
  \uses{def:moduleKSheaf, def:HModule_linearEquiv0}
  \leanok
  The \(k\)-linear equivalence \(H^0\bigl(X, \struct{X/k}\bigr) \cong \Gamma(X, \struct X)\).
\end{definition}
\begin{proof}
  \leanok
  The open \(X\) is a terminal object of the small Zariski site; apply
  \ref{def:HModule_linearEquiv0}.
\end{proof}

\section{\v{C}ech cobounding on affine opens}

For a scheme \(X\), a section \(g\) of the structure sheaf on an open \(V \subseteq X\) and
an open \(U \subseteq V\), we write \(D(g)\) for the basic open locus of \(V\) where \(g\) is
invertible, and \(g|_U\), \(s|_U\) for restrictions of sections.

\begin{lemma}[Denominator clearing on an affine open]
  \label{lem:denominator_clearing}
  \dcref{custom}
  \lean{AlgebraicGeometry.IsAffineOpen.exists_res_eq_pow_mul}
  \leanok
  Let \(X\) be a scheme, \(U \subseteq X\) an affine open, \(V \subseteq X\) an open with
  \(U \subseteq V\), and \(g \in \Gamma(X, V)\) a section of the structure sheaf on \(V\) (so
  that \(D(g) \subseteq V\) is its basic open). For every \(c \in \Gamma(X, U \cap D(g))\)
  there exist \(N \ge 0\) and \(e \in \Gamma(X, U)\) such that
  \[
    e|_{U \cap D(g)} \;=\; \bigl(g|_{U}\bigr)|_{U \cap D(g)}^{N} \cdot c .
  \]
\end{lemma}
\begin{proof}
  \leanok
  \(U \cap D(g)\) is the basic open of the affine \(U\) defined by \(g|_U\), and restriction
  identifies the sections over the basic open of an affine open with the localization:
  \(\Gamma(X, U \cap D(g)) = \Gamma(X, U)\bigl[1/(g|_U)\bigr]\) (quasi-coherence of the
  structure sheaf). Writing \(c\) as a fraction \(e / (g|_U)^N\) and clearing the denominator
  gives the claim.
\end{proof}

\begin{lemma}[Annihilation on an affine open]
  \label{lem:section_annihilation}
  \dcref{custom}
  \lean{AlgebraicGeometry.IsAffineOpen.exists_pow_mul_eq_zero_of_res_eq_zero}
  \leanok
  Let \(X\) be a scheme, \(U \subseteq X\) an affine open, \(V \subseteq X\) an open with
  \(U \subseteq V\), and \(g \in \Gamma(X, V)\) a section of the structure sheaf on \(V\) (so
  that \(D(g) \subseteq V\) is its basic open). If \(t \in \Gamma(X, U)\) restricts to zero
  on \(U \cap D(g)\), then \((g|_U)^{M} \cdot t = 0\) in \(\Gamma(X, U)\) for some \(M \ge 0\).
\end{lemma}
\begin{proof}
  \leanok
  As in \ref{lem:denominator_clearing}, \(\Gamma(X, U \cap D(g))\) is the localization of
  \(\Gamma(X, U)\) at the powers of \(g|_U\), and the kernel of the localization map consists
  of the elements annihilated by a power of \(g|_U\).
\end{proof}

\begin{theorem}[Degree-one \v{C}ech cobounding on an affine open]
  \label{thm:affine_open_cech_cobounding}
  \lean{AlgebraicGeometry.IsAffineOpen.exists_cech_cobounding}
  \leanok
  \dcref{custom}
  Let \(X\) be a scheme and \(U\) an affine open of \(X\). Let \(f_1, \dots, f_n \in
  \Gamma(X, U)\) with \(D(f_1) \cup \dots \cup D(f_n) = U\). Let
  \(a_{ij} \in \Gamma\bigl(X, D(f_i) \cap D(f_j)\bigr)\) be a \v{C}ech \(1\)-cocycle: for all
  \(i, j, l\),
  \[
    a_{il} \;=\; a_{ij} + a_{jl}
    \qquad \text{on } D(f_i) \cap D(f_j) \cap D(f_l).
  \]
  Then there are sections \(b_i \in \Gamma\bigl(X, D(f_i)\bigr)\) with
  \[
    a_{ij} \;=\; b_i - b_j \qquad \text{on } D(f_i) \cap D(f_j)
  \]
  for all \(i, j\).
\end{theorem}
\begin{proof}
  \leanok
  \uses{lem:denominator_clearing, lem:section_annihilation}
  All restriction maps are suppressed from the notation. Note that each \(D(f_i)\) and each
  pairwise intersection \(D(f_i) \cap D(f_j) = D(f_i f_j)\) is an affine open of \(U\), hence
  of \(X\).

  \emph{Step 1: clear denominators with a uniform exponent.} For each pair \((i,j)\) apply
  \ref{lem:denominator_clearing} with affine open \(D(f_i)\), ambient open \(V = U\) and
  \(g = f_j\): there are \(N_{ij} \ge 0\) and \(e^0_{ij} \in \Gamma(X, D(f_i))\) with
  \(e^0_{ij} = f_j^{N_{ij}} a_{ij}\) on \(D(f_i) \cap D(f_j)\). Let \(N\) be the maximum of
  the \(N_{ij}\) and set \(e_{ij} = f_j^{\,N - N_{ij}}\, e^0_{ij}\), so that
  \(e_{ij} = f_j^{\,N} a_{ij}\) on \(D(f_i) \cap D(f_j)\).

  \emph{Step 2: kill the defects.} For \(i, j, l\) define, on \(D(f_i) \cap D(f_j)\),
  \[
    t_{ijl} \;=\; e_{il} - e_{jl} - f_l^{\,N} a_{ij}.
  \]
  On the triple intersection \(D(f_i) \cap D(f_j) \cap D(f_l)\) we have
  \(e_{il} = f_l^N a_{il}\) and \(e_{jl} = f_l^N a_{jl}\), so the cocycle identity gives
  \(t_{ijl} = f_l^N (a_{il} - a_{jl} - a_{ij}) = 0\) there. Since
  \(D(f_i) \cap D(f_j)\) is an affine open (of \(U\), hence of \(X\)) and the triple
  intersection is its basic open defined by \(f_l\), \ref{lem:section_annihilation} (with
  ambient open \(V = U\)) yields \(M_{ijl} \ge 0\) with
  \(f_l^{\,M_{ijl}}\, t_{ijl} = 0\) on \(D(f_i) \cap D(f_j)\). Let \(M\) be the maximum of
  the \(M_{ijl}\); then \(f_l^{\,M} t_{ijl} = 0\) for all \(i,j,l\).

  \emph{Step 3: partition of unity.} Since the \(D(f_i)\) cover the affine open \(U\), the
  \(f_i\) generate the unit ideal of \(\Gamma(X, U)\); hence so do the powers
  \(f_i^{\,N+M}\), and we may write \(1 = \sum_l h_l\, f_l^{\,N+M}\) with
  \(h_l \in \Gamma(X, U)\).

  \emph{Step 4: assemble the cobounding.} Set
  \(b_i = \sum_l h_l\, f_l^{\,M}\, e_{il} \in \Gamma\bigl(X, D(f_i)\bigr)\). On
  \(D(f_i) \cap D(f_j)\),
  \[
    b_i - b_j
      = \sum_l h_l f_l^{\,M} \bigl(e_{il} - e_{jl}\bigr)
      = \sum_l h_l f_l^{\,M} \bigl(f_l^{\,N} a_{ij} + t_{ijl}\bigr)
      = \Bigl(\sum_l h_l f_l^{\,N+M}\Bigr) a_{ij}
      = a_{ij},
  \]
  using \(f_l^{\,M} t_{ijl} = 0\) in the third equality.
\end{proof}

\begin{theorem}[Degree-one \v{C}ech cobounding on an affine scheme]
  \label{thm:cech_cobounding}
  \dcref{custom}
  \lean{AlgebraicGeometry.Scheme.exists_cech_cobounding}
  \uses{thm:affine_open_cech_cobounding}
  \leanok
  Let \(X\) be an affine scheme and let \(f_1, \dots, f_n\) be global sections of the
  structure sheaf with \(D(f_1) \cup \dots \cup D(f_n) = X\). Let
  \(a_{ij} \in \Gamma\bigl(X, D(f_i) \cap D(f_j)\bigr)\) be a \v{C}ech \(1\)-cocycle as in
  \ref{thm:affine_open_cech_cobounding}. Then there are sections
  \(b_i \in \Gamma\bigl(X, D(f_i)\bigr)\) with \(a_{ij} = b_i - b_j\) on
  \(D(f_i) \cap D(f_j)\) for all \(i, j\).
\end{theorem}
\begin{proof}
  \leanok
  The special case \(U = X\) of \ref{thm:affine_open_cech_cobounding}: an affine scheme is
  an affine open of itself.
\end{proof}

\section{Vanishing of degree-one cohomology on affine opens}

\begin{lemma}[A vanishing criterion for \(\Ext^1\)]
  \label{lem:ext_one_vanishing_criterion}
  \dcref{custom}
  \lean{CategoryTheory.Abelian.Ext.subsingleton_one_of_injective_of_surjective}
  \leanok
  Let \(\mathcal A\) be an abelian category (with \(\Ext\)-groups), let
  \(\iota : A \to I\) be a monomorphism with \(I\) injective, and let \(L\) be an object. If
  every morphism \(L \to \coker \iota\) lifts along the projection
  \(\pi : I \to \coker \iota\), then \(\Ext^1(L, A) = 0\).
\end{lemma}
\begin{proof}
  \leanok
  Consider the short exact sequence \(0 \to A \to I \to \coker \iota \to 0\) and the induced
  long exact sequence
  \[
    \Hom(L, I) \xrightarrow{\;\pi_*\;} \Hom(L, \coker \iota)
      \xrightarrow{\;\delta\;} \Ext^1(L, A) \longrightarrow \Ext^1(L, I).
  \]
  Let \(z \in \Ext^1(L, A)\). Its image in \(\Ext^1(L, I)\) vanishes because \(I\) is
  injective, so by exactness \(z = \delta(x)\) for some \(x : L \to \coker \iota\). By
  hypothesis \(x = \pi \circ \psi\) for some \(\psi : L \to I\), and
  \(\delta \circ \pi_* = 0\) (composite of consecutive maps of the long exact sequence);
  hence \(z = \delta(\pi_*(\psi)) = 0\).
\end{proof}

\begin{lemma}[Sections are left exact]
  \label{lem:sections_left_exact}
  \dcref{custom}
  \lean{CategoryTheory.Sheaf.exists_app_eq_of_cokernelπ_app_eq_zero}
  \leanok
  Let \(\iota : F \to G\) be a monomorphism of sheaves of \(R\)-modules on a site and let
  \(U\) be an object of the site. If \(c \in G(U)\) maps to zero in \((\coker \iota)(U)\),
  then \(c = \iota_U(a)\) for a (unique) \(a \in F(U)\).
\end{lemma}
\begin{proof}
  \leanok
  In an abelian category, a monomorphism is a kernel of its cokernel. Evaluation of sheaves
  at \(U\) preserves limits: it is the composite of the inclusion of sheaves into presheaves
  (a right adjoint, sheafification being left adjoint to it) with evaluation of presheaves at
  \(U\) (limits of presheaves are computed objectwise). Hence
  \(F(U) \to G(U)\) is a kernel of \(G(U) \to (\coker \iota)(U)\), which is the claim.
\end{proof}

\begin{lemma}[Monomorphisms are injective on sections]
  \label{lem:app_injective_of_mono}
  \dcref{custom}
  \lean{CategoryTheory.Sheaf.app_injective_of_mono}
  \leanok
  A monomorphism of sheaves of \(R\)-modules on a site is injective on sections over every
  object of the site.
\end{lemma}
\begin{proof}
  \leanok
  Evaluation of sheaves at an object preserves limits (as in
  \ref{lem:sections_left_exact}), hence preserves monomorphisms, and a monomorphism of
  \(R\)-modules is an injective map.
\end{proof}

\begin{theorem}[Sections surjectivity on an affine open]
  \label{thm:cokernel_app_surjective_affine_open}
  \dcref{custom}
  \lean{AlgebraicGeometry.IsAffineOpen.cokernel_app_surjective}
  \uses{def:moduleKSheaf}
  \leanok
  Let \(k\) be a commutative ring, \(X\) a scheme over \(\Spec k\), \(U\) an affine open of
  \(X\), \(G\) a sheaf of \(k\)-modules on the small Zariski site of \(X\), and
  \(\iota : \struct{X/k} \to G\) a monomorphism. Then every section of
  \(\coker \iota\) over \(U\) lifts to a section of \(G\) over \(U\): the map
  \(G(U) \to (\coker \iota)(U)\) is surjective.
\end{theorem}
\begin{proof}
  \leanok
  \uses{lem:sections_left_exact, lem:app_injective_of_mono, thm:affine_open_cech_cobounding}
  Let \(q \in (\coker \iota)(U)\) and write \(\pi : G \to \coker \iota\) for the projection.

  (a) \(\pi\) is an epimorphism of sheaves, hence locally surjective on sections: every point
  of \(U\) has an open neighbourhood on which the restriction of \(q\) lifts to \(G\).
  Shrinking, and since the basic opens of \(U\) form a basis of \(U\) (as \(U\) is affine),
  every point of \(U\) lies in some basic open \(D(g) \subseteq U\) carrying a lift of
  \(q|_{D(g)}\).

  (b) \(U\) is quasi-compact (it is affine), so there are finitely many sections
  \(f_1, \dots, f_n \in \Gamma(X, U)\) with \(D(f_1) \cup \dots \cup D(f_n) = U\) and lifts
  \(s_i \in G(D(f_i))\) with \(\pi(s_i) = q|_{D(f_i)}\).

  (c) On \(D(f_i) \cap D(f_j)\) the difference \(s_i - s_j\) is killed by \(\pi\), so by
  \ref{lem:sections_left_exact} there is a unique
  \(a_{ij} \in \Gamma\bigl(X, D(f_i) \cap D(f_j)\bigr)\) with
  \(\iota(a_{ij}) = s_i - s_j\).

  (d) The \(a_{ij}\) form a \v{C}ech \(1\)-cocycle: on triple intersections,
  \(\iota(a_{il}) = s_i - s_l = (s_i - s_j) + (s_j - s_l) = \iota(a_{ij} + a_{jl})\), and
  \(\iota\) is injective on sections (\ref{lem:app_injective_of_mono}).

  (e) By \ref{thm:affine_open_cech_cobounding} (applied to the affine open \(U\)) there are
  \(b_i \in \Gamma\bigl(X, D(f_i)\bigr)\) with \(a_{ij} = b_i - b_j\) on overlaps. The
  corrected lifts \(s_i - \iota(b_i)\) then agree on all overlaps, so by the sheaf condition
  they glue to a section \(s \in G(U)\).

  (f) Finally \(\pi(s) = q\): on each \(D(f_i)\) we have
  \(\pi(s) = \pi(s_i) - \pi(\iota(b_i)) = q|_{D(f_i)}\) since \(\pi \circ \iota = 0\), and
  sections of the sheaf \(\coker \iota\) that agree on a cover of \(U\) agree on \(U\).
\end{proof}

\begin{theorem}[Vanishing of \(H^1{}'\) on affine opens]
  \label{thm:affine_open_h1_vanishing}
  \lean{AlgebraicGeometry.IsAffineOpen.subsingleton_moduleKSheaf_hModule'_one}
  \uses{def:HModule', def:moduleKSheaf, thm:cokernel_app_surjective_affine_open}
  \leanok
  \dcref{custom}
  Let \(k\) be a commutative ring, \(X\) a scheme over \(\Spec k\) and \(U\) an affine open
  of \(X\). Then
  \[
    H'^1\bigl(U, \struct{X/k}\bigr) \;=\; 0 .
  \]
  No hypothesis on \(X\) beyond the affineness of \(U\) is needed.
\end{theorem}
\begin{proof}
  \leanok
  \uses{lem:ext_one_vanishing_criterion, def:freeModuleSheafHomEquiv}
  The category of sheaves of \(k\)-modules on the small Zariski site of \(X\) is Grothendieck
  abelian, hence has enough injectives: choose a monomorphism
  \(\iota : \struct{X/k} \to G\) with \(G\) injective. By
  \ref{lem:ext_one_vanishing_criterion} (with \(L = k[U]\) the free sheaf of \(k\)-modules on
  \(U\)), it suffices to lift every morphism \(k[U] \to \coker \iota\) along
  \(G \to \coker \iota\). By \ref{def:freeModuleSheafHomEquiv}, morphisms out of \(k[U]\) are,
  compatibly with postcomposition, sections over \(U\); so the required lifting is exactly
  the surjectivity of \(G(U) \to (\coker \iota)(U)\), which is
  \ref{thm:cokernel_app_surjective_affine_open}.
\end{proof}

\begin{theorem}[Vanishing of \(H^1\) on affine schemes]
  \label{thm:affine_vanishing}
  \lean{AlgebraicGeometry.Scheme.subsingleton_moduleKSheaf_hModule_one}
  \uses{def:HModule, def:moduleKSheaf}
  \leanok
  \dcref{custom}
  Let \(k\) be a commutative ring and \(X\) an affine scheme over \(\Spec k\). Then
  \[
    H^1\bigl(X, \struct{X/k}\bigr) \;=\; 0 .
  \]
\end{theorem}
\begin{proof}
  \leanok
  \uses{lem:ext_one_vanishing_criterion, thm:cokernel_app_surjective_affine_open,
    def:constHomEquiv, lem:constHomEquiv_naturality}
  The category of sheaves of \(k\)-modules on the small Zariski site of \(X\) is Grothendieck
  abelian, hence has enough injectives: choose a monomorphism
  \(\iota : \struct{X/k} \to G\) with \(G\) injective. By
  \ref{lem:ext_one_vanishing_criterion} (with \(L = k_J\) the constant sheaf), it suffices to
  lift every morphism \(k_J \to \coker \iota\) along \(G \to \coker \iota\). By
  \ref{def:constHomEquiv} and its naturality \ref{lem:constHomEquiv_naturality}, morphisms
  out of \(k_J\) are global sections, compatibly with postcomposition; so the required
  lifting is exactly the surjectivity of \(G(X) \to (\coker \iota)(X)\), i.e.\ of
  \(G(\top) \to (\coker \iota)(\top)\), which is
  \ref{thm:cokernel_app_surjective_affine_open} applied to the affine open \(U = \top\)
  (an affine scheme is an affine open of itself).
\end{proof}

\section{Twisted affine vanishing: quasi-coherent coefficients}
\label{sec:qcoh_vanishing}

The vanishing of the previous section is stated for the structure sheaf, but Serre's argument
only ever uses two features of \(\struct{X/k}\): that ambient sections act on sections, and that
sections localise on affine opens. This section isolates those two features as a packaging on an
arbitrary sheaf of \(k\)-modules and reruns the argument, obtaining degree-one vanishing on an
affine open with those coefficients. The intended application is to line bundles trivialised on
an affine chart, for which the packaging holds chart by chart.

For a sheaf \(F\) of \(k\)-modules on the small Zariski site of \(X\) and an inclusion of opens
\(W \subseteq V\), write \(s|_W\) (\emph{section restriction}) for the image of \(s \in F(V)\)
under the \(k\)-linear restriction \(F(V) \to F(W)\); it is functorial in the inclusion and
natural in \(F\).

\begin{definition}[Quasi-coherence packaging on an open]
  \label{def:qcohOn}
  \dcref{custom}
  \lean{AlgebraicGeometry.Scheme.QcohOn}
  \leanok
  Let \(X\) be a scheme, \(U\) an open of \(X\), and \(F\) a sheaf of \(k\)-modules on the small
  Zariski site of \(X\). A \emph{quasi-coherence packaging} of \(F\) on \(U\) is the data of:
  \begin{itemize}
    \item \textbf{(P1) an action} of the fixed ring \(\Gamma(X, U)\): for every open \(W \le U\)
      a scaling operation \(r \cdot m\) of \(r \in \Gamma(X, U)\) on \(m \in F(W)\), satisfying
      the module laws over \(\Gamma(X, U)\) (unit, associativity, and both distributivities) and
      natural under section restriction, \((r \cdot m)|_{W'} = r \cdot (m|_{W'})\) for \(W' \le
      W\);
    \item \textbf{(P2) localisation on affine opens} \(V \le U\), for every \(g \in \Gamma(X,
      U)\), in the two-lemma form of localisation away from \(g\):
      \begin{itemize}
        \item \emph{denominator clearing} — every \(c \in F(V \cap D(g))\) is \((g^N \cdot
          e)|_{V \cap D(g)}\) for some \(N \ge 0\) and \(e \in F(V)\);
        \item \emph{defect annihilation} — every \(t \in F(V)\) with \(t|_{V \cap D(g)} = 0\) is
          killed by a power of \(g\): \(g^M \cdot t = 0\) for some \(M \ge 0\).
      \end{itemize}
  \end{itemize}
  The action is by the single ring \(\Gamma(X, U)\) (not a per-open module structure), with the
  inclusion \(W \le U\) an explicit argument.
\end{definition}

\begin{theorem}[The structure sheaf is quasi-coherently packaged]
  \label{thm:moduleKSheaf_qcohOn}
  \dcref{custom}
  \lean{AlgebraicGeometry.Scheme.moduleKSheaf.instQcohOn}
  \uses{def:qcohOn, def:moduleKSheaf, lem:denominator_clearing, lem:section_annihilation}
  \leanok
  For \(X\) a scheme over \(\Spec k\), the structure sheaf \(\struct{X/k}\), viewed as a sheaf of
  \(k\)-modules, carries a quasi-coherence packaging on every open \(U\): the action of \(r \in
  \Gamma(X, U)\) on \(s \in \Gamma(X, W)\) is \(r|_W \cdot s\), the ordinary
  restriction-then-multiplication.
\end{theorem}
\begin{proof}
  \leanok
  The (P1) module laws are the ring axioms of \(\Gamma(X, W)\), and naturality under restriction
  holds because restriction is a ring homomorphism. The two (P2) axioms, with the action spelled
  out, are exactly denominator clearing (\ref{lem:denominator_clearing}) and annihilation
  (\ref{lem:section_annihilation}) for the structure sheaf.
\end{proof}

\begin{theorem}[Degree-one \v{C}ech cobounding, quasi-coherent coefficients]
  \label{thm:cech_cobounding_qcoh}
  \dcref{custom}
  \lean{AlgebraicGeometry.IsAffineOpen.exists_cech_cobounding_of_qcoh}
  \uses{def:qcohOn}
  \leanok
  Let \(U\) be an affine open of a scheme \(X\), covered by basic opens \(D(f_1), \dots, D(f_n)\),
  and let \(F\) be a sheaf of \(k\)-modules with a quasi-coherence packaging on \(U\). Every
  \v{C}ech \(1\)-cocycle \(a_{ij} \in F(D(f_i) \cap D(f_j))\) (with \(a_{il} = a_{ij} + a_{jl}\)
  on triple overlaps) is a coboundary: there are \(b_i \in F(D(f_i))\) with \(a_{ij} = b_i|_{D(f_i)
  \cap D(f_j)} - b_j|_{D(f_i) \cap D(f_j)}\).
\end{theorem}
\begin{proof}
  \leanok
  Verbatim the structure-sheaf argument \ref{thm:affine_open_cech_cobounding}, reading the
  packaging axioms in place of the structure-sheaf localisation. Clear denominators with one
  uniform exponent \(N\) (denominator clearing of the packaging), producing \(e_{ij} \in
  F(D(f_i))\) with \(e_{ij}|_{D(f_i) \cap D(f_j)} = f_j^N \cdot a_{ij}\). The triple-overlap
  defects \(t_{ijl} = e_{il} - e_{jl} - f_l^N \cdot a_{ij}\) vanish on the triple intersections by
  the cocycle identity, hence are killed by a uniform power \(f_l^M\) (defect annihilation). With
  a partition of unity \(1 = \sum_l h_l f_l^{N+M}\) of the affine \(U\), set \(b_i = \sum_l h_l
  f_l^M \cdot e_{il}\); the assembly \(b_i - b_j = \sum_l h_l f_l^M(e_{il} - e_{jl}) = (\sum_l h_l
  f_l^{N+M}) \cdot a_{ij} = a_{ij}\) uses only that the action distributes over restriction and
  sums, i.e.\ the (P1) mixin.
\end{proof}

\begin{theorem}[Sections surjectivity on an affine open, quasi-coherent coefficients]
  \label{thm:cokernel_app_surjective_qcoh}
  \dcref{custom}
  \lean{AlgebraicGeometry.IsAffineOpen.cokernel_app_surjective_of_qcoh}
  \uses{def:qcohOn, thm:cech_cobounding_qcoh, lem:sections_left_exact, lem:app_injective_of_mono}
  \leanok
  Let \(U\) be an affine open of a scheme \(X\), let \(F\) be a sheaf of \(k\)-modules with a
  quasi-coherence packaging on \(U\), and let \(\iota : F \to G\) be a monomorphism of sheaves of
  \(k\)-modules. Then every section of \(\coker \iota\) over \(U\) lifts to a section of \(G\)
  over \(U\).
\end{theorem}
\begin{proof}
  \leanok
  Verbatim \ref{thm:cokernel_app_surjective_affine_open} with \(F\) in place of \(\struct{X/k}\).
  Local lifts of \(q \in (\coker \iota)(U)\) exist on a finite basic cover \(D(f_i)\) of \(U\) by
  local surjectivity of the epimorphism \(\coker\iota\)-projection and quasi-compactness. On
  overlaps the differences of the lifts are killed by the projection, so by left-exactness of
  sections (\ref{lem:sections_left_exact}) they come from \(F\)-sections \(a_{ij}\); these form a
  \v{C}ech \(1\)-cocycle (\(\iota\) is injective on sections, \ref{lem:app_injective_of_mono}),
  which cobounds by \ref{thm:cech_cobounding_qcoh}. The corrected lifts glue to a section of \(G\)
  over \(U\) mapping to \(q\).
\end{proof}

\begin{theorem}[Twisted affine vanishing of \(H^1{}'\)]
  \label{thm:hModule1_vanishing_qcoh}
  \dcref{custom}
  \lean{AlgebraicGeometry.IsAffineOpen.subsingleton_hModule'_one_of_qcoh}
  \uses{def:qcohOn, def:HModule', thm:cokernel_app_surjective_qcoh,
    lem:ext_one_vanishing_criterion, def:freeModuleSheafHomEquiv}
  \leanok
  Let \(U\) be an affine open of a scheme \(X\) and \(F\) a sheaf of \(k\)-modules with a
  quasi-coherence packaging on \(U\). Then the degree-one cohomology of the object \(U\) with
  coefficients in \(F\) vanishes:
  \[
    H'^1(U, F) \;=\; 0 .
  \]
  Taking \(F = \struct{X/k}\) recovers \ref{thm:affine_open_h1_vanishing}.
\end{theorem}
\begin{proof}
  \leanok
  As in \ref{thm:affine_open_h1_vanishing}: choose a monomorphism \(\iota : F \to G\) with \(G\)
  injective; by the \(\Ext^1\) vanishing criterion \ref{lem:ext_one_vanishing_criterion} (with
  \(L = k[U]\) the free sheaf of \(k\)-modules on \(U\)) it suffices to lift every morphism
  \(k[U] \to \coker \iota\) along \(G \to \coker \iota\), which by \ref{def:freeModuleSheafHomEquiv}
  is the surjectivity of \(G(U) \to (\coker \iota)(U)\) — Theorem
  \ref{thm:cokernel_app_surjective_qcoh}.
\end{proof}

\begin{theorem}[Twisted affine vanishing of \(H^1\)]
  \label{thm:hModule_vanishing_qcoh}
  \dcref{custom}
  \lean{AlgebraicGeometry.Scheme.subsingleton_hModule_one_of_qcoh}
  \uses{def:qcohOn, def:HModule, thm:hModule1_vanishing_qcoh, def:HModule_linearEquivHModule'}
  \leanok
  Let \(X\) be an affine scheme and \(F\) a sheaf of \(k\)-modules with a quasi-coherence
  packaging on \(\top\). Then the degree-one cohomology of the site vanishes,
  \(H^1(X, F) = 0\). Taking \(F = \struct{X/k}\) recovers \ref{thm:affine_vanishing}.
\end{theorem}
\begin{proof}
  \leanok
  Apply \ref{thm:hModule1_vanishing_qcoh} to the affine open \(U = \top\) (an affine scheme is an
  affine open of itself) and transport the object-cohomology vanishing to the site cohomology
  along the equivalence \ref{def:HModule_linearEquivHModule'}.
\end{proof}

The two (P2) axioms of the packaging are not merely \emph{sufficient} for Serre's argument:
they say precisely that restriction from an affine chart to a basic open of the chart is a
localisation of modules. The remaining nodes of this section record that bridge, which lets
any packaged sheaf on an affine chart be read through the general theory of localised
modules --- in particular, it is how the two-lattice description of the two-cover complex
consumes the packaging.

\begin{definition}[The chart-ring module structure on packaged sections]
  \label{def:qcohOn_module}
  \dcref{custom}
  \lean{AlgebraicGeometry.Scheme.QcohOn.moduleOfLE}
  \uses{def:qcohOn}
  \leanok
  Let \(F\) carry a quasi-coherence packaging on \(U\) and let \(V \le U\). The (P1) action
  makes the \(k\)-module \(F(V)\) a module over the chart ring \(\Gamma(X, U)\): the unit,
  associativity and distributivity laws are the (P1) axioms, and the remaining law
  \(0 \cdot m = 0\) follows by adding \(0 \cdot m\) to both sides of
  \((0 + 0) \cdot m = 0 \cdot m\) and cancelling.
\end{definition}

\begin{definition}[Restriction is linear over the chart ring]
  \label{def:qcohOn_secRes_linear}
  \dcref{custom}
  \lean{AlgebraicGeometry.Scheme.QcohOn.secResₗ}
  \uses{def:qcohOn_module}
  \leanok
  For opens \(V \le W \le U\), section restriction \(F(W) \to F(V)\) is
  \(\Gamma(X, U)\)-linear for the module structures of \ref{def:qcohOn_module}: it is
  additive, being a map of \(k\)-modules, and it commutes with the action by the (P1)
  restriction-naturality \((r \cdot m)|_V = r \cdot (m|_V)\).
\end{definition}

\begin{theorem}[The packaging axioms are the localisation axioms]
  \label{thm:qcohOn_isLocalizedModule}
  \dcref{custom}
  \lean{AlgebraicGeometry.Scheme.QcohOn.isLocalizedModule_secResₗ}
  \uses{def:qcohOn, def:qcohOn_module, def:qcohOn_secRes_linear}
  \leanok
  Let \(U\) be an affine open of a scheme \(X\), let \(F\) be a sheaf of \(k\)-modules with
  a quasi-coherence packaging on \(U\), and let \(g \in \Gamma(X, U)\) act invertibly on
  \(F(U \cap D(g))\) (for the module structure of \ref{def:qcohOn_module}). Then the
  restriction map \(F(U) \to F(U \cap D(g))\) exhibits \(F(U \cap D(g))\) as the
  localisation of the \(\Gamma(X, U)\)-module \(F(U)\) at the multiplicative set
  \(S = \{1, g, g^2, \dots\}\) of powers of \(g\); that is:
  \begin{enumerate}
    \item every element of \(S\) acts invertibly on \(F(U \cap D(g))\);
    \item every \(c \in F(U \cap D(g))\) is of the form \(g^{-N} \cdot (e|_{U \cap D(g)})\)
      for some \(N \ge 0\) and \(e \in F(U)\);
    \item any two sections of \(F(U)\) with equal restriction to \(U \cap D(g)\) differ by
      an element killed by a power of \(g\).
  \end{enumerate}
  These three properties characterise the localisation map \(M \to S^{-1}M\) up to unique
  isomorphism under \(F(U)\).
\end{theorem}
\begin{proof}
  \leanok
  Each axiom is the term-for-term image of one piece of the packaging data. (1) Powers of
  an invertibly-acting element act invertibly. (2) Denominator clearing on the affine open
  \(U\) itself produces, for \(c \in F(U \cap D(g))\), an exponent \(N\) and a section
  \(e \in F(U)\) with \(e|_{U \cap D(g)} = g^N \cdot c\); acting by the inverse of
  \(g^N\) gives \(c = g^{-N} \cdot (e|_{U \cap D(g)})\). (3) If
  \(x_1|_{U \cap D(g)} = x_2|_{U \cap D(g)}\) then \(x_1 - x_2\) restricts to \(0\) on
  \(U \cap D(g)\), so defect annihilation yields \(M \ge 0\) with
  \(g^M \cdot (x_1 - x_2) = 0\).
\end{proof}

\begin{theorem}[Sections over a basic open are the localised sections]
  \label{thm:qcohOn_isLocalizedModule_moduleKSheaf}
  \dcref{custom}
  \lean{AlgebraicGeometry.Scheme.QcohOn.isLocalizedModule_secResₗ_moduleKSheaf}
  \uses{thm:qcohOn_isLocalizedModule, thm:moduleKSheaf_qcohOn, def:moduleKSheaf}
  \leanok
  Let \(X\) be a scheme over \(\Spec k\), \(U\) an affine open, and \(g \in \Gamma(X, U)\).
  Then restriction exhibits \(\Gamma(X, U \cap D(g))\) as the localisation of
  \(\Gamma(X, U)\) at the powers of \(g\) --- the classical
  \(\Gamma(X, D(g)) = \Gamma(X, U)_g\), with no hypothesis on \(g\).
\end{theorem}
\begin{proof}
  \leanok
  The structure sheaf carries the packaging on \(U\) (\ref{thm:moduleKSheaf_qcohOn}), with
  the action of \(r \in \Gamma(X, U)\) given by restriction-then-multiplication. The unit
  hypothesis of \ref{thm:qcohOn_isLocalizedModule} is discharged for free: \(g\) restricts
  to a unit on \(D(g)\), hence on \(U \cap D(g)\), so multiplication by
  \(g|_{U \cap D(g)}\) --- which is the packaging action of \(g\) --- is invertible on
  \(\Gamma(X, U \cap D(g))\). Apply \ref{thm:qcohOn_isLocalizedModule}.
\end{proof}

\section{The Mayer--Vietoris \((0,1)\)-slice}
\label{sec:mv-in-text}

For a Mayer--Vietoris square \(S\) (typically: opens \(X_1 = U \cap V\), \(X_2 = U\),
\(X_3 = V\), \(X_4 = U \cup V\) of a topological space) and a sheaf of \(R\)-modules \(F\) on
the site, this section produces the exactness of
\[
  0 \to F(X_4) \to F(X_2) \times F(X_3) \to F(X_1)
    \xrightarrow{\;\delta\;} H'^1(X_4, F) \to H'^1(X_2, F) \times H'^1(X_3, F)
\]
and derives the two-cover cokernel computation: if \(H'^1(X_2, F)\) and \(H'^1(X_3, F)\)
vanish, then \(F(X_1)/\mathrm{range}(\text{difference of restrictions}) \cong H'^1(X_4, F)\),
the quotient map being the connecting homomorphism \(\delta\).

\begin{definition}[Mayer--Vietoris square]
  \label{def:MayerVietorisSquare}
  \dcref{custom}
  \lean{CategoryTheory.GrothendieckTopology.MayerVietorisSquare}
  \mathlibok
  \textit{Provided by Mathlib.}
  Let \((\mathcal C, J)\) be a small site. A Mayer--Vietoris square consists of four objects
  \(X_1, X_2, X_3, X_4\) of \(\mathcal C\) and morphisms \(f_{12} : X_1 \to X_2\),
  \(f_{13} : X_1 \to X_3\), \(f_{24} : X_2 \to X_4\), \(f_{34} : X_3 \to X_4\) such that the
  two composites \(X_1 \to X_2 \to X_4\) and \(X_1 \to X_3 \to X_4\) agree, \(f_{13}\) is a
  monomorphism, and, after applying the representable-presheaf-then-sheafification functor
  \(\mathcal C \to \mathrm{Sh}(\mathcal C, J; \mathrm{Set})\), the resulting square of
  sheaves of sets is a pushout square: the image of \(X_4\) is the pushout, over the image of
  \(X_1\), of the images of \(X_2\) and \(X_3\). The prototypical example is
  \(X_1 = U \cap V\), \(X_2 = U\), \(X_3 = V\), \(X_4 = U \cup V\) for two opens \(U, V\) of a
  topological space, with the evident inclusions.
\end{definition}

\begin{definition}[The sheaf condition of a Mayer--Vietoris square]
  \label{def:mv_sheaf_condition}
  \dcref{custom}
  \lean{CategoryTheory.GrothendieckTopology.MayerVietorisSquare.SheafCondition,
    CategoryTheory.GrothendieckTopology.MayerVietorisSquare.sheafCondition_of_sheaf}
  \uses{def:MayerVietorisSquare}
  \mathlibok
  \textit{Provided by Mathlib.}
  Let \(S = (X_1,X_2,X_3,X_4,f_{12},f_{13},f_{24},f_{34})\) be a Mayer--Vietoris square on
  \((\mathcal C, J)\) and let \(P\) be a presheaf on \(\mathcal C\). Say \(P\) satisfies the
  sheaf condition for \(S\) when the square \(P(X_4) \to P(X_2) \to P(X_1)\),
  \(P(X_4) \to P(X_3) \to P(X_1)\) (restriction along \(f_{24}, f_{12}\) and along
  \(f_{34}, f_{13}\)) is a pullback square: equivalently, restriction identifies \(P(X_4)\)
  with the pairs \((u,v) \in P(X_2) \times P(X_3)\) of equal restriction to \(X_1\). Every
  sheaf satisfies the sheaf condition for every Mayer--Vietoris square of its site.
\end{definition}

\begin{definition}[The Mayer--Vietoris short complex of free sheaves]
  \label{def:moduleShortComplex}
  \dcref{custom}
  \lean{CategoryTheory.GrothendieckTopology.MayerVietorisSquare.moduleShortComplex}
  \uses{def:MayerVietorisSquare, def:freeModuleSheaf, def:freeModuleSheafMap}
  \leanok
  For a Mayer--Vietoris square \(S\) on \((\mathcal C, J)\), the complex of sheaves of
  \(R\)-modules
  \[
    R[X_1] \xrightarrow{\ (R[f_{12}],\, -R[f_{13}])\ } R[X_2] \oplus R[X_3]
      \xrightarrow{\ R[f_{24}] + R[f_{34}]\ } R[X_4],
  \]
  whose composite is zero by the commutativity of \(S\)'s square.
\end{definition}

\begin{theorem}[The Mayer--Vietoris short complex is short exact]
  \label{thm:moduleShortComplex_shortExact}
  \dcref{custom}
  \lean{CategoryTheory.GrothendieckTopology.MayerVietorisSquare.moduleShortComplex_shortExact}
  \uses{def:moduleShortComplex}
  \leanok
  For every Mayer--Vietoris square \(S\), the complex of \ref{def:moduleShortComplex} is a
  short exact sequence of sheaves of \(R\)-modules:
  \[
    0 \to R[X_1] \to R[X_2] \oplus R[X_3] \to R[X_4] \to 0 .
  \]
\end{theorem}
\begin{proof}
  \leanok
  \emph{The square of free sheaves is a pushout.} The free-\(R\)-module-sheaf functor
  \(U \mapsto R[U]\) (\ref{def:freeModuleSheafMap} on morphisms) is a left adjoint of the
  forgetful functor to sheaves of sets (it factors the free-\(R\)-module-on-a-set functor,
  itself a left adjoint, through sheafification, also a left adjoint). Left adjoints
  preserve colimits, in particular pushout squares; applying this functor to the defining
  pushout square of \(S\) (Definition~\ref{def:MayerVietorisSquare}: the square of
  sheafified representables at \(X_1,X_2,X_3,X_4\)) shows that the square
  \(R[X_1] \to R[X_2]\), \(R[X_1] \to R[X_3]\), \(R[X_2] \to R[X_4]\),
  \(R[X_3] \to R[X_4]\) is again a pushout, now in the category of sheaves of
  \(R\)-modules.

  \emph{Pushout along a mono is short exact.} In an additive category with biproducts, a
  commuting square with legs \(p : A \to B\), \(q : A \to C\), \(r : B \to D\),
  \(s : C \to D\) is a pushout square exactly when \(D\), together with the map
  \((r,s) : B \oplus C \to D\), is a cokernel of the difference map
  \((p,-q) : A \to B \oplus C\). Applied to the pushout square above, this shows that
  \(R[X_4]\), with the sum map, is a cokernel of the difference map
  \(R[X_1] \to R[X_2] \oplus R[X_3]\); that is, the complex of \ref{def:moduleShortComplex}
  is exact at \(R[X_2] \oplus R[X_3]\) and at \(R[X_4]\), with the map to \(R[X_4]\) an
  epimorphism.

  \emph{The difference map is a monomorphism.} By definition of a Mayer--Vietoris square,
  \(f_{13} : X_1 \to X_3\) is a monomorphism of \(\mathcal C\); postcomposition with a
  monomorphism is injective on morphism sets, so the representable-presheaf map it induces
  is a monomorphism of presheaves of sets, and both the free-\(R\)-module functor (on a set,
  mapping along an injection of index sets is injective on finitely supported functions) and
  sheafification preserve monomorphisms. Hence \(R[f_{13}] : R[X_1] \to R[X_3]\) is a
  monomorphism of sheaves of \(R\)-modules, and so is its negative. The second component of
  the difference map \((R[f_{12}], -R[f_{13}]) : R[X_1] \to R[X_2] \oplus R[X_3]\) (its
  composite with the projection to \(R[X_3]\)) is exactly \(-R[f_{13}]\); since a composite
  \(g \circ h\) being a monomorphism forces \(h\) to be a monomorphism, the difference map
  itself is a monomorphism.

  Together, the three paragraphs give short exactness.
\end{proof}

\begin{definition}[The restriction-difference map]
  \label{def:moduleDiff}
  \dcref{custom}
  \lean{CategoryTheory.GrothendieckTopology.MayerVietorisSquare.moduleDiff}
  \uses{def:MayerVietorisSquare}
  \leanok
  For a Mayer--Vietoris square \(S\) and a sheaf \(F\) of \(R\)-modules, the \(R\)-linear map
  \[
    F(X_2) \times F(X_3) \longrightarrow F(X_1), \qquad
    (s_2, s_3) \longmapsto s_2|_{X_1} - s_3|_{X_1},
  \]
  restricting along \(f_{12}\) and \(f_{13}\) respectively and taking the difference.
\end{definition}

\begin{lemma}[Degree-zero exactness]
  \label{lem:mv_exact_degree0}
  \dcref{custom}
  \lean{CategoryTheory.GrothendieckTopology.MayerVietorisSquare.sections_ext,
    CategoryTheory.GrothendieckTopology.MayerVietorisSquare.exists_glue_of_moduleDiff_eq_zero}
  \uses{def:moduleDiff, def:mv_sheaf_condition}
  \leanok
  For a Mayer--Vietoris square \(S\) and a sheaf \(F\) of \(R\)-modules, the sequence
  \[
    0 \to F(X_4) \xrightarrow{\ (\mathrm{res},\,\mathrm{res})\ } F(X_2) \times F(X_3)
      \xrightarrow{\ \mathrm{moduleDiff}\ } F(X_1)
  \]
  is exact: the pair of restrictions to \(X_2, X_3\) is injective, and its image is exactly
  the kernel of \ref{def:moduleDiff}.
\end{lemma}
\begin{proof}
  \leanok
  This is the sheaf condition of \(S\) for (the sheaf of sets underlying) \(F\)
  (Definition~\ref{def:mv_sheaf_condition}): the map \(F(X_4) \to F(X_2) \times F(X_3)\) has
  image exactly the pairs with equal restriction to \(X_1\) (i.e.\
  \(\ker(\mathrm{moduleDiff})\)), and two elements of \(F(X_4)\) with the same pair of
  restrictions are equal.
\end{proof}

\begin{lemma}[Contravariant long exact sequence of \(\Ext\)]
  \label{lem:ext_contravariant_les}
  \dcref{custom}
  \lean{CategoryTheory.Abelian.Ext.contravariant_sequence_exact₁,
    CategoryTheory.Abelian.Ext.contravariant_sequence_exact₃,
    CategoryTheory.ShortComplex.ShortExact.extClass,
    CategoryTheory.Abelian.Ext.biprodAddEquiv}
  \mathlibok
  \textit{Provided by Mathlib.}
  Let \(0 \to A \xrightarrow{f} B \xrightarrow{g} C \to 0\) be a short exact sequence in an
  abelian category with \(\Ext\)-groups and let \(Y\) be any object. There is a canonical
  extension class \(\delta \in \Ext^1(C, A)\) (built from the connecting homomorphism of the
  sequence), and for every \(n \ge 0\) the sequence of abelian groups
  \[
    \Ext^n(C,Y) \xrightarrow{g^*} \Ext^n(B,Y) \xrightarrow{f^*} \Ext^n(A,Y)
      \xrightarrow{\ \delta \cdot (-)\ } \Ext^{n+1}(C,Y) \xrightarrow{g^*} \Ext^{n+1}(B,Y)
      \xrightarrow{f^*} \Ext^{n+1}(A,Y)
  \]
  is exact at each of the four inner terms, where \(f^*, g^*\) are precomposition
  (contravariant functoriality of \(\Ext(-,Y)\)) and \(\delta \cdot (-)\) is composition
  (Yoneda product) with \(\delta\). (Convention: \(\Ext^0 = \Hom\).) If \(B = X_2 \oplus X_3\)
  is a direct sum, \(\Ext^n(B,Y) \cong \Ext^n(X_2,Y) \times \Ext^n(X_3,Y)\) naturally,
  compatibly with \(f^*, g^*\).
\end{lemma}

\begin{definition}[The connecting map]
  \label{def:moduleDelta}
  \dcref{custom}
  \lean{CategoryTheory.GrothendieckTopology.MayerVietorisSquare.moduleDelta}
  \uses{def:HModule'_linearEquiv0, thm:moduleShortComplex_shortExact,
    lem:ext_contravariant_les}
  \leanok
  For a Mayer--Vietoris square \(S\) and a sheaf \(F\) of \(R\)-modules, the connecting map
  \(F(X_1) \to H'^1(X_4, F)\): the identification
  \(F(X_1) \cong H'^0(X_1,F) = \Ext^0(R[X_1], F)\) of \ref{def:HModule'_linearEquiv0},
  followed by the map \(\delta \cdot (-) : \Ext^0(R[X_1],F) \to \Ext^1(R[X_4],F) = H'^1(X_4,F)\)
  of \ref{lem:ext_contravariant_les} for the short exact sequence of
  \ref{thm:moduleShortComplex_shortExact} (\(A = R[X_1]\), \(B = R[X_2] \oplus R[X_3]\),
  \(C = R[X_4]\)).
\end{definition}

\begin{lemma}[Exactness at \(F(X_1)\)]
  \label{lem:mv_exact_at_X1}
  \dcref{custom}
  \lean{CategoryTheory.GrothendieckTopology.MayerVietorisSquare.moduleDelta_moduleDiff,
    CategoryTheory.GrothendieckTopology.MayerVietorisSquare.exists_of_moduleDelta_eq_zero}
  \uses{def:moduleDelta, def:moduleDiff}
  \leanok
  For a Mayer--Vietoris square \(S\) and a sheaf \(F\) of \(R\)-modules,
  \(\ker(\mathrm{moduleDelta}) = \mathrm{range}(\mathrm{moduleDiff})\) as submodules of
  \(F(X_1)\): the connecting class of a difference of restrictions vanishes, and conversely
  every section of \(F(X_1)\) with vanishing connecting class is such a difference.
\end{lemma}
\begin{proof}
  \leanok
  \uses{lem:ext_contravariant_les, thm:moduleShortComplex_shortExact,
    lem:freeModuleSheafHomEquiv_naturality}
  This is exactness of the sequence of \ref{lem:ext_contravariant_les} (for \(n=0\), the
  short exact sequence of \ref{thm:moduleShortComplex_shortExact}, \(Y=F\)) at the term
  \(\Ext^0(A,F) = F(X_1)\), transported along the identifications
  \(\Ext^0(B,F) \cong F(X_2) \times F(X_3)\) (\ref{lem:ext_contravariant_les}, direct sum) and
  \(\Ext^0(A,F) \cong F(X_1)\) (\ref{def:HModule'_linearEquiv0}): under these identifications
  \(f^* : \Ext^0(B,F) \to \Ext^0(A,F)\) becomes \(\mathrm{moduleDiff}\) and
  \(\delta \cdot (-) : \Ext^0(A,F) \to \Ext^1(C,F)\) becomes \(\mathrm{moduleDelta}\), by
  naturality of the section equivalence in the site variable
  (\ref{lem:freeModuleSheafHomEquiv_naturality}), which identifies precomposition at the
  level of morphisms out of a free sheaf with restriction at the level of sections.
\end{proof}

\begin{lemma}[Exactness at \(H'^1(X_4,F)\)]
  \label{lem:mv_exact_at_H1X4}
  \dcref{custom}
  \lean{CategoryTheory.GrothendieckTopology.MayerVietorisSquare.res_moduleDelta₂,
    CategoryTheory.GrothendieckTopology.MayerVietorisSquare.res_moduleDelta₃,
    CategoryTheory.GrothendieckTopology.MayerVietorisSquare.exists_moduleDelta_eq}
  \uses{def:moduleDelta, def:HModule'_res}
  \leanok
  For a Mayer--Vietoris square \(S\) and a sheaf \(F\) of \(R\)-modules,
  \(\ker(\mathrm{res}_{f_{24}}) \cap \ker(\mathrm{res}_{f_{34}}) = \mathrm{range}(\mathrm{moduleDelta})\)
  as submodules of \(H'^1(X_4,F)\), where
  \(\mathrm{res}_{f_{24}}, \mathrm{res}_{f_{34}} : H'^1(X_4,F) \to H'^1(X_2,F), H'^1(X_3,F)\)
  are the restrictions of \ref{def:HModule'_res}: the connecting class of any section of
  \(F(X_1)\) restricts to zero on both \(X_2\) and \(X_3\), and conversely every class
  restricting to zero on both is a connecting class.
\end{lemma}
\begin{proof}
  \leanok
  \uses{lem:ext_contravariant_les, thm:moduleShortComplex_shortExact}
  As in \ref{lem:mv_exact_at_X1}, this is exactness of the sequence of
  \ref{lem:ext_contravariant_les} (for \(n=0\), i.e.\ at the term \(\Ext^1(C,F)\) between
  \(\delta \cdot (-)\) and \(g^*\)) transported along the same identifications, together
  with the splitting \(\Ext^1(B,F) \cong H'^1(X_2,F) \times H'^1(X_3,F)\) under which
  \(g^*\) becomes the pair \((\mathrm{res}_{f_{24}}, \mathrm{res}_{f_{34}})\).
\end{proof}

\begin{theorem}[Surjectivity of the connecting map]
  \label{thm:moduleDelta_surjective}
  \dcref{custom}
  \lean{CategoryTheory.GrothendieckTopology.MayerVietorisSquare.moduleDelta_surjective}
  \uses{lem:mv_exact_at_H1X4}
  \leanok
  If \(H'^1(X_2,F)\) and \(H'^1(X_3,F)\) are both zero, then
  \(\mathrm{moduleDelta} : F(X_1) \to H'^1(X_4,F)\) is surjective.
\end{theorem}
\begin{proof}
  \leanok
  Immediate from \ref{lem:mv_exact_at_H1X4}:
  \(\ker(\mathrm{res}_{f_{24}}) \cap \ker(\mathrm{res}_{f_{34}}) = H'^1(X_4,F)\) once both
  targets vanish, so \(\mathrm{range}(\mathrm{moduleDelta}) = H'^1(X_4,F)\).
\end{proof}

\begin{definition}[The connecting map on the cokernel]
  \label{def:moduleDeltaQuotient}
  \dcref{custom}
  \lean{CategoryTheory.GrothendieckTopology.MayerVietorisSquare.moduleDeltaQuotient}
  \uses{def:moduleDelta, lem:mv_exact_at_X1}
  \leanok
  Since \(\mathrm{range}(\mathrm{moduleDiff}) \subseteq \ker(\mathrm{moduleDelta})\)
  (\ref{lem:mv_exact_at_X1}), \(\mathrm{moduleDelta}\) descends to an \(R\)-linear map
  \[
    F(X_1)/\mathrm{range}(\mathrm{moduleDiff}) \longrightarrow H'^1(X_4,F).
  \]
\end{definition}

\begin{theorem}[The Mayer--Vietoris two-cover cokernel computation]
  \label{thm:h1LinearEquiv}
  \lean{CategoryTheory.GrothendieckTopology.MayerVietorisSquare.h1LinearEquiv}
  \uses{def:moduleDeltaQuotient, lem:mv_exact_at_X1, thm:moduleDelta_surjective}
  \leanok
  \dcref{custom}
  If \(H'^1(X_2,F)\) and \(H'^1(X_3,F)\) are both zero, the map of
  \ref{def:moduleDeltaQuotient} is an \(R\)-linear isomorphism
  \[
    F(X_1)/\mathrm{range}(\mathrm{moduleDiff}) \;\cong\; H'^1(X_4,F).
  \]
\end{theorem}
\begin{proof}
  \leanok
  Injectivity: if the class of \(s \in F(X_1)\) maps to \(0\), then
  \(\mathrm{moduleDelta}(s) = 0\), so by \ref{lem:mv_exact_at_X1}
  \(s \in \mathrm{range}(\mathrm{moduleDiff})\), i.e.\ the class of \(s\) is already \(0\).
  Surjectivity is \ref{thm:moduleDelta_surjective}.
\end{proof}

\begin{lemma}[Compatibility with the connecting map]
  \label{lem:h1LinearEquiv_mk}
  \dcref{custom}
  \lean{CategoryTheory.GrothendieckTopology.MayerVietorisSquare.h1LinearEquiv_mk}
  \uses{thm:h1LinearEquiv}
  \leanok
  Under the isomorphism of \ref{thm:h1LinearEquiv}, the class of \(s \in F(X_1)\)
  corresponds to \(\mathrm{moduleDelta}(s)\).
\end{lemma}
\begin{proof}
  \leanok
  Immediate from the construction of \ref{def:moduleDeltaQuotient} and
  \ref{thm:h1LinearEquiv}.
\end{proof}

\section{The two-affine-cover cokernel bridge}

For a scheme \(X\) over \(\Spec k\) covered by two affine opens \(U_0 \cup U_1 = X\), this
section computes the degree-one cohomology of the structure sheaf as
\[
  H^1\bigl(X, \struct{X/k}\bigr) \;\cong\;
    \Gamma(X, U_0 \cap U_1) \big/
      \mathrm{range}\bigl((s_0,s_1) \mapsto s_0|_{U_0 \cap U_1} - s_1|_{U_0 \cap U_1}\bigr),
\]
the computational access point for the genus of a curve. No hypothesis on the intersection
\(U_0 \cap U_1\) is needed: the derived Mayer--Vietoris sequence of Section~\ref{sec:mv-in-text}
is exact regardless, and affineness of the two pieces alone kills their degree-one
cohomology (\ref{thm:affine_open_h1_vanishing}) --- the classical affine-intersection
hypothesis is needed only for the \v{C}ech-complex identification of \(H^1\), which this
route bypasses.

\begin{definition}[Mayer--Vietoris square of two opens]
  \label{def:opens_mayerVietorisSquare'}
  \dcref{custom}
  \lean{Opens.mayerVietorisSquare'}
  \uses{def:MayerVietorisSquare}
  \mathlibok
  \textit{Provided by Mathlib.}
  For a topological space \(T\) and opens \(X_1,X_2,X_3,X_4 \subseteq T\) with
  \(X_1 = X_2 \cap X_3\) and \(X_4 = X_2 \cup X_3\), the square \((X_1,X_2,X_3,X_4)\) with the
  inclusion morphisms is a Mayer--Vietoris square (Definition~\ref{def:MayerVietorisSquare})
  on the small Zariski site of \(T\).
\end{definition}

\begin{definition}[The two-cover square]
  \label{def:twoCoverSquare}
  \dcref{custom}
  \lean{AlgebraicGeometry.Scheme.twoCoverSquare}
  \uses{def:opens_mayerVietorisSquare'}
  \leanok
  For a scheme \(X\) and opens \(U_0, U_1 \subseteq X\) with \(U_0 \cup U_1 = X\), the
  Mayer--Vietoris square \((U_0 \cap U_1,\, U_0,\, U_1,\, X)\) on the small Zariski site of
  \(X\) (Definition~\ref{def:opens_mayerVietorisSquare'}, specialized to \(X_4 = X\)), with
  the inclusion morphisms.
\end{definition}

\begin{theorem}[The two-cover \(H^1\) computation, general coefficients]
  \label{thm:twoCoverH1LinearEquiv}
  \lean{AlgebraicGeometry.Scheme.twoCoverH1LinearEquiv}
  \uses{def:twoCoverSquare, def:HModule_linearEquivHModule', thm:h1LinearEquiv}
  \leanok
  \dcref{custom}
  Let \(k\) be a commutative ring, \(X\) a scheme over \(\Spec k\), \(U_0, U_1 \subseteq X\)
  opens with \(U_0 \cup U_1 = X\), and \(F\) a sheaf of \(k\)-modules on the small Zariski
  site of \(X\) with \(H'^1(U_0,F) = H'^1(U_1,F) = 0\). Then
  \[
    H^1(X,F) \;\cong\;
      F(U_0 \cap U_1) \big/ \mathrm{range}\bigl((s_0,s_1) \mapsto
        s_0|_{U_0 \cap U_1} - s_1|_{U_0 \cap U_1}\bigr) .
  \]
\end{theorem}
\begin{proof}
  \leanok
  Compose the isomorphism \(H^1(X,F) \cong H'^1(X,F)\) of
  \ref{def:HModule_linearEquivHModule'} (\(X\) is a terminal object of its own small Zariski
  site) with the inverse of the isomorphism of \ref{thm:h1LinearEquiv} for the two-cover
  square of \ref{def:twoCoverSquare}, whose \(\mathrm{moduleDiff}\) is exactly the displayed
  restriction-difference map.
\end{proof}

\begin{definition}[The two-cover restriction-difference map]
  \label{def:TwoCover_diff}
  \dcref{custom}
  \lean{AlgebraicGeometry.TwoCover.diff}
  \uses{def:moduleKSheaf}
  \leanok
  Let \(k\) be a commutative ring, \(X\) a scheme over \(\Spec k\), and \(U_0, U_1 \subseteq X\)
  opens. The \(k\)-linear map
  \[
    \Gamma(X,U_0) \times \Gamma(X,U_1) \longrightarrow \Gamma(X, U_0 \cap U_1), \qquad
    (s_0,s_1) \longmapsto s_0|_{U_0 \cap U_1} - s_1|_{U_0 \cap U_1} .
  \]
  When \(U_0 \cup U_1 = X\), this is the map \ref{def:moduleDiff} of the two-cover square
  (Definition~\ref{def:twoCoverSquare}) for the structure sheaf \(\struct{X/k}\).
\end{definition}

\begin{definition}[The two-cover cokernel carrier]
  \label{def:TwoCover_H1Cok}
  \dcref{custom}
  \lean{AlgebraicGeometry.TwoCover.H1Cok}
  \uses{def:TwoCover_diff}
  \leanok
  The \(k\)-module \(\Gamma(X, U_0 \cap U_1)/\mathrm{range}(\mathrm{diff})\), the cokernel of
  the map of \ref{def:TwoCover_diff}.
\end{definition}

\begin{lemma}[Sections extending to one piece vanish in the cokernel]
  \label{lem:TwoCover_h1Cok_mk_res}
  \dcref{custom}
  \lean{AlgebraicGeometry.TwoCover.h1Cok_mk_resHom_left,
    AlgebraicGeometry.TwoCover.h1Cok_mk_resHom_right}
  \uses{def:TwoCover_H1Cok}
  \leanok
  For \(t \in \Gamma(X,U_0)\), the class of \(t|_{U_0 \cap U_1}\) in
  \ref{def:TwoCover_H1Cok} is zero; symmetrically, for \(t \in \Gamma(X,U_1)\) the class of
  \(t|_{U_0 \cap U_1}\) is zero.
\end{lemma}
\begin{proof}
  \leanok
  The first is the class of \(\mathrm{diff}(t,0)\), the second of \(\mathrm{diff}(0,-t)\)
  (Definition~\ref{def:TwoCover_diff}), both in the range of \(\mathrm{diff}\).
\end{proof}

\begin{definition}[The two-cover connecting map]
  \label{def:TwoCover_delta}
  \dcref{custom}
  \lean{AlgebraicGeometry.TwoCover.delta}
  \uses{def:twoCoverSquare, def:HModule_linearEquivHModule', def:moduleDelta}
  \leanok
  Assume \(U_0 \cup U_1 = X\). The \(k\)-linear map
  \(\Gamma(X, U_0 \cap U_1) \to H^1(X, \struct{X/k})\): the connecting map
  \ref{def:moduleDelta} of the two-cover square (Definition~\ref{def:twoCoverSquare}) for
  \(F = \struct{X/k}\), followed by the inverse of the isomorphism
  \(H^1(X,\struct{X/k}) \cong H'^1(X,\struct{X/k})\) of
  \ref{def:HModule_linearEquivHModule'}.
\end{definition}

\begin{lemma}[The connecting class of a difference vanishes]
  \label{lem:TwoCover_delta_diff}
  \dcref{custom}
  \lean{AlgebraicGeometry.TwoCover.delta_diff}
  \uses{def:TwoCover_delta, def:TwoCover_diff, lem:mv_exact_at_X1}
  \leanok
  For \(t \in \Gamma(X,U_0) \times \Gamma(X,U_1)\), \(\mathrm{delta}(\mathrm{diff}(t)) = 0\).
\end{lemma}
\begin{proof}
  \leanok
  \(\mathrm{diff}(t)\) equals the two-cover square's \(\mathrm{moduleDiff}(t)\)
  (Definition~\ref{def:TwoCover_diff}), whose connecting class vanishes by
  \ref{lem:mv_exact_at_X1}; transport along the (linear, hence class-zero-preserving)
  isomorphism of \ref{def:HModule_linearEquivHModule'} used in the definition of
  \(\mathrm{delta}\).
\end{proof}

\begin{lemma}[The connecting class of an extending section vanishes]
  \label{lem:TwoCover_delta_resHom}
  \dcref{custom}
  \lean{AlgebraicGeometry.TwoCover.delta_resHom_left,
    AlgebraicGeometry.TwoCover.delta_resHom_right}
  \uses{lem:TwoCover_delta_diff}
  \leanok
  For \(t \in \Gamma(X,U_0)\), \(\mathrm{delta}(t|_{U_0 \cap U_1}) = 0\); symmetrically, for
  \(t \in \Gamma(X,U_1)\), \(\mathrm{delta}(t|_{U_0 \cap U_1}) = 0\).
\end{lemma}
\begin{proof}
  \leanok
  \(t|_{U_0 \cap U_1} = \mathrm{diff}(t,0)\), resp.\ \(-\mathrm{diff}(0,t)\)
  (Definition~\ref{def:TwoCover_diff}), so this is \ref{lem:TwoCover_delta_diff} (using
  \(\mathrm{delta}\)'s linearity in the second case).
\end{proof}

\begin{theorem}[Surjectivity of the two-cover connecting map]
  \label{thm:TwoCover_delta_surjective}
  \dcref{custom}
  \lean{AlgebraicGeometry.TwoCover.delta_surjective}
  \uses{def:TwoCover_delta, thm:affine_open_h1_vanishing, thm:moduleDelta_surjective}
  \leanok
  If \(U_0\) and \(U_1\) are affine opens of \(X\) (still with \(U_0 \cup U_1 = X\)), then
  \(\mathrm{delta} : \Gamma(X, U_0 \cap U_1) \to H^1(X,\struct{X/k})\) is surjective.
\end{theorem}
\begin{proof}
  \leanok
  By \ref{thm:affine_open_h1_vanishing}, \(H'^1(U_0,\struct{X/k}) = H'^1(U_1,\struct{X/k}) = 0\);
  by \ref{thm:moduleDelta_surjective} applied to the two-cover square, its connecting map
  \(\mathrm{moduleDelta}\) is surjective onto \(H'^1(X,\struct{X/k})\), and transporting
  along the isomorphism of \ref{def:HModule_linearEquivHModule'} used to define
  \(\mathrm{delta}\) preserves surjectivity.
\end{proof}

\begin{theorem}[The two-affine-cover \(H^1\) cokernel bridge]
  \label{thm:TwoCover_h1CokEquiv}
  \lean{AlgebraicGeometry.TwoCover.h1CokEquiv}
  \uses{thm:affine_open_h1_vanishing, thm:h1LinearEquiv, def:HModule_linearEquivHModule',
    def:TwoCover_H1Cok}
  \leanok
  \dcref{custom}
  Let \(k\) be a commutative ring, \(X\) a scheme over \(\Spec k\), and \(U_0, U_1\) affine
  opens of \(X\) with \(U_0 \cup U_1 = X\). There is a \(k\)-linear isomorphism
  \[
    H^1\bigl(X, \struct{X/k}\bigr) \;\cong\; \Gamma(X, U_0 \cap U_1)\big/\mathrm{range}(\mathrm{diff}),
  \]
  where \(\mathrm{diff}(s_0,s_1) = s_0|_{U_0 \cap U_1} - s_1|_{U_0 \cap U_1}\)
  (\ref{def:TwoCover_diff}). No hypothesis on \(U_0 \cap U_1\) (affineness, separatedness, or
  even quasi-compactness) is needed: the derived Mayer--Vietoris sequence of
  \ref{thm:moduleShortComplex_shortExact} is exact regardless, and affineness of \(U_0, U_1\)
  alone kills their degree-one cohomology (\ref{thm:affine_open_h1_vanishing}); the classical
  hypothesis on the intersection is needed only for the \v{C}ech-complex identification of
  \(H^1\), which this route bypasses.
\end{theorem}
\begin{proof}
  \leanok
  By \ref{thm:affine_open_h1_vanishing}, \(H'^1(U_0,\struct{X/k}) = H'^1(U_1,\struct{X/k}) = 0\).
  Apply \ref{thm:h1LinearEquiv} to the two-cover square of \ref{def:twoCoverSquare} with
  \(F = \struct{X/k}\), giving
  \(\Gamma(X, U_0 \cap U_1)/\mathrm{range}(\mathrm{diff}) \cong H'^1(X,\struct{X/k})\), and
  compose with the inverse of the isomorphism
  \(H^1(X,\struct{X/k}) \cong H'^1(X,\struct{X/k})\) of
  \ref{def:HModule_linearEquivHModule'}.
\end{proof}

\begin{lemma}[Compatibility with the connecting map]
  \label{lem:TwoCover_h1CokEquiv_delta}
  \dcref{custom}
  \lean{AlgebraicGeometry.TwoCover.h1CokEquiv_delta}
  \uses{thm:TwoCover_h1CokEquiv, def:TwoCover_delta, lem:h1LinearEquiv_mk}
  \leanok
  Under the isomorphism of \ref{thm:TwoCover_h1CokEquiv}, the class
  \(\mathrm{delta}(s)\) (Definition~\ref{def:TwoCover_delta}) corresponds, for every
  \(s \in \Gamma(X, U_0 \cap U_1)\), to the class of \(s\) in
  \(\Gamma(X,U_0 \cap U_1)/\mathrm{range}(\mathrm{diff})\).
\end{lemma}
\begin{proof}
  \leanok
  Unwind the definitions of \ref{thm:TwoCover_h1CokEquiv} and \ref{def:TwoCover_delta} (both
  built from the same two ingredients, \ref{thm:h1LinearEquiv} and
  \ref{def:HModule_linearEquivHModule'}) and apply the compatibility \ref{lem:h1LinearEquiv_mk}
  of the Mayer--Vietoris cokernel isomorphism with its connecting map.
\end{proof}

\begin{lemma}[Compatibility with the connecting map, inverse form]
  \label{lem:TwoCover_h1CokEquiv_symm_mk}
  \dcref{custom}
  \lean{AlgebraicGeometry.TwoCover.h1CokEquiv_symm_mk}
  \uses{lem:TwoCover_h1CokEquiv_delta}
  \leanok
  The inverse of the isomorphism of \ref{thm:TwoCover_h1CokEquiv} sends the class of
  \(s \in \Gamma(X, U_0 \cap U_1)\) to \(\mathrm{delta}(s)\).
\end{lemma}
\begin{proof}
  \leanok
  Apply the isomorphism of \ref{thm:TwoCover_h1CokEquiv} to both sides of
  \ref{lem:TwoCover_h1CokEquiv_delta} and use injectivity.
\end{proof}

\section{Base change of the two-cover cohomology over an affine test ring}

Fix a field \(k\) and an object \(C\) of \(\Over(\Spec k)\). For a commutative \(k\)-algebra \(R\)
--- the \emph{affine test ring} --- base change produces the relative curve
\(C_R = C \otimes \Spec R\) over \(\Spec R\); an affine two-chart cover of \(C\)
(Section~\ref{sec:mv-in-text} and Chapter~\ref{chap:Curves}) pulls back to one of \(C_R\), and the
two-affine-cover bridge of the previous section then computes \(H^1(C_R, \struct{C_R/R})\) as an
explicit two-cover cokernel of \(R\)-modules. When \(C\) is a smooth curve, the degree-zero
cohomology of the relative curve is the test ring itself, \(H^0(C_R, \struct{C_R/R}) = R\).

\begin{definition}[The relative curve over the test ring]
  \label{def:relCurve}
  \dcref{custom}
  \lean{AlgebraicGeometry.relCurve, AlgebraicGeometry.relCurve.instOver}
  \leanok
  For \(C \in \Over(\Spec k)\) and a commutative \(k\)-algebra \(R\), the \emph{relative curve}
  \(C_R\) is the underlying scheme of the fibre product \(C \otimes \Spec R\) of \(C\) with the
  object \(\Spec R\) of \(\Over(\Spec k)\), regarded as a scheme over \(\Spec R\) via the second
  projection \(\mathrm{pr}_2 : C \otimes \Spec R \to \Spec R\).
\end{definition}

\begin{definition}[The base-changed affine two-cover]
  \label{def:relCover}
  \dcref{custom}
  \lean{AlgebraicGeometry.relCover}
  \uses{def:relCurve, def:AffineTwoCover, thm:AffineTwoCover_pullbackProd}
  \leanok
  Let \(D = (V_0, V_1)\) be an affine two-chart cover (\ref{def:AffineTwoCover}) of the underlying
  scheme of \(C\). The \emph{base-changed cover} of the relative curve \(C_R\) is the pair of
  preimages \(\bigl(\mathrm{pr}_1^{-1} V_0,\ \mathrm{pr}_1^{-1} V_1\bigr)\) of the charts of \(D\)
  under the first projection \(\mathrm{pr}_1 : C_R \to C\); by
  \ref{thm:AffineTwoCover_pullbackProd} it is again an affine two-chart cover, now of \(C_R\).
\end{definition}

\begin{lemma}[The base-changed charts are affine and cover the relative curve]
  \label{lem:relCover_charts}
  \dcref{custom}
  \lean{AlgebraicGeometry.relCover_isAffineOpen₀, AlgebraicGeometry.relCover_isAffineOpen₁,
    AlgebraicGeometry.relCover_sup}
  \uses{def:relCover}
  \leanok
  The two charts \(\mathrm{pr}_1^{-1} V_0\) and \(\mathrm{pr}_1^{-1} V_1\) of the base-changed cover
  are affine open, and their union is all of \(C_R\).
\end{lemma}
\begin{proof}
  \leanok
  \uses{thm:AffineTwoCover_pullbackProd}
  These are the affineness of the two charts and the covering property recorded in the
  affine-two-chart-cover structure \(\mathrm{pr}_1^{-1} D\) of \ref{thm:AffineTwoCover_pullbackProd}.
\end{proof}

\begin{lemma}[The overlap of the base-changed charts]
  \label{lem:relCover_inf}
  \dcref{custom}
  \lean{AlgebraicGeometry.relCover_inf}
  \uses{def:relCover}
  \leanok
  The overlap of the two base-changed charts is the preimage of the overlap of \(D\):
  \[
    \mathrm{pr}_1^{-1} V_0 \cap \mathrm{pr}_1^{-1} V_1 = \mathrm{pr}_1^{-1}(V_0 \cap V_1).
  \]
\end{lemma}
\begin{proof}
  \leanok
  Taking preimages under \(\mathrm{pr}_1\) commutes with intersection of opens.
\end{proof}

\begin{definition}[The relative two-cover cohomology carrier]
  \label{def:relTwoCoverH1}
  \dcref{custom}
  \lean{AlgebraicGeometry.relTwoCoverH1}
  \uses{def:relCurve, def:relCover, def:moduleKSheaf, def:TwoCover_diff, def:TwoCover_H1Cok}
  \leanok
  For the structure sheaf of the relative curve \(C_R\) viewed as a sheaf of \(R\)-modules, its
  degree-one cohomology over \(R\) is the two-cover cokernel of the base-changed cover: an
  \(R\)-linear isomorphism
  \[
    H^1\bigl(C_R, \struct{C_R/R}\bigr) \;\cong\;
      \Gamma\bigl(C_R,\ \mathrm{pr}_1^{-1}V_0 \cap \mathrm{pr}_1^{-1}V_1\bigr)\big/
        \mathrm{range}(\mathrm{diff}),
  \]
  where \(\mathrm{diff}(s_0,s_1) = s_0|_\cap - s_1|_\cap\) (\ref{def:TwoCover_diff}).
\end{definition}
\begin{proof}
  \leanok
  \uses{thm:TwoCover_h1CokEquiv, lem:relCover_charts}
  The base-changed charts are affine and cover \(C_R\) (\ref{lem:relCover_charts}), so this is the
  two-affine-cover \(H^1\) cokernel bridge \ref{thm:TwoCover_h1CokEquiv} applied over the base ring
  \(R\) to the scheme \(C_R\) with its base-changed cover; the affine-vanishing inputs required by
  that bridge are exactly the affineness of the two charts.
\end{proof}

\begin{theorem}[Degree-zero cohomology of the relative curve is the test ring]
  \label{thm:relStructureSectionsTop}
  \dcref{custom}
  \lean{AlgebraicGeometry.relStructureSectionsTop}
  \uses{def:relCurve, thm:isIso_appLE_snd}
  \leanok
  Let \(C\) be a smooth curve over \(k\) (proper, smooth of relative dimension one, geometrically
  irreducible), so that \(C\) is in particular geometrically reduced. For every commutative
  \(k\)-algebra \(R\), the degree-zero cohomology of the structure sheaf of the relative curve is
  the test ring itself:
  \[
    H^0\bigl(C_R, \struct{C_R/R}\bigr) \;=\; \Gamma(C_R, \top) \;\cong\; R,
  \]
  the isomorphism being pullback of global sections along the second projection
  \(\mathrm{pr}_2 : C_R \to \Spec R\). This is the degree-zero case of ``cohomology commutes with
  base change'', universal in the test ring \(R\).
\end{theorem}
\begin{proof}
  \leanok
  \uses{def:moduleKSheafHZero}
  Degree-zero cohomology is global sections, \(H^0(C_R, \struct{C_R/R}) \cong \Gamma(C_R, \top)\)
  (\ref{def:moduleKSheafHZero}, over the base ring \(R\)). Specialize the base-change isomorphism
  \ref{thm:isIso_appLE_snd} to the test object \(T = \Spec R\) and its affine open \(V = \top\):
  pullback along the second projection is an isomorphism
  \(\Gamma(\Spec R, \top) \xrightarrow{\sim} \Gamma(C_R, \mathrm{pr}_2^{-1}\top) = \Gamma(C_R, \top)\)
  (the hypotheses of \ref{thm:isIso_appLE_snd} --- \(C\) proper, geometrically irreducible and
  geometrically reduced --- hold for a smooth curve). Composing with the identification
  \(\Gamma(\Spec R, \top) \cong R\) gives \(\Gamma(C_R, \top) \cong R\).
\end{proof}

\section{The projective line: inputs for the finiteness argument}

This section collects the facts about \(\PP^1\) and about morphisms over \(\Spec k\) that
the finiteness argument of the next section consumes. It builds on the description of
\(\PP^1\), its standard charts \(U_0, U_1\), the chart coordinates
\(t_{ij} = X_j/X_i \in A_{(X_i)}\), and the section-ring identifications
\(\gamma_0, \gamma_1, \gamma_{01}\) developed in Chapter~\ref{chap:Curves}
(Definitions~\ref{def:P1_chartOpen}, \ref{def:P1_chartCoord},
\ref{def:P1_chartSectionsEquiv0}--\ref{def:P1_overlapSectionsEquiv}); the statements here
are recorded in this chapter because they are private inputs of the finiteness glue rather
than general theory of the projective line. Throughout, for a homogeneous \(f\) of positive
degree we write \(A_{(f)} \to \Gamma(\PP^1, D_+(f))\) for the canonical map from the
degree-zero localization to the sections over the basic open (an isomorphism).

\begin{lemma}[Morphisms over the base intertwine the structure algebra maps]
  \label{lem:appTop_map_appLE}
  \dcref{custom}
  \lean{AlgebraicGeometry.Scheme.Hom.appTop_map_appLE}
  \leanok
  Let \(k\) be a commutative ring, let \(X, Y\) be schemes over \(\Spec k\) and let
  \(\pi : X \to Y\) be a morphism compatible with the structure morphisms. For opens
  \(U \subseteq Y\) and \(V \subseteq X\) with \(V \subseteq \pi^{-1}(U)\), the two
  composite ring homomorphisms
  \[
    \Gamma(\Spec k, \top) \to \Gamma(Y, \top) \to \Gamma(Y, U)
      \xrightarrow{\ \pi^\sharp\ } \Gamma(X, V)
    \qquad\text{and}\qquad
    \Gamma(\Spec k, \top) \to \Gamma(X, \top) \to \Gamma(X, V)
  \]
  agree, where the unlabeled maps are induced by the structure morphisms and by
  restriction.
\end{lemma}
\begin{proof}
  \leanok
  Pulling back sections along \(\pi\) commutes with restriction, so the first composite is
  the map on sections induced by \(\pi\) composed with the structure morphism of \(Y\),
  restricted to \(V\); by compatibility this composite of morphisms is the structure
  morphism of \(X\).
\end{proof}

\begin{lemma}[Pulling back the structure algebra map]
  \label{lem:appLE_overAlgebraMap}
  \dcref{custom}
  \lean{AlgebraicGeometry.Scheme.Hom.appLE_overAlgebraMap}
  \uses{def:overAlgebraMap, lem:appTop_map_appLE}
  \leanok
  In the situation of \ref{lem:appTop_map_appLE}, for every \(r \in k\) the map
  \(\pi^\sharp : \Gamma(Y, U) \to \Gamma(X, V)\) sends the image of \(r\) under the
  structure algebra map of \(Y\) (Definition~\ref{def:overAlgebraMap}) to the image of
  \(r\) under the structure algebra map of \(X\).
\end{lemma}
\begin{proof}
  \leanok
  Evaluate the equality of ring homomorphisms of \ref{lem:appTop_map_appLE} at the section
  corresponding to \(r\).
\end{proof}

\begin{lemma}[The basic-open models lie over the base]
  \label{lem:P1_awayi_structureMap}
  \dcref{custom}
  \lean{AlgebraicGeometry.P1.awayι_structureMap}
  \uses{def:P1_structureMap}
  \leanok
  Let \(f \in k[X_0, X_1]\) be homogeneous of positive degree. The composite of the affine
  model \(\Spec A_{(f)} \to \PP^1\) of the basic open \(D_+(f)\) with the structure
  morphism \(\PP^1 \to \Spec k\) is the morphism induced by the algebra map
  \(k \to A_{(f)}\). (This generalizes \ref{lem:P1_chartIota_structureMap} from the two
  standard charts to all standard basic opens.)
\end{lemma}
\begin{proof}
  \leanok
  \uses{def:P1_gradeZero}
  As for \ref{lem:P1_chartIota_structureMap}: the canonical morphism
  \(\Proj A \to \Spec A_0\), restricted to the affine model of a basic open, is induced by
  the ring homomorphism \(A_0 \to A_{(f)}\), and composing with the isomorphism
  \(k \cong A_0\) of \ref{def:P1_gradeZero} gives the claim.
\end{proof}

\begin{lemma}[The basic-open inclusion factors through its affine model]
  \label{lem:P1_basicOpen_iota_eq}
  \dcref{custom}
  \lean{AlgebraicGeometry.P1.basicOpen_ι_eq}
  \uses{def:P1}
  \leanok
  Let \(f \in k[X_0, X_1]\) be homogeneous of positive degree. The inclusion
  \(D_+(f) \to \PP^1\) is the composite of the canonical identification
  \(D_+(f) \cong \Spec A_{(f)}\) with the affine model \(\Spec A_{(f)} \to \PP^1\).
\end{lemma}
\begin{proof}
  \leanok
  The affine model is by construction an open immersion with image \(D_+(f)\), and the
  canonical identification is the induced isomorphism onto the image; composing the two
  recovers the inclusion.
\end{proof}

\begin{lemma}[The structure algebra map on a basic open]
  \label{lem:P1_structureMap_appTop_awayToSection}
  \dcref{custom}
  \lean{AlgebraicGeometry.P1.structureMap_appTop_awayToSection}
  \uses{def:P1_structureMap, lem:P1_awayi_structureMap, lem:P1_basicOpen_iota_eq}
  \leanok
  Let \(f \in k[X_0, X_1]\) be homogeneous of positive degree. The composite
  \[
    \Gamma(\Spec k, \top) \to \Gamma(\PP^1, \top) \to \Gamma\bigl(\PP^1, D_+(f)\bigr)
  \]
  (structure morphism on global sections, then restriction) is the composite of the
  identification \(\Gamma(\Spec k, \top) \cong k\), the algebra map \(k \to A_{(f)}\), and
  the canonical map \(A_{(f)} \to \Gamma(\PP^1, D_+(f))\).
\end{lemma}
\begin{proof}
  \leanok
  By \ref{lem:P1_basicOpen_iota_eq} and \ref{lem:P1_awayi_structureMap}, the composite of
  the inclusion \(D_+(f) \to \PP^1\) with the structure morphism is the identification
  \(D_+(f) \cong \Spec A_{(f)}\) followed by the morphism induced by \(k \to A_{(f)}\).
  Taking global sections of this equality of scheme morphisms, and identifying the global
  sections of the open subscheme \(D_+(f)\) with \(\Gamma(\PP^1, D_+(f))\) --- under which
  the map on sections induced by the identification with \(\Spec A_{(f)}\) becomes the
  canonical map \(A_{(f)} \to \Gamma(\PP^1, D_+(f))\) --- yields the claim.
\end{proof}

\begin{lemma}[The distinguished fraction of a chart pair is the chart coordinate]
  \label{lem:P1_isLocalizationElem_X_eq}
  \dcref{custom}
  \lean{AlgebraicGeometry.P1.isLocalizationElem_X_eq}
  \uses{def:P1_chartCoord}
  \leanok
  For \(i, j \in \{0,1\}\), the canonical degree-zero fraction
  \(X_j^{\deg X_i} / X_i^{\deg X_j}\) attached to the pair \((X_i, X_j)\) of homogeneous
  elements equals the chart coordinate \(t_{ij} = X_j / X_i \in A_{(X_i)}\).
\end{lemma}
\begin{proof}
  \leanok
  Both degrees are one, and a first power is the element itself.
\end{proof}

\begin{theorem}[The chart coordinate cuts out the overlap]
  \label{thm:P1_basicOpen_chartCoord}
  \dcref{custom}
  \lean{AlgebraicGeometry.P1.basicOpen_awayToSection_chartCoord}
  \uses{def:P1_chartOpen, def:P1_chartCoord, lem:P1_basicOpen_iota_eq,
    lem:P1_isLocalizationElem_X_eq}
  \leanok
  For \(i, j \in \{0,1\}\), regard the chart coordinate \(t_{ij} \in A_{(X_i)}\) as a
  section of the structure sheaf over the chart \(U_i = D_+(X_i)\), via the canonical map
  \(A_{(X_i)} \to \Gamma(\PP^1, U_i)\). Then the basic open of this section in the scheme
  \(\PP^1\) --- the locus of \(U_i\) where it is invertible --- is the chart overlap:
  \[
    D(t_{ij}) \;=\; U_i \cap U_j .
  \]
\end{theorem}
\begin{proof}
  \leanok
  Both sides are open subsets of \(U_i\) (the basic open of a section over \(U_i\) is
  contained in \(U_i\)), so it suffices to show that they have the same preimage under the
  identification \(U_i \cong \Spec A_{(X_i)}\), and to transport the equality forward
  along this open immersion. No stalks are involved: everything is transported through
  the affine model.

  On the one hand, the preimage of the basic open \(D(t_{ij})\) is the basic open of
  \(\Spec A_{(X_i)}\) at the ring element \(t_{ij}\): under the identification of an
  affine open with the spectrum of its section ring, scheme-level basic opens of sections
  correspond to spectrum-level basic opens of ring elements.

  On the other hand, the preimage of \(U_j = D_+(X_j)\) under the affine model
  \(\Spec A_{(X_i)} \to \PP^1\) (through which the inclusion of \(U_i\) factors,
  \ref{lem:P1_basicOpen_iota_eq}) is the basic open at the canonical degree-zero fraction
  attached to the pair \((X_i, X_j)\): a homogeneous prime avoiding \(X_i\) avoids
  \(X_j\) exactly when the degree-zero fraction \(X_j/X_i\) is invertible in the
  corresponding localization. By \ref{lem:P1_isLocalizationElem_X_eq} this fraction is
  \(t_{ij}\), so the preimage of \(U_i \cap U_j\) is again the basic open at
  \(t_{ij}\).
\end{proof}

\begin{lemma}[The Laurent generators as restrictions from the charts]
  \label{lem:P1_laurent_generators}
  \dcref{custom}
  \lean{AlgebraicGeometry.P1.overlapSectionsEquiv_symm_T,
    AlgebraicGeometry.P1.overlapSectionsEquiv_symm_T_neg,
    AlgebraicGeometry.P1.overlapSectionsEquiv_symm_algebraMap}
  \uses{def:P1_overlapSectionsEquiv, def:P1_chartCoord}
  \leanok
  Under the identification
  \(\gamma_{01} : \Gamma(\PP^1, U_0 \cap U_1) \cong k[T, T^{-1}]\) of
  \ref{def:P1_overlapSectionsEquiv}:
  \begin{enumerate}
    \item \(T\) corresponds to the restriction to the overlap of the chart coordinate
      \(t_{01}\), regarded as a section over \(U_0\);
    \item \(T^{-1}\) corresponds to the restriction to the overlap of the chart
      coordinate \(t_{10}\), regarded as a section over \(U_1\);
    \item each scalar \(r \in k\) corresponds to the section over the overlap given by
      the image of \(r\) under the algebra map \(k \to A_{(X_0X_1)}\).
  \end{enumerate}
\end{lemma}
\begin{proof}
  \leanok
  \uses{lem:P1_res_left, lem:P1_res_right}
  (1) and (2): by \ref{lem:P1_res_left} (resp.\ \ref{lem:P1_res_right}), restriction from
  the chart \(U_0\) (resp.\ \(U_1\)) to the overlap corresponds, under the
  identifications \(\gamma_0, \gamma_1, \gamma_{01}\), to \(t \mapsto T\) (resp.\
  \(t \mapsto T^{-1}\)); apply this to the sections corresponding to the polynomial
  \(t\), which are the chart coordinates. (3): \(\gamma_{01}\) is built from the
  \(k\)-algebra isomorphism \(\theta\) of \ref{def:P1_overlapAlgEquiv}, hence matches the
  two algebra maps from \(k\).
\end{proof}

\section{Finiteness of the first cohomology}

Combining the two-affine-cover cokernel bridge with the two-lattice ladder of
Chapter~\ref{chap:Algebra} along a finite morphism to the projective line yields the
finiteness of \(H^1\). Classically this computation goes back to the explicit \v{C}ech
computation of the cohomology of projective space --- for \(\PP^1\) itself, the
description of \(H^1\) as the cokernel of the two-chart restriction-difference map into
the Laurent ring, spanned by the finitely many ``negative'' monomials --- extended to a
curve by pushing forward along a finite morphism; the two-lattice ladder is exactly the
abstraction of that monomial bookkeeping.
Hartshorne's cited calculation is for projective space over a Noetherian base; it is the
model for the explicit two-chart argument below, whose hypotheses are stated independently.

\begin{theorem}[Finiteness of \(H^1\) along a finite morphism to the projective line]
  \label{thm:moduleFinite_hModule_one_of_isFinite_toP1}
  \alignmentcheck{TO CHECK: Hartshorne pages~0242--0243 treat projective space over a Noetherian
  base, while this theorem is stated for a finite map from an arbitrary scheme over a field;
  verify that the cited calculation supports this scope.}
  \lean{AlgebraicGeometry.moduleFinite_hModule_one_of_isFinite_toP1}
  \uses{def:HModule, def:moduleKSheaf, def:P1_structureMap}
  \leanok
  \dcref{custom}
  Let \(k\) be a field, \(X\) a scheme over \(\Spec k\), and \(\pi : X \to \PP^1\) a
  finite morphism over \(k\) (its composite with the structure morphism of \(\PP^1\) is
  the structure morphism of \(X\)). Then \(H^1\bigl(X, \struct{X/k}\bigr)\) is a finite
  \(k\)-module.
\end{theorem}
\begin{proof}
  \leanok
  \uses{lem:preimage_chart_affine, lem:preimage_chart_sup, lem:finite_app_overlap,
    lem:P1_chartOpen_inf, thm:TwoCover_h1CokEquiv, def:TwoCover_diff, def:TwoCover_H1Cok,
    def:P1_overlapSectionsEquiv, def:P1_chartCoord, lem:P1_laurent_generators,
    lem:appLE_overAlgebraMap, lem:P1_structureMap_appTop_awayToSection,
    thm:P1_basicOpen_chartCoord, cor:two_lattice_loc_pow}
  Write \(U_0, U_1\) for the standard charts of \(\PP^1\) and set
  \(V_i = \pi^{-1}(U_i)\).

  (a) \emph{The cover and the cokernel.} Since \(\pi\) is finite, hence affine, the
  \(V_i\) are affine opens of \(X\) (\ref{lem:preimage_chart_affine}), and they cover
  \(X\) (\ref{lem:preimage_chart_sup}). By the two-affine-cover bridge
  \ref{thm:TwoCover_h1CokEquiv},
  \[
    H^1\bigl(X, \struct{X/k}\bigr) \;\cong\;
      N \big/ \mathrm{range}(\mathrm{diff}),
    \qquad N := \Gamma(X, V_0 \cap V_1),
  \]
  with \(\mathrm{diff}(s_0, s_1) = s_0|_{V_0 \cap V_1} - s_1|_{V_0 \cap V_1}\)
  (\ref{def:TwoCover_diff}). It therefore suffices to prove that this quotient is a
  finite \(k\)-module.

  (b) \emph{The Laurent module structure.} Preimages commute with intersections, and
  \(U_0 \cap U_1 = D_+(X_0X_1)\) (\ref{lem:P1_chartOpen_inf}), so
  \(V_0 \cap V_1 = \pi^{-1}(U_0 \cap U_1)\). Make \(N\) an algebra over the Laurent
  polynomial ring via the composite
  \[
    k[T, T^{-1}] \xrightarrow{\ \gamma_{01}^{-1}\ } \Gamma(\PP^1, U_0 \cap U_1)
      \xrightarrow{\ \pi^\sharp\ } \Gamma\bigl(X, \pi^{-1}(U_0 \cap U_1)\bigr) = N
  \]
  (\ref{def:P1_overlapSectionsEquiv}). This structure is compatible with the \(k\)-module
  structure of \(N\): a scalar \(r \in k[T,T^{-1}]\) corresponds under \(\gamma_{01}\) to
  the image of \(r\) under the algebra map into the overlap ring
  (\ref{lem:P1_laurent_generators}(3)), which is the structure algebra map of \(\PP^1\)
  on \(U_0 \cap U_1\) (\ref{lem:P1_structureMap_appTop_awayToSection}), and \(\pi\) pulls
  the structure algebra map of \(\PP^1\) back to that of \(X\)
  (\ref{lem:appLE_overAlgebraMap}). Hence the \(k\)-module structure of \(N\) is the
  restriction of scalars of the \(k[T,T^{-1}]\)-structure along the canonical map
  \(k \to k[T,T^{-1}]\).

  (c) \emph{Module-finiteness over the Laurent ring.} Since \(\pi\) is finite and
  \(U_0 \cap U_1\) is an affine open of \(\PP^1\), the map
  \(\pi^\sharp : \Gamma(\PP^1, U_0 \cap U_1) \to N\) is module-finite
  (\ref{lem:finite_app_overlap}); composing with the isomorphism \(\gamma_{01}^{-1}\),
  \(N\) is a finite \(k[T,T^{-1}]\)-module.

  (d) \emph{The two lattices, and the range of the difference map.} Let
  \(Q_0, Q_1 \subseteq N\) be the images of the (\(k\)-linear) restriction maps
  \(\Gamma(X, V_0) \to N\) and \(\Gamma(X, V_1) \to N\). Then
  \(\mathrm{range}(\mathrm{diff}) = Q_0 + Q_1\): every value
  \(s_0| - s_1|\) of \(\mathrm{diff}\) is the sum of \(s_0| \in Q_0\) and
  \(-(s_1|) = (-s_1)| \in Q_1\); conversely \(s_0| = \mathrm{diff}(s_0, 0) \) and
  \(s_1| = \mathrm{diff}(0, -s_1)\) exhibit every element of \(Q_0\) and of \(Q_1\) as a
  value of \(\mathrm{diff}\) --- the sign is absorbed because \(Q_1\) is the image of a
  linear map.

  (e) \emph{Lattice stability.} Let
  \(\tau_0 = \pi^\sharp(t_{01}) \in \Gamma(X, V_0)\) and
  \(\tau_1 = \pi^\sharp(t_{10}) \in \Gamma(X, V_1)\) be the pullbacks of the chart
  coordinates, regarded as sections over the charts
  (Definition~\ref{def:P1_chartCoord}). By \ref{lem:P1_laurent_generators}(1) and the
  compatibility of \(\pi^\sharp\) with restriction, \(T\) acts on \(N\) as
  multiplication by \(\tau_0|_{V_0 \cap V_1}\); symmetrically \(T^{-1}\) acts as
  multiplication by \(\tau_1|_{V_0 \cap V_1}\) (\ref{lem:P1_laurent_generators}(2)).
  Hence \(Q_0\) is stable under multiplication by \(T\): for \(s \in \Gamma(X, V_0)\),
  \[
    T \cdot s|_{V_0 \cap V_1} \;=\; \tau_0|\cdot s| \;=\; (\tau_0 s)|_{V_0 \cap V_1}
      \;\in\; Q_0,
  \]
  and symmetrically \(Q_1\) is stable under multiplication by \(T^{-1}\).

  (f) \emph{Absorption up to powers.} The overlap is a basic open inside each chart
  preimage: pulling back the identity \(D(t_{ij}) = U_i \cap U_j\) of
  \ref{thm:P1_basicOpen_chartCoord} along \(\pi\) (basic opens of sections pull back to
  basic opens of the pulled-back sections) gives
  \[
    V_0 \cap V_1 = D(\tau_0) \subseteq V_0,
    \qquad
    V_0 \cap V_1 = D(\tau_1) \subseteq V_1 .
  \]
  Since \(V_0\) is affine and \(V_0 \cap V_1\) is its basic open defined by \(\tau_0\),
  restriction identifies \(N\) with the localization of \(\Gamma(X, V_0)\) at the powers
  of \(\tau_0\) (quasi-coherence of the structure sheaf). Hence for every \(n \in N\)
  there are \(m \ge 0\) and \(a \in \Gamma(X, V_0)\) with
  \(n \cdot (\tau_0|)^m = a|\), i.e.\ \(T^m \cdot n \in Q_0\). Symmetrically, every
  \(n \in N\) satisfies \(T^{-m} \cdot n \in Q_1\) for some \(m \ge 0\).

  (g) \emph{The ladder.} By (b), (c), (e) and (f), the module \(N\) over
  \(k[T,T^{-1}]\) with the \(k\)-submodules \(Q_0, Q_1\) satisfies all hypotheses of the
  two-lattice finiteness theorem in its localization form via powers of the generator
  (\ref{cor:two_lattice_loc_pow}); hence \(N/(Q_0 + Q_1)\) is a finite \(k\)-module. By
  (d) this quotient is \(N/\mathrm{range}(\mathrm{diff})\), and by (a) it is
  \(H^1\bigl(X, \struct{X/k}\bigr)\).
\end{proof}

\begin{theorem}[Finiteness of \(H^1\) of the curve]
  \label{thm:moduleFinite_hModule_one}
  \dcref{custom}
  \lean{AlgebraicGeometry.moduleFinite_hModule_one}
  \uses{def:HModule, def:moduleKSheaf}
  \leanok
  Let \(k\) be a field and \(C\) a smooth curve over \(k\) (proper, smooth of relative
  dimension one, geometrically irreducible). Then
  \(H^1\bigl(C, \struct{C/k}\bigr)\) is a finite \(k\)-module; equivalently, it is a
  finite-dimensional \(k\)-vector space.
\end{theorem}
\begin{proof}
  \leanok
  \uses{thm:exists_finite_toP1, thm:moduleFinite_hModule_one_of_isFinite_toP1}
  By \ref{thm:exists_finite_toP1} there is a finite morphism \(\pi : C \to \PP^1\) over
  \(k\); apply \ref{thm:moduleFinite_hModule_one_of_isFinite_toP1}.
\end{proof}

\begin{remark}[The genus is the dimension of a finite-dimensional space]
  \label{rem:genus_finite}
  \dcref{custom}
  \uses{def:genus, thm:moduleFinite_hModule_one}
  \leanok
  For a smooth curve \(C\) over \(k\), the genus of Definition~\ref{def:genus} is the
  \(k\)-dimension of \(H^1(C, \struct C)\), which by
  \ref{thm:moduleFinite_hModule_one} is a finite-dimensional \(k\)-vector space; the
  genus is therefore a well-defined natural number, attained as the cardinality of any
  basis.
\end{remark}

\section{The cocycle-glued twisted sheaf on a two-cover}
\label{sec:twisted_sheaf}

Let \(k\) be a commutative ring and \(X\) a scheme over \(\Spec k\), its structure sheaf
viewed as a sheaf of \(k\)-modules (\ref{def:moduleKSheaf}). Fix two opens
\(V_0, V_1 \subseteq X\) and a \emph{unit cocycle}: an invertible section
\(g \in \Gamma(X, V_0 \cap V_1)^{\times}\). (On a two-chart cover a \v{C}ech
\(1\)-cocycle has a single component and the cocycle condition on triple overlaps is
vacuous.) This section glues from \(g\) a sheaf \(F_g\) of \(k\)-modules --- for
\(V_0 \cup V_1 = X\), the line bundle with transition function \(g\), trivialised on each
chart --- and computes its degree-one cohomology as a two-cover cokernel, by transporting
the quasi-coherence packaging of Section~\ref{sec:qcoh_vanishing} through the chart
trivialisations.

\begin{definition}[Transport of the packaging along a sections-level trivialisation]
  \label{def:qcohOn_transport}
  \dcref{custom}
  \lean{AlgebraicGeometry.Scheme.QcohOn.ofSectionsEquiv}
  \uses{def:qcohOn}
  \leanok
  Let \(F\) and \(G\) be sheaves of \(k\)-modules on the small Zariski site of \(X\), let
  \(U\) be an open, and suppose given \(k\)-linear equivalences
  \(e_W : F(W) \cong G(W)\) for every open \(W \le U\), commuting with section
  restriction. If \(G\) carries a quasi-coherence packaging on \(U\)
  (\ref{def:qcohOn}), then so does \(F\), with the conjugated action
  \(r \cdot m := e_W^{-1}\bigl(r \cdot e_W(m)\bigr)\).
\end{definition}
\begin{proof}
  \leanok
  Each (P1) law for \(F\) follows by conjugating the corresponding law for \(G\) through
  the equivalences, and naturality under restriction combines the (P1) naturality for
  \(G\) with the compatibility of \(e\) with restriction. For (P2): given
  \(c \in F(V \cap D(g))\) on an affine \(V \le U\), denominator clearing for \(G\)
  applied to \(e(c)\) produces \(N\) and \(e' \in G(V)\) with
  \(e'|_{V \cap D(g)} = g^N \cdot e(c)\); then \(e_V^{-1}(e')\) clears the denominator
  for \(c\), since \(e\) commutes with restriction and the action is conjugated. Defect
  annihilation transports the same way.
\end{proof}

\begin{definition}[Twisted sections of a unit cocycle]
  \label{def:twistSections}
  \dcref{custom}
  \lean{AlgebraicGeometry.twistDefect, AlgebraicGeometry.twistSubmodule}
  \uses{def:moduleKSheaf}
  \leanok
  For an open \(W \subseteq X\), the \emph{defect map} of the cocycle \(g\) over \(W\) is
  the \(k\)-linear map
  \[
    \Gamma(X, W \cap V_0) \times \Gamma(X, W \cap V_1) \longrightarrow
      \Gamma(X, W \cap V_0 \cap V_1), \qquad
    (s_0, s_1) \longmapsto s_0| - g \cdot s_1|,
  \]
  where both restrictions are to \(W \cap V_0 \cap V_1\) and \(g\) acts through its
  restriction there. The module of \emph{twisted sections} \(F_g(W)\) is the kernel of
  the defect map: the pairs \((s_0, s_1)\) that match through \(g\) on the overlap.
\end{definition}

\begin{lemma}[Twisted sections form a presheaf]
  \label{lem:twistSections_presheaf}
  \dcref{custom}
  \lean{AlgebraicGeometry.twistRes, AlgebraicGeometry.twistSubmodule_res,
    AlgebraicGeometry.twistPresheaf}
  \uses{def:twistSections}
  \leanok
  For \(W' \le W\), componentwise restriction carries \(F_g(W)\) into \(F_g(W')\), and
  the resulting \(k\)-linear maps are functorial in the inclusion; so \(W \mapsto F_g(W)\)
  is a presheaf of \(k\)-modules.
\end{lemma}
\begin{proof}
  \leanok
  Restrict the matching identity \(s_0| = g \cdot s_1|\) from \(W \cap V_0 \cap V_1\) to
  \(W' \cap V_0 \cap V_1\): restriction of the structure sheaf is a ring homomorphism
  commuting with further restriction, so the restricted pair again matches.
  Functoriality is functoriality of restriction in each component.
\end{proof}

\begin{theorem}[The twisted sheaf]
  \label{thm:twistSheaf}
  \dcref{custom}
  \lean{AlgebraicGeometry.isSheaf_twistPresheaf, AlgebraicGeometry.twistSheaf}
  \uses{lem:twistSections_presheaf}
  \leanok
  The presheaf \(F_g\) of twisted sections is a sheaf of \(k\)-modules on the small
  Zariski site of \(X\).
\end{theorem}
\begin{proof}
  \leanok
  Let \(\{U_i\}\) be an open cover of \(W\) and \((s_0^i, s_1^i) \in F_g(U_i)\) sections
  agreeing on overlaps. The first components are sections of the structure sheaf over the
  cover \(\{U_i \cap V_0\}\) of \(W \cap V_0\), compatible on its overlaps; by the sheaf
  property of the structure sheaf they glue to a unique \(t_0 \in \Gamma(X, W \cap V_0)\),
  and likewise the second components glue to a unique \(t_1 \in \Gamma(X, W \cap V_1)\).
  The matching \(t_0| = g \cdot t_1|\) on \(W \cap V_0 \cap V_1\) may be checked on the
  cover \(\{U_i \cap V_0 \cap V_1\}\), where it restricts to the matching of
  \((s_0^i, s_1^i)\); by separation of the structure sheaf it holds on the union. So
  \((t_0, t_1) \in F_g(W)\) restricts to every \((s_0^i, s_1^i)\), and it is the unique
  such element because each of its components is determined by separation.
\end{proof}

\begin{definition}[Chart trivialisations of the twisted sheaf]
  \label{def:twistTriv}
  \dcref{custom}
  \lean{AlgebraicGeometry.twistTriv₀, AlgebraicGeometry.twistTriv₁}
  \uses{def:twistSections}
  \leanok
  Let \(W \le V_0\) (so \(W \cap V_0 = W\) and \(W \cap V_0 \cap V_1 = W \cap V_1\)).
  Projection to the first component is a \(k\)-linear equivalence
  \[
    \mathrm{triv}_0 : F_g(W) \;\cong\; \Gamma(X, W), \qquad (s_0, s_1) \mapsto s_0,
  \]
  with inverse \(s \mapsto (s,\ g^{-1} \cdot s|_{W \cap V_1})\). Symmetrically, for
  \(W \le V_1\), projection to the second component is a \(k\)-linear equivalence
  \(\mathrm{triv}_1 : F_g(W) \cong \Gamma(X, W)\) with inverse
  \(s \mapsto (g \cdot s|_{W \cap V_0},\ s)\).
\end{definition}
\begin{proof}
  \leanok
  For \(W \le V_0\) the matching relation of a pair \((s_0, s_1) \in F_g(W)\) reads
  \(s_0| = g \cdot s_1\) on \(W \cap V_1\), i.e.\ \(s_1 = g^{-1} \cdot s_0|\): the second
  component is rebuilt from the first, so projection and the displayed inverse are
  mutually inverse; the pair \((s, g^{-1} \cdot s|)\) matches by construction. Both maps
  are \(k\)-linear, multiplication by the fixed sections \(g^{\pm 1}\) being
  \(k\)-linear. The second chart is symmetric.
\end{proof}

\begin{lemma}[The trivialisations commute with restriction]
  \label{lem:twistTriv_res}
  \dcref{custom}
  \lean{AlgebraicGeometry.twistTriv₀_res, AlgebraicGeometry.twistTriv₁_res}
  \uses{def:twistTriv, lem:twistSections_presheaf}
  \leanok
  For \(W' \le W \le V_0\), trivialising after restricting equals restricting after
  trivialising: \(\mathrm{triv}_0(p|_{W'}) = (\mathrm{triv}_0\, p)|_{W'}\); likewise on
  the second chart.
\end{lemma}
\begin{proof}
  \leanok
  Both sides are the first (resp.\ second) component of \(p\) restricted to \(W'\);
  componentwise restriction is functorial.
\end{proof}

\begin{theorem}[The twisted sheaf is quasi-coherently packaged on each chart]
  \label{thm:twistSheaf_qcohOn}
  \dcref{custom}
  \lean{AlgebraicGeometry.twistSheaf.instQcohOn₀, AlgebraicGeometry.twistSheaf.instQcohOn₁}
  \uses{thm:twistSheaf, def:qcohOn_transport, def:twistTriv, lem:twistTriv_res,
    thm:moduleKSheaf_qcohOn}
  \leanok
  The twisted sheaf \(F_g\) carries a quasi-coherence packaging on \(V_0\) and on
  \(V_1\).
\end{theorem}
\begin{proof}
  \leanok
  The chart trivialisations give \(k\)-linear equivalences
  \(F_g(W) \cong \Gamma(X, W)\) for every \(W \le V_0\) (resp.\ \(W \le V_1\)),
  compatible with restriction (\ref{lem:twistTriv_res}); the structure sheaf is packaged
  on every open (\ref{thm:moduleKSheaf_qcohOn}); transport the packaging along the
  trivialisation (\ref{def:qcohOn_transport}).
\end{proof}

\begin{theorem}[Twisted affine vanishing]
  \label{thm:twistSheaf_affine_vanishing}
  \dcref{custom}
  \lean{AlgebraicGeometry.subsingleton_hModule'_twistSheaf_one₀,
    AlgebraicGeometry.subsingleton_hModule'_twistSheaf_one₁}
  \uses{def:HModule', thm:twistSheaf_qcohOn, thm:hModule1_vanishing_qcoh}
  \leanok
  If \(V_0\) (resp.\ \(V_1\)) is an affine open of \(X\), then
  \(H'^1(V_0, F_g) = 0\) (resp.\ \(H'^1(V_1, F_g) = 0\)).
\end{theorem}
\begin{proof}
  \leanok
  The twisted affine vanishing theorem \ref{thm:hModule1_vanishing_qcoh}, applied to the
  packaging of \ref{thm:twistSheaf_qcohOn}.
\end{proof}

\begin{definition}[The two-cover \(H^1\) carrier of the twisted sheaf]
  \label{def:twistTwoCoverH1}
  \dcref{custom}
  \lean{AlgebraicGeometry.twistTwoCoverH1}
  \uses{thm:twistSheaf, def:twoCoverSquare, thm:twistSheaf_affine_vanishing,
    thm:twoCoverH1LinearEquiv}
  \leanok
  Let \(V_0, V_1\) be affine opens with \(V_0 \cup V_1 = X\). Then the degree-one
  cohomology of the twisted sheaf is the two-cover cokernel of twisted sections: a
  \(k\)-linear isomorphism
  \[
    H^1(X, F_g) \;\cong\;
      F_g(V_0 \cap V_1) \big/ \mathrm{range}\bigl((p_0, p_1) \mapsto
        p_0|_{V_0 \cap V_1} - p_1|_{V_0 \cap V_1}\bigr),
  \]
  the restriction-difference map being that of the two-cover square
  (\ref{def:twoCoverSquare}) with coefficients in \(F_g\). (The overlap sections
  \(F_g(V_0 \cap V_1)\) are further identified with \(\Gamma(X, V_0 \cap V_1)\) by the
  chart trivialisation at \(V_0 \cap V_1 \le V_0\), \ref{def:twistTriv}.)
\end{definition}
\begin{proof}
  \leanok
  The two charts are affine, so \(H'^1(V_i, F_g) = 0\)
  (\ref{thm:twistSheaf_affine_vanishing}); apply the two-cover \(H^1\) computation with
  general coefficients (\ref{thm:twoCoverH1LinearEquiv}) to \(F = F_g\).
\end{proof}

\begin{definition}[The relative twisted sheaf]
  \label{def:relTwistSheaf}
  \dcref{custom}
  \lean{AlgebraicGeometry.relUnitCocycle, AlgebraicGeometry.relTwistSheaf}
  \uses{def:relCurve, def:relCover, lem:relCover_inf, def:twistSections, thm:twistSheaf}
  \leanok
  Let \(k\) be a field, \(C \in \Over(\Spec k)\), \(R\) a commutative \(k\)-algebra, and
  \(D = (V_0, V_1)\) an affine two-chart cover of \(C\). A unit cocycle
  \(g_k \in \Gamma(C, V_0 \cap V_1)^{\times}\) pulls back along the first projection to a
  unit cocycle on the base-changed cover of the relative curve:
  \(\mathrm{pr}_1^{\sharp}(g_k) \in
  \Gamma\bigl(C_R,\ \mathrm{pr}_1^{-1}V_0 \cap \mathrm{pr}_1^{-1}V_1\bigr)^{\times}\)
  (pullback of sections is a ring homomorphism, so it preserves units; the overlap of the
  base-changed charts is the preimage of \(V_0 \cap V_1\) by \ref{lem:relCover_inf}). The
  \emph{relative twisted sheaf} of a unit cocycle \(g\) on the base-changed cover is the
  twisted sheaf \(F_g\) on \(C_R\), a sheaf of \(R\)-modules.
\end{definition}

\begin{definition}[The relative twisted \(H^1\) carrier]
  \label{def:relTwistTwoCoverH1}
  \dcref{custom}
  \lean{AlgebraicGeometry.relTwistTwoCoverH1}
  \uses{def:relTwistSheaf, def:twistTwoCoverH1, lem:relCover_charts}
  \leanok
  In the setting of \ref{def:relTwistSheaf}, the degree-one cohomology of the relative
  twisted sheaf is the explicit two-cover cokernel of \(R\)-modules: an \(R\)-linear
  isomorphism
  \[
    H^1(C_R, F_g) \;\cong\;
      F_g\bigl(\mathrm{pr}_1^{-1}V_0 \cap \mathrm{pr}_1^{-1}V_1\bigr) \big/
        \mathrm{range}(\mathrm{diff}) .
  \]
\end{definition}
\begin{proof}
  \leanok
  The base-changed charts are affine and cover \(C_R\) (\ref{lem:relCover_charts});
  apply \ref{def:twistTwoCoverH1} over the base ring \(R\).
\end{proof}

\section{Base change of the two-term complex and of \(H^1\) on the relative curve}
\label{sec:relative_h1_base_change}

Fix a field \(k\), an object \(C \in \Over(\Spec k)\), and a homomorphism \(R \to R'\) of
commutative \(k\)-algebras --- a change of affine test rings. This section identifies the
sections of the relative curve as the free base change of the sections of \(C\), builds
the comparison morphism \(C_{R'} \to C_R\), base-changes the two-term two-cover complex
of Section~\ref{sec:mv-in-text} along \(R \to R'\) termwise, and concludes that the
degree-one cohomology of the structure sheaf of the relative curve commutes with
\emph{every} change of test ring --- no flatness of \(R \to R'\) is needed, because the
two-cover \(H^1\) is a cokernel and cokernels commute with arbitrary base change
(Chapter~\ref{chap:Algebra}).

\begin{theorem}[Sections of the relative curve are the base change of sections]
  \label{thm:sectionsBaseChange}
  \dcref{custom}
  \lean{AlgebraicGeometry.Over.sectionsBaseChange, AlgebraicGeometry.Over.sectionsBaseChange_tmul}
  \uses{def:relCurve}
  \leanok
  Let \(k\) be a field, \(C \in \Over(\Spec k)\), \(R\) a commutative \(k\)-algebra, and
  \(V \subseteq C\) a quasi-compact quasi-separated open (for instance any affine open,
  or a finite intersection or union of affine opens of a quasi-separated scheme). Write
  \(V_R := \mathrm{pr}_1^{-1} V \subseteq C_R\). Then there is a ring isomorphism
  \[
    \Gamma(C, V) \otimes_k R \;\cong\; \Gamma(C_R, V_R), \qquad
    s \otimes a \longmapsto \mathrm{pr}_1^{\sharp}(s) \cdot \mathrm{pr}_2^{\sharp}(a),
  \]
  where \(\mathrm{pr}_1^{\sharp}\) is pullback of sections along the first projection,
  and \(\mathrm{pr}_2^{\sharp}(a)\) is the pullback along the second projection of the
  global section \(a \in R = \Gamma(\Spec R, \top)\). (The \(k\)-algebra structure on
  \(\Gamma(C, V)\) is restriction of scalars along the structure morphism.)
\end{theorem}
\begin{proof}
  \leanok
  The displayed assignment is a well-defined \(k\)-algebra map \(\chi_V\) out of the
  tensor product: both pullbacks are \(k\)-algebra homomorphisms into the commutative
  ring \(\Gamma(C_R, V_R)\), and the tensor product of commutative \(k\)-algebras is
  their pushout.

  \emph{Step 1: \(V\) affine.} The open \(V_R = \mathrm{pr}_1^{-1}(V)\) of the fibre
  product \(C_R = C \times_{\Spec k} \Spec R\) is the fibre product
  \(V \times_{\Spec k} \Spec R\) of the open subscheme \(V\), a fibre product of affine
  schemes, hence the affine scheme \(\Spec(\Gamma(V) \otimes_k R)\): \(\Spec\) converts
  the pushout of commutative rings into the fibre product of affine schemes. Under this
  identification \(\chi_V\) is the canonical isomorphism.

  \emph{Step 2: \(V\) quasi-compact quasi-separated.} Choose a finite affine cover
  \(V = U_1 \cup \dots \cup U_n\) (quasi-compactness); each intersection
  \(U_i \cap U_j\) is quasi-compact (quasi-separatedness), so it in turn has a finite
  affine cover \(\{W^{ij}_l\}\). The sheaf axiom writes \(\Gamma(C, V)\) as the
  equaliser of the two restriction maps between the finite products
  \(\prod_i \Gamma(U_i)\) and \(\prod_{i,j,l} \Gamma(W^{ij}_l)\). The preimages
  \(\mathrm{pr}_1^{-1} U_i\) form an affine cover of \(V_R\) (Step 1) whose pairwise
  intersections are covered by the affine \(\mathrm{pr}_1^{-1} W^{ij}_l\) (preimage
  commutes with unions and intersections), so \(\Gamma(C_R, V_R)\) is the equaliser of
  the corresponding restriction maps for the preimage cover. Since \(R\) is a free,
  hence flat, \(k\)-module (\(k\) is a field), the functor \(- \otimes_k R\) is exact
  and commutes with finite products, so it carries the first equaliser diagram to an
  equaliser diagram. The maps \(\chi\) on the terms commute with the two restriction
  maps (pullback commutes with restriction) and are isomorphisms on the product terms by
  Step 1; hence the induced map on equalisers, which is \(\chi_V\), is an isomorphism.
\end{proof}

\begin{definition}[Relative sections as a free \(R\)-module base change]
  \label{def:relSectionsBaseChange}
  \dcref{custom}
  \lean{AlgebraicGeometry.relPullbackSection, AlgebraicGeometry.relSectionsBaseChange,
    AlgebraicGeometry.relSectionsBaseChange_tmul,
    AlgebraicGeometry.Over.sectionsBaseChange_one_tmul_overAlgebraMap}
  \uses{thm:sectionsBaseChange}
  \leanok
  In the setting of \ref{thm:sectionsBaseChange}, endow \(\Gamma(C_R, V_R)\) with the
  \(R\)-module structure \(r \cdot y := \mathrm{pr}_2^{\sharp}(r) \cdot y\) given by
  pullback along the structure morphism \(\mathrm{pr}_2 : C_R \to \Spec R\). The
  isomorphism of \ref{thm:sectionsBaseChange}, read with the tensor factors flipped, is
  an \(R\)-linear equivalence
  \[
    R \otimes_k \Gamma(C, V) \;\cong\; \Gamma(C_R, V_R), \qquad
    a \otimes s \longmapsto \mathrm{pr}_2^{\sharp}(a) \cdot \mathrm{pr}_1^{\sharp}(s),
  \]
  with \(R\) acting on the left tensor factor; in particular
  \(1 \otimes r \mapsto \mathrm{pr}_2^{\sharp}(r)\).
\end{definition}
\begin{proof}
  \leanok
  Only \(R\)-linearity needs comment: on a pure tensor,
  \(r \cdot (a \otimes s) = (ra) \otimes s\) maps to
  \(\mathrm{pr}_2^{\sharp}(ra) \cdot \mathrm{pr}_1^{\sharp}(s) =
  \mathrm{pr}_2^{\sharp}(r) \cdot \bigl(\mathrm{pr}_2^{\sharp}(a) \cdot
  \mathrm{pr}_1^{\sharp}(s)\bigr)\), because \(\mathrm{pr}_2^{\sharp}\) is a ring
  homomorphism --- and the right-hand side is \(r\) acting on the image.
\end{proof}

\begin{definition}[The comparison morphism of relative curves]
  \label{def:relCurveMap}
  \dcref{custom}
  \lean{AlgebraicGeometry.overSpecMap, AlgebraicGeometry.relCurveMap}
  \uses{def:relCurve}
  \leanok
  Let \(\varphi : R \to R'\) be a homomorphism of commutative \(k\)-algebras. Applying
  \(\Spec\) gives a morphism \(\Spec R' \to \Spec R\) over \(\Spec k\), and the induced
  morphism of fibre products
  \[
    \rho \;=\; \mathrm{id}_C \times \Spec(\varphi) \;:\; C_{R'} \longrightarrow C_R
  \]
  is the \emph{comparison morphism} between the relative curves.
\end{definition}

\begin{lemma}[Projection identities of the comparison morphism]
  \label{lem:relCurveMap_proj}
  \dcref{custom}
  \lean{AlgebraicGeometry.relCurveMap_fst, AlgebraicGeometry.relCurveMap_snd,
    AlgebraicGeometry.relCurveMap_preimage}
  \uses{def:relCurveMap}
  \leanok
  The comparison morphism satisfies
  \(\mathrm{pr}_1 \circ \rho = \mathrm{pr}_1\) and
  \(\mathrm{pr}_2 \circ \rho = \Spec(\varphi) \circ \mathrm{pr}_2\); consequently
  \(\rho^{-1}(V_R) = V_{R'}\) for every open \(V \subseteq C\).
\end{lemma}
\begin{proof}
  \leanok
  The first two identities are the compatibility of the induced morphism of fibre
  products with the two projections; the third follows from the first by taking
  preimages of \(V\).
\end{proof}

\begin{definition}[The sections comparison map]
  \label{def:relSectionsMap}
  \dcref{custom}
  \lean{AlgebraicGeometry.relSectionsMap}
  \uses{def:relCurveMap, lem:relCurveMap_proj}
  \leanok
  For \(V \subseteq C\) open, pullback of sections along \(\rho\) (allowed since
  \(\rho^{-1}(V_R) = V_{R'}\), \ref{lem:relCurveMap_proj}) is a ring homomorphism
  \[
    \rho^{\sharp} \;:\; \Gamma(C_R, V_R) \longrightarrow \Gamma(C_{R'}, V_{R'}) .
  \]
\end{definition}

\begin{lemma}[Governing identities of the sections comparison map]
  \label{lem:relSectionsMap_props}
  \dcref{custom}
  \lean{AlgebraicGeometry.relSectionsMap_pullback,
    AlgebraicGeometry.relSectionsMap_overAlgebraMap,
    AlgebraicGeometry.relSectionsMap_smul, AlgebraicGeometry.relSectionsMap_resHom}
  \uses{def:relSectionsMap, lem:relCurveMap_proj, def:relSectionsBaseChange}
  \leanok
  The sections comparison map \(\rho^{\sharp}\):
  \begin{enumerate}
    \item commutes with the curve pullbacks:
      \(\rho^{\sharp}(\mathrm{pr}_1^{\sharp} s) = \mathrm{pr}_1^{\sharp} s\) for
      \(s \in \Gamma(C, V)\);
    \item intertwines the structure pullbacks:
      \(\rho^{\sharp}(\mathrm{pr}_2^{\sharp} r) = \mathrm{pr}_2^{\sharp}(\varphi(r))\)
      for \(r \in R\);
    \item hence is semilinear along \(\varphi\):
      \(\rho^{\sharp}(r \cdot y) = \varphi(r) \cdot \rho^{\sharp}(y)\) for the module
      structures of \ref{def:relSectionsBaseChange};
    \item commutes with restriction along \(W \le V\).
  \end{enumerate}
\end{lemma}
\begin{proof}
  \leanok
  (1) is pullback along \(\mathrm{pr}_1 \circ \rho = \mathrm{pr}_1\)
  (\ref{lem:relCurveMap_proj}); (2) is pullback along
  \(\mathrm{pr}_2 \circ \rho = \Spec(\varphi) \circ \mathrm{pr}_2\), and the pullback of
  the global section \(r\) along \(\Spec(\varphi)\) is \(\varphi(r)\); (3) is (2)
  together with multiplicativity of \(\rho^{\sharp}\); (4) is functoriality of pullback
  of sections with respect to inclusions of opens.
\end{proof}

\begin{theorem}[CBC-1: termwise base change of relative sections]
  \label{thm:relTermBaseChange}
  \dcref{custom}
  \lean{AlgebraicGeometry.relTermBaseChange, AlgebraicGeometry.relTermBaseChange_tmul}
  \uses{def:relSectionsBaseChange, lem:relSectionsMap_props}
  \leanok
  Let \(R \to R'\) be a homomorphism of test rings and \(V \subseteq C\) a quasi-compact
  quasi-separated open. There is an \(R'\)-linear equivalence
  \[
    R' \otimes_R \Gamma(C_R, V_R) \;\cong\; \Gamma(C_{R'}, V_{R'}), \qquad
    r' \otimes y \longmapsto r' \cdot \rho^{\sharp}(y) .
  \]
\end{theorem}
\begin{proof}
  \leanok
  Compose the free identifications of \ref{def:relSectionsBaseChange} over \(R\) and
  over \(R'\) with the tensor cancellation
  \(R' \otimes_R (R \otimes_k M) \cong R' \otimes_k M\):
  \[
    R' \otimes_R \Gamma(C_R, V_R) \;\cong\; R' \otimes_R \bigl(R \otimes_k \Gamma(C, V)\bigr)
      \;\cong\; R' \otimes_k \Gamma(C, V) \;\cong\; \Gamma(C_{R'}, V_{R'}) .
  \]
  For the formula, both sides are additive in \(y\) and \(R'\)-linear in \(r'\), so it
  suffices to treat \(y = \mathrm{pr}_2^{\sharp}(a) \cdot \mathrm{pr}_1^{\sharp}(s)\)
  (the images of pure tensors, which span). Chasing \(r' \otimes y\) through the
  composite gives \(\mathrm{pr}_2^{\sharp}\bigl(r' \varphi(a)\bigr) \cdot
  \mathrm{pr}_1^{\sharp}(s) = r' \cdot \bigl(\mathrm{pr}_2^{\sharp}(\varphi(a)) \cdot
  \mathrm{pr}_1^{\sharp}(s)\bigr)\), which by (1) and (2) of
  \ref{lem:relSectionsMap_props} is \(r' \cdot \rho^{\sharp}(y)\).
\end{proof}

\begin{lemma}[CBC-1: the two-cover complex base-changes]
  \label{lem:relDiffBaseChange}
  \dcref{custom}
  \lean{AlgebraicGeometry.relDiffDomBaseChange, AlgebraicGeometry.relDiffCodBaseChange,
    AlgebraicGeometry.relDiffBaseChange, AlgebraicGeometry.relDiffBaseChange_range}
  \uses{thm:relTermBaseChange, def:TwoCover_diff, def:relCover, lem:relCover_charts,
    lem:relCover_inf, lem:relSectionsMap_props}
  \leanok
  Let \(D = (V_0, V_1)\) be an affine two-chart cover of \(C\), and write
  \(\mathrm{diff}_R\) for the two-cover restriction-difference map
  (\ref{def:TwoCover_diff}) of \(C_R\) with its base-changed cover
  (\ref{def:relCover}). Under the term equivalences of \ref{thm:relTermBaseChange} ---
  the product of the two chart equivalences on the domain, the overlap equivalence on the
  codomain --- the base change \(\mathrm{diff}_R \otimes_R R'\) corresponds to
  \(\mathrm{diff}_{R'}\): the evident square commutes. In particular the overlap term
  equivalence carries \(\mathrm{range}(\mathrm{diff}_R \otimes_R R')\) onto
  \(\mathrm{range}(\mathrm{diff}_{R'})\).
\end{lemma}
\begin{proof}
  \leanok
  Both composites are additive, so it suffices to compare them on
  \(r' \otimes (s_0, s_1)\). Along the top and right, the result is
  \(r' \cdot \rho^{\sharp}\bigl(s_0|_{\cap} - s_1|_{\cap}\bigr)\); along the left and
  bottom, it is \(\mathrm{diff}_{R'}\bigl(r' \cdot \rho^{\sharp} s_0,\ r' \cdot
  \rho^{\sharp} s_1\bigr) = \bigl(r' \cdot \rho^{\sharp} s_0\bigr)|_{\cap} -
  \bigl(r' \cdot \rho^{\sharp} s_1\bigr)|_{\cap}\). These agree because
  \(\rho^{\sharp}\) commutes with restriction (\ref{lem:relSectionsMap_props}(4)) and
  the \(R'\)-action commutes with restriction (the structure pullback restricts to the
  structure pullback). The range statement follows because the domain equivalence is
  surjective.
\end{proof}

\begin{theorem}[CBC-2, cokernel form: the two-cover \(H^1\) carrier base-changes]
  \label{thm:relH1CokBaseChange}
  \dcref{custom}
  \lean{AlgebraicGeometry.relH1CokBaseChange, AlgebraicGeometry.relH1CokBaseChange_tmul_mk}
  \uses{thm:quot_range_baseChange, lem:relDiffBaseChange, def:TwoCover_H1Cok}
  \leanok
  In the setting of \ref{lem:relDiffBaseChange} there is an \(R'\)-linear equivalence
  \[
    R' \otimes_R \Bigl(\Gamma\bigl(C_R, \cap\bigr)\big/\mathrm{range}(\mathrm{diff}_R)\Bigr)
    \;\cong\;
    \Gamma\bigl(C_{R'}, \cap\bigr)\big/\mathrm{range}(\mathrm{diff}_{R'}),
    \qquad r' \otimes [s] \longmapsto [\,r' \cdot \rho^{\sharp} s\,],
  \]
  where \(\cap\) denotes the overlap of the respective base-changed covers
  (\ref{def:TwoCover_H1Cok}). No hypothesis on \(R \to R'\) is needed.
\end{theorem}
\begin{proof}
  \leanok
  Base change commutes with cokernels (\ref{thm:quot_range_baseChange}, with
  \(N = \Gamma(C_R, \cap)\), \(f = \mathrm{diff}_R\), \(S = R'\)):
  \(R' \otimes_R \bigl(N/\mathrm{range}\,f\bigr) \cong (R' \otimes_R N)/
  \mathrm{range}(f \otimes_R R')\). Then pass along the overlap term equivalence, which
  identifies \(\mathrm{range}(\mathrm{diff}_R \otimes_R R')\) with
  \(\mathrm{range}(\mathrm{diff}_{R'})\) (\ref{lem:relDiffBaseChange}), hence descends to
  the quotients. The composite sends \(r' \otimes [s]\) to
  \([\,r' \cdot \rho^{\sharp} s\,]\) by \ref{thm:relTermBaseChange}.
\end{proof}

\begin{theorem}[CBC-2: \(H^1\) of the relative curve commutes with change of the test ring]
  \label{thm:relH1BaseChange}
  \lean{AlgebraicGeometry.relH1BaseChange}
  \uses{def:relTwoCoverH1, thm:relH1CokBaseChange}
  \leanok
  \dcref{custom}
  For every homomorphism \(R \to R'\) of commutative \(k\)-algebras there is an
  \(R'\)-linear equivalence
  \[
    R' \otimes_R H^1\bigl(C_R, \struct{C_R/R}\bigr) \;\cong\;
      H^1\bigl(C_{R'}, \struct{C_{R'}/R'}\bigr) :
  \]
  degree-one cohomology of the structure sheaf of the relative curve commutes with
  \emph{every} base change of the test ring. No flatness of \(R \to R'\) is required:
  the two-cover \(H^1\) is the cokernel of a two-term complex whose terms base-change
  freely, and cokernels commute with arbitrary base change.
\end{theorem}
\begin{proof}
  \leanok
  Conjugate the cokernel equivalence of \ref{thm:relH1CokBaseChange} by the relative
  two-cover identifications \(H^1(C_R, \struct{C_R/R}) \cong
  \Gamma(C_R, \cap)/\mathrm{range}(\mathrm{diff}_R)\) of \ref{def:relTwoCoverH1}, over
  \(R\) and over \(R'\).
\end{proof}

\begin{definition}[Base change of the unit cocycle]
  \label{def:relCocycleBaseChange}
  \dcref{custom}
  \lean{AlgebraicGeometry.relCocycleBaseChange}
  \uses{def:relTwistSheaf, def:relSectionsMap, lem:relCurveMap_proj}
  \leanok
  For a unit cocycle \(g\) on the base-changed cover of \(C_R\), the image
  \(\rho^{\sharp}(g)\) is a unit cocycle on the base-changed cover of \(C_{R'}\): the
  ring homomorphism \(\rho^{\sharp}\) preserves units, and the overlaps correspond under
  \(\rho\) (\ref{lem:relCurveMap_proj}).
\end{definition}

\begin{definition}[CBC-1 for the twisted sheaf: term equivalences]
  \label{def:relTwistTermBaseChange}
  \dcref{custom}
  \lean{AlgebraicGeometry.relTwistTermBaseChange₀, AlgebraicGeometry.relTwistTermBaseChange₁,
    AlgebraicGeometry.relTwistOverlapBaseChange}
  \uses{def:relTwistSheaf, def:relCocycleBaseChange, def:twistTriv, thm:relTermBaseChange,
    lem:relDiffBaseChange}
  \leanok
  Let \(g\) be a unit cocycle on the base-changed cover of \(C_R\) and
  \(g' = \rho^{\sharp}(g)\) its base change (\ref{def:relCocycleBaseChange}). For each
  term of the twisted two-cover complex --- the two charts and their overlap --- the
  composite of the chart trivialisation over \(R\) (\ref{def:twistTriv}), the
  structure-sheaf term base change of \ref{thm:relTermBaseChange} (at the overlap: the
  codomain equivalence of \ref{lem:relDiffBaseChange}), and the inverse chart
  trivialisation over \(R'\) is an \(R'\)-linear equivalence
  \[
    R' \otimes_R F_g(\text{term over } R) \;\cong\; F_{g'}(\text{term over } R') .
  \]
  On the overlap, the first-chart trivialisation (at
  \(\mathrm{pr}_1^{-1}V_0 \cap \mathrm{pr}_1^{-1}V_1 \le \mathrm{pr}_1^{-1}V_0\)) is
  used.
\end{definition}

\section{The base-field transition and its sections base change}
\label{sec:base_field_transition}

Fix a field \(k\) and an object \(C \in \Over(\Spec k)\). For a field extension \(K/k\)
write \(C_K := (C \otimes \Spec K)_{\mathrm{left}}\) for the base-changed curve, a scheme
over \(\Spec K\) through the second projection. A \(k\)-algebra homomorphism
\(\varphi : K_1 \to K_2\) between field extensions of \(k\) induces the
\emph{base-field transition}
\[
  \pi \;=\; \bigl(C \lhd \Spec(\varphi)\bigr)_{\mathrm{left}} \;:\; C_{K_2}
  \longrightarrow C_{K_1},
\]
the whiskering of \(C\) with \(\Spec\varphi\) between the test objects
(\ref{def:overSpecMap}). This section provides the transition's geometric toolkit and
identifies the sections of \(C_{K_2}\) over the preimage of an affine open with the base
change of the sections of \(C_{K_1}\).

A design point governs everything that follows: the ``base-field shuffle''
\(C_{K_1} \otimes_{K_1} K_2 \cong C_{K_2}\) is \emph{never constructed as an isomorphism
of schemes}. All statements are made against the transition square below, which exhibits
\(\pi\) as a base change of \(\Spec\varphi\) directly. In particular, the sections base
change of this section is deliberately \emph{not} obtained by instantiating
Theorem~\ref{thm:sectionsBaseChange} (whose tensor factor is anchored at the base field
\(k\)) over the base \(K_1\) and transporting along a shuffle isomorphism --- the short
pushout argument is re-run against the pasted square instead, with the
\(K_1\)-algebra structure on sections carried by the second projection.

\begin{lemma}[The product bridge]
  \label{lem:over_isPullback_left}
  \dcref{custom}
  \lean{AlgebraicGeometry.Over.isPullback_left}
  \leanok
  Let \(S\) be a scheme and \(X, T \in \Over(S)\). The underlying scheme of the
  cartesian-monoidal product, with its two projections, is the fibre product of the
  structure morphisms:
  \[
    \begin{array}{ccc}
      (X \otimes T)_{\mathrm{left}} & \xrightarrow{\ \mathrm{fst}\ } & X_{\mathrm{left}} \\
      \downarrow \mathrm{snd} & & \downarrow X.\mathrm{hom} \\
      T_{\mathrm{left}} & \xrightarrow{\ T.\mathrm{hom}\ } & S
    \end{array}
  \]
  is a pullback square of schemes. This is the single point where the product of the over
  category is unfolded to a scheme-level pullback.
\end{lemma}
\begin{proof}
  \leanok
  The cartesian-monoidal product of the over category is constructed as the chosen pullback
  of the structure morphisms, with the two projections as its legs; the square is that
  chosen pullback square.
\end{proof}

\begin{theorem}[The pasted square of a morphism of test objects]
  \label{thm:isPullback_whiskerLeft_left}
  \dcref{custom}
  \lean{AlgebraicGeometry.Over.isPullback_whiskerLeft_left}
  \uses{lem:over_isPullback_left}
  \leanok
  Let \(S\) be a scheme, \(X \in \Over(S)\), and \(t : T' \to T\) a morphism of
  \(\Over(S)\). The square
  \[
    \begin{array}{ccc}
      (X \otimes T')_{\mathrm{left}} & \xrightarrow{(X \lhd t)_{\mathrm{left}}} &
        (X \otimes T)_{\mathrm{left}} \\
      \downarrow \mathrm{pr}_2' & & \downarrow \mathrm{pr}_2 \\
      T'_{\mathrm{left}} & \xrightarrow{\ t_{\mathrm{left}}\ } & T_{\mathrm{left}}
    \end{array}
  \]
  (verticals the second projections) is a pullback of schemes. This generalises the
  base-change square of the product (\ref{thm:isPullback_whiskerLeft_snd}) from the base
  \(\Spec k\) to an arbitrary base \(S\); it is stated for an arbitrary morphism of test
  objects so that later campaigns can re-instantiate it without re-proving.
\end{theorem}
\begin{proof}
  \leanok
  Paste the two cartesian squares expressing \(X \otimes T\) and \(X \otimes T'\) as
  products over \(S\). The first-projection triangle
  \((X \lhd t)_{\mathrm{left}} \circ \mathrm{fst} = \mathrm{fst}'\) and the base identity
  \(t_{\mathrm{left}} \circ (T.\mathrm{hom}) = T'.\mathrm{hom}\) exhibit the outer
  rectangle over \(X_{\mathrm{left}} \times S\) as the product square of
  \(X \otimes T'\); the right-cancellation property of pullback squares applied to the
  product square of \(X \otimes T\) then makes the displayed square a pullback.
\end{proof}

\begin{theorem}[The base-field transition square]
  \label{thm:isPullback_baseFieldTransition}
  \dcref{custom}
  \lean{AlgebraicGeometry.isPullback_baseFieldTransition}
  \uses{thm:isPullback_whiskerLeft_left, def:overSpecMap}
  \leanok
  Let \(k\) be a commutative ring, \(C \in \Over(\Spec k)\), and
  \(\varphi : K_1 \to K_2\) a \(k\)-algebra homomorphism. The square
  \[
    \begin{array}{ccc}
      C_{K_2} & \xrightarrow{\ \pi\ } & C_{K_1} \\
      \downarrow \mathrm{pr}_2 & & \downarrow \mathrm{pr}_2 \\
      \Spec K_2 & \xrightarrow{\Spec \varphi} & \Spec K_1
    \end{array}
  \]
  is a pullback of schemes: the base-field transition is the base change of
  \(\Spec\varphi\) along the second projection of the \(K_1\)-curve.
\end{theorem}
\begin{proof}
  \leanok
  Instantiate the pasted square (\ref{thm:isPullback_whiskerLeft_left}) at the morphism
  of test objects \(\Spec\varphi : \Spec K_2 \to \Spec K_1\) over \(\Spec k\)
  (\ref{def:overSpecMap}).
\end{proof}

The pasted square above compares two tests over one base. The next squares cross a
morphism of \emph{bases} \(f : S' \to S\): the slice product of the base change
\(f^{*}X\) with a test over \(S'\) agrees, as a scheme, with the slice product of \(X\)
itself with the test pushed forward to \(S\). Throughout, \(f^{*}X \in \Over(S')\)
denotes the base change of \(X \in \Over(S)\) along \(f\) --- its underlying scheme is
the fibre product \(X_{\mathrm{left}} \times_S S'\), its structure morphism the second
projection --- and \(f_{!}T \in \Over(S)\) denotes a test \(T \in \Over(S')\) regarded
over \(S\) through \(f\): the same underlying scheme, with structure morphism
\(f \circ T.\mathrm{hom}\).

\begin{theorem}[The cross-base pasted square]
  \label{thm:over_isPullback_crossBase}
  \dcref{custom}
  \lean{AlgebraicGeometry.Over.isPullback_crossBase}
  \uses{lem:over_isPullback_left}
  \leanok
  Let \(f : S' \to S\) be a morphism of schemes, \(X \in \Over(S)\), and
  \(T \in \Over(S')\). The square
  \[
    \begin{array}{ccc}
      (f^{*}X \otimes T)_{\mathrm{left}} &
        \xrightarrow{\ \mathrm{pr}_1 \circ\, \mathrm{fst}\ } & X_{\mathrm{left}} \\
      \downarrow \mathrm{snd} & & \downarrow X.\mathrm{hom} \\
      T_{\mathrm{left}} & \xrightarrow{\ f \circ T.\mathrm{hom}\ } & S
    \end{array}
  \]
  is a pullback of schemes; here the top edge is the first slice projection
  \(\mathrm{fst} : (f^{*}X \otimes T)_{\mathrm{left}} \to (f^{*}X)_{\mathrm{left}}
  = X_{\mathrm{left}} \times_S S'\) followed by the first projection of the fibre
  product, and the bottom edge \(f \circ T.\mathrm{hom}\) is the structure morphism of
  \(f_{!}T\). Thus the slice product \(f^{*}X \otimes T\), formed over \(S'\), is at the
  same time a pullback of the \emph{cross-base} cospan
  \(X_{\mathrm{left}} \to S \leftarrow T_{\mathrm{left}}\).
\end{theorem}
\begin{proof}
  \leanok
  Paste horizontally the slice-product square of \(f^{*}X\) and \(T\) over \(S'\)
  (\ref{lem:over_isPullback_left}) --- horizontals \(\mathrm{fst}\) and
  \(T.\mathrm{hom}\), verticals \(\mathrm{snd}\) and the structure morphism of
  \(f^{*}X\) --- with the defining base-change square of \(f^{*}X\),
  \[
    \begin{array}{ccc}
      X_{\mathrm{left}} \times_S S' & \xrightarrow{\ \mathrm{pr}_1\ } &
        X_{\mathrm{left}} \\
      \downarrow \mathrm{pr}_2 & & \downarrow X.\mathrm{hom} \\
      S' & \xrightarrow{\ f\ } & S
    \end{array}
  \]
  whose left vertical \(\mathrm{pr}_2\) \emph{is} that structure morphism. The two
  pullback squares share this middle edge, and a horizontal paste of two pullback
  squares is a pullback; the outer rectangle is the displayed square, its top the
  composite \(\mathrm{pr}_1 \circ \mathrm{fst}\) and its bottom the composite
  \(f \circ T.\mathrm{hom}\).
\end{proof}

\begin{definition}[The cross-base comparison isomorphism]
  \label{def:crossBaseIso}
  \dcref{custom}
  \lean{AlgebraicGeometry.Over.crossBaseIso}
  \uses{thm:over_isPullback_crossBase, lem:over_isPullback_left}
  \leanok
  In the situation of \ref{thm:over_isPullback_crossBase}, the \emph{cross-base
  comparison isomorphism} is the isomorphism of schemes
  \[
    (f^{*}X \otimes T)_{\mathrm{left}} \;\cong\; (X \otimes f_{!}T)_{\mathrm{left}}
  \]
  identifying the two pullbacks of the common cospan
  \(X_{\mathrm{left}} \to S \leftarrow T_{\mathrm{left}}\): the left-hand side by the
  cross-base pasted square (\ref{thm:over_isPullback_crossBase}), the right-hand side
  by the product bridge for \(X\) and \(f_{!}T\) over \(S\)
  (\ref{lem:over_isPullback_left}). It is compatible with the projections: composed
  with the first projection of the right-hand side it is the top edge
  \(\mathrm{pr}_1 \circ \mathrm{fst}\) of the pasted square, and composed with the
  second projection it is \(\mathrm{snd}\).
\end{definition}
\begin{proof}
  \leanok
  Two pullbacks of one cospan are canonically isomorphic: each limiting cone factors
  uniquely through the other, and the two factorizations compose to endomorphisms of
  limiting cones commuting with the projections, hence to the identities. The stated
  projection compatibilities are the defining property of the factorization.
\end{proof}

\begin{theorem}[Naturality of the cross-base isomorphism in the test]
  \label{thm:crossBaseIso_naturality}
  \dcref{custom}
  \lean{AlgebraicGeometry.Over.crossBaseIso_naturality}
  \uses{def:crossBaseIso, lem:over_isPullback_left}
  \leanok
  For a morphism \(t : T' \to T\) of \(\Over(S')\), the comparison isomorphisms of
  \ref{def:crossBaseIso} intertwine the two whiskered morphisms: the square
  \[
    \begin{array}{ccc}
      (f^{*}X \otimes T')_{\mathrm{left}} &
        \xrightarrow{(f^{*}X \lhd t)_{\mathrm{left}}} &
        (f^{*}X \otimes T)_{\mathrm{left}} \\
      \downarrow \cong & & \downarrow \cong \\
      (X \otimes f_{!}T')_{\mathrm{left}} &
        \xrightarrow{(X \lhd f_{!}t)_{\mathrm{left}}} &
        (X \otimes f_{!}T)_{\mathrm{left}}
    \end{array}
  \]
  commutes, where \(f_{!}t\) is \(t\) regarded as a morphism \(f_{!}T' \to f_{!}T\)
  over \(S\).
\end{theorem}
\begin{proof}
  \leanok
  Two morphisms into the pullback \((X \otimes f_{!}T)_{\mathrm{left}}\)
  (\ref{lem:over_isPullback_left}) are equal once they agree after composing with both
  projections. With the first projection: whiskering leaves the first factor untouched,
  and the comparison isomorphisms identify the first projections with the top edges of
  the pasted squares (\ref{def:crossBaseIso}), so both composites equal the top edge
  \(\mathrm{pr}_1 \circ \mathrm{fst}'\) of the square of \(T'\). With the second
  projection: whiskering acts on the second factor by \(t\) --- whose underlying
  morphism of schemes is also that of \(f_{!}t\) --- and the comparison isomorphisms
  identify the second projections, so both composites equal
  \(t_{\mathrm{left}} \circ \mathrm{snd}'\). Hence the square commutes.
\end{proof}

\begin{theorem}[The cross-base square at a base-changed curve]
  \label{thm:crossBase_baseChange}
  \dcref{custom}
  \lean{AlgebraicGeometry.isPullback_crossBase, AlgebraicGeometry.crossBaseIso,
    AlgebraicGeometry.crossBaseIso_naturality}
  \uses{thm:over_isPullback_crossBase, def:crossBaseIso, thm:crossBaseIso_naturality,
    def:baseChange}
  \leanok
  Let \(k \to L\) be a field extension, \(\sigma : \Spec L \to \Spec k\) the induced
  morphism of spectra, and \(C \in \Over(\Spec k)\), with base change
  \(\sigma^{*}C \in \Over(\Spec L)\) along \(L/k\) (\ref{def:baseChange}). Then for
  every test object \(T \in \Over(\Spec L)\): the slice product
  \((\sigma^{*}C \otimes T)_{\mathrm{left}}\), formed over \(\Spec L\), is a pullback
  of \(C.\mathrm{hom}\) against the structure morphism of \(\sigma_{!}T\), and the
  comparison isomorphism of schemes
  \[
    (\sigma^{*}C \otimes T)_{\mathrm{left}} \;\cong\;
      (C \otimes \sigma_{!}T)_{\mathrm{left}}
  \]
  is natural in \(T\).
\end{theorem}
\begin{proof}
  \leanok
  Base change along \(L/k\) is the slice pullback along \(\sigma\)
  (\ref{def:baseChange}), so this is \ref{thm:over_isPullback_crossBase},
  \ref{def:crossBaseIso} and \ref{thm:crossBaseIso_naturality} instantiated at
  \(f = \sigma\) and \(X = C\).
\end{proof}

\begin{definition}[Pushing an affine test down a field extension]
  \label{def:mapOverSpecIso}
  \dcref{custom}
  \lean{AlgebraicGeometry.mapOverSpecIso}
  \leanok
  Let \(k \to L\) be a field extension, \(\sigma : \Spec L \to \Spec k\) the induced
  morphism, and \(A\) an \(L\)-algebra whose \(k\)-algebra structure is the restriction
  through \(L\) (a scalar tower \(k \to L \to A\)). Write \(\Spec A / \Spec F\) for the
  affine test attached to an \(F\)-algebra \(A\): the object \(\Spec A\) of
  \(\Over(\Spec F)\) with structure morphism \(\Spec\) of the algebra map. Then pushing
  the affine test of \(A\) over \(L\) forward along \(\sigma\) gives the affine test of
  \(A\) over \(k\): there is an isomorphism
  \[
    \sigma_{!}\bigl(\Spec A / \Spec L\bigr) \;\cong\; \Spec A / \Spec k
    \quad \text{in } \Over(\Spec k),
  \]
  whose underlying morphism of schemes is the identity of \(\Spec A\).
\end{definition}
\begin{proof}
  \leanok
  The push-forward keeps the underlying scheme \(\Spec A\) and composes the structure
  morphism with \(\sigma\). The algebra map \(k \to A\) is the composite
  \(k \to L \to A\) (the tower identity), so by contravariant functoriality of
  \(\Spec\) the composed structure morphism
  \(\sigma \circ \Spec(L \to A)\) is \(\Spec\) of the algebra map \(k \to A\) --- which
  is exactly the structure morphism of the affine test of \(A\) over \(k\). Hence the
  identity of \(\Spec A\) is an isomorphism between the two objects of
  \(\Over(\Spec k)\).
\end{proof}

The cross-base comparison and the affine matching combine at an affine test: for an
algebra \(A\) in a scalar tower \(k \to L \to A\), the base-changed curve tensored with
the affine test of \(A\) over \(L\) and the original curve tensored with the affine test
of \(A\) over \(k\) have canonically isomorphic underlying schemes --- a
\emph{same-carrier} comparison, compatible with the two projections onto the common
carrier \(\Spec A\) and natural in \(A\).

\begin{definition}[The same-carrier comparison at an affine test]
  \label{def:crossBaseAffineIso}
  \dcref{custom}
  \lean{AlgebraicGeometry.crossBaseAffineIso, AlgebraicGeometry.crossBaseAffineIso_hom}
  \uses{thm:crossBase_baseChange, def:mapOverSpecIso, def:baseChange}
  \leanok
  Let \(k \to L\) be a field extension, \(\sigma : \Spec L \to \Spec k\) the induced
  morphism, \(C \in \Over(\Spec k)\) with base change \(\sigma^{*}C \in \Over(\Spec L)\)
  (\ref{def:baseChange}), and \(A\) an algebra in a scalar tower \(k \to L \to A\). The
  \emph{same-carrier comparison} is the isomorphism of schemes
  \[
    \bigl(\sigma^{*}C \otimes (\Spec A/\Spec L)\bigr)_{\mathrm{left}}
      \;\cong\; \bigl(C \otimes (\Spec A/\Spec k)\bigr)_{\mathrm{left}}
  \]
  (the left slice product formed over \(\Spec L\), the right one over \(\Spec k\))
  defined as the composite of the cross-base comparison at the test
  \(T = \Spec A/\Spec L\) (\ref{thm:crossBase_baseChange}),
  \[
    \bigl(\sigma^{*}C \otimes (\Spec A/\Spec L)\bigr)_{\mathrm{left}} \;\cong\;
    \bigl(C \otimes \sigma_{!}(\Spec A/\Spec L)\bigr)_{\mathrm{left}},
  \]
  with the underlying morphism of \(C\) whiskered with the affine matching
  \(\sigma_{!}(\Spec A/\Spec L) \cong \Spec A/\Spec k\) (\ref{def:mapOverSpecIso}).
\end{definition}
\begin{proof}
  \leanok
  Both factors are isomorphisms of schemes: the first by
  \ref{thm:crossBase_baseChange}; the second because whiskering with \(C\) is a functor
  on \(\Over(\Spec k)\) and taking underlying schemes is a functor to schemes, and
  functors carry the isomorphism of \ref{def:mapOverSpecIso} to an isomorphism. The
  composite of the two is therefore an isomorphism, and by construction its underlying
  morphism is the composite of the underlying morphisms of the two factors.
\end{proof}

\begin{lemma}[The same-carrier comparison lies over \(\Spec A\)]
  \label{lem:crossBaseAffineIso_snd}
  \dcref{custom}
  \lean{AlgebraicGeometry.crossBaseAffineIso_hom_snd,
    AlgebraicGeometry.crossBaseAffineIso_inv_snd}
  \uses{def:crossBaseAffineIso, def:crossBaseIso, def:mapOverSpecIso}
  \leanok
  Both slice products of \ref{def:crossBaseAffineIso} project onto the same scheme
  \(\Spec A\) by their second projections, and the same-carrier comparison \(c\)
  commutes with them:
  \[
    \mathrm{snd}_{k} \circ c \;=\; \mathrm{snd}_{L},
    \qquad
    \mathrm{snd}_{L} \circ c^{-1} \;=\; \mathrm{snd}_{k},
  \]
  where \(\mathrm{snd}_{L}\) and \(\mathrm{snd}_{k}\) denote the underlying morphisms of
  the second projections of \(\sigma^{*}C \otimes (\Spec A/\Spec L)\) and of
  \(C \otimes (\Spec A/\Spec k)\). In particular \(c\) is an isomorphism of schemes over
  \(\Spec A\).
\end{lemma}
\begin{proof}
  \leanok
  Write \(c = w \circ b\) with \(b\) the cross-base comparison and \(w\) the whiskered
  affine matching (\ref{def:crossBaseAffineIso}). Whiskering acts on the second factor
  by the whiskered morphism: \(\mathrm{snd}_{k} \circ w\) is the underlying morphism of
  the affine matching composed with the second projection of
  \(C \otimes \sigma_{!}(\Spec A/\Spec L)\), and the underlying morphism of the affine
  matching is the identity of \(\Spec A\) (\ref{def:mapOverSpecIso}). The cross-base
  comparison identifies the second projections (\ref{def:crossBaseIso}): the second
  projection of \(C \otimes \sigma_{!}(\Spec A/\Spec L)\) composed with \(b\) is
  \(\mathrm{snd}_{L}\). Chaining, \(\mathrm{snd}_{k} \circ c = \mathrm{snd}_{L}\);
  composing this identity with \(c^{-1}\) on the right gives
  \(\mathrm{snd}_{L} \circ c^{-1} = \mathrm{snd}_{k}\). Finally, the two second
  projections are the structure morphisms of the two products regarded over
  \(\Spec A\), so \(c\) is an isomorphism over \(\Spec A\).
\end{proof}

\begin{theorem}[Naturality of the same-carrier comparison in the algebra]
  \label{thm:crossBaseAffineIso_naturality}
  \dcref{custom}
  \lean{AlgebraicGeometry.crossBaseAffineIso_naturality}
  \uses{def:crossBaseAffineIso, thm:crossBase_baseChange, thm:crossBaseIso_naturality,
    def:mapOverSpecIso, def:overSpecMap}
  \leanok
  Let \(A, B\) be algebras in scalar towers \(k \to L \to A\), \(k \to L \to B\) and
  \(f : A \to B\) a homomorphism of \(L\)-algebras; through the towers, \(f\) is then
  also a homomorphism of \(k\)-algebras. Writing \(c_A, c_B\) for the same-carrier
  comparisons (\ref{def:crossBaseAffineIso}) and \(\Spec_L f : \Spec B/\Spec L \to
  \Spec A/\Spec L\), \(\Spec_k f : \Spec B/\Spec k \to \Spec A/\Spec k\) for the
  morphisms of affine tests induced by \(f\) over the two bases (\ref{def:overSpecMap}),
  the square
  \[
    \begin{array}{ccc}
      \bigl(\sigma^{*}C \otimes (\Spec B/\Spec L)\bigr)_{\mathrm{left}} &
        \xrightarrow{\ c_B\ } &
        \bigl(C \otimes (\Spec B/\Spec k)\bigr)_{\mathrm{left}} \\
      \downarrow (\sigma^{*}C \lhd \Spec_L f)_{\mathrm{left}} & &
        \downarrow (C \lhd \Spec_k f)_{\mathrm{left}} \\
      \bigl(\sigma^{*}C \otimes (\Spec A/\Spec L)\bigr)_{\mathrm{left}} &
        \xrightarrow{\ c_A\ } &
        \bigl(C \otimes (\Spec A/\Spec k)\bigr)_{\mathrm{left}}
    \end{array}
  \]
  commutes.
\end{theorem}
\begin{proof}
  \leanok
  \emph{The affine matchings intertwine the two \(\Spec f\).} The push-forward
  \(\sigma_{!}(\Spec_L f)\) and the morphism \(\Spec_k f\) have the same underlying
  morphism of schemes \(\Spec f\), and the affine matchings
  \(\mathrm{match}_A, \mathrm{match}_B\) have underlying morphism the identity of the
  respective carrier (\ref{def:mapOverSpecIso}). Hence the two composites
  \(\mathrm{match}_A \circ \sigma_{!}(\Spec_L f)\) and
  \(\Spec_k f \circ \mathrm{match}_B\), both morphisms
  \(\sigma_{!}(\Spec B/\Spec L) \to \Spec A/\Spec k\) of \(\Over(\Spec k)\), have equal
  underlying morphisms \(\Spec f\); morphisms of a slice category with equal underlying
  morphisms are equal, so the matchings intertwine \(\sigma_{!}(\Spec_L f)\) with
  \(\Spec_k f\).

  \emph{Pasting.} Factor \(c_A = w_A \circ b_A\) and \(c_B = w_B \circ b_B\) as in
  \ref{def:crossBaseAffineIso}. Whiskering with \(C\) and passing to underlying schemes
  is functorial, so the intertwining above whiskers to
  \((C \lhd \Spec_k f)_{\mathrm{left}} \circ w_B
    = w_A \circ \bigl(C \lhd \sigma_{!}(\Spec_L f)\bigr)_{\mathrm{left}}\).
  Naturality of the cross-base comparison in the test
  (\ref{thm:crossBaseIso_naturality}, at the frozen instantiation of
  \ref{thm:crossBase_baseChange}) applied to \(\Spec_L f\) gives
  \(\bigl(C \lhd \sigma_{!}(\Spec_L f)\bigr)_{\mathrm{left}} \circ b_B
    = b_A \circ (\sigma^{*}C \lhd \Spec_L f)_{\mathrm{left}}\).
  Chaining,
  \[
    (C \lhd \Spec_k f)_{\mathrm{left}} \circ c_B
      = w_A \circ \bigl(C \lhd \sigma_{!}(\Spec_L f)\bigr)_{\mathrm{left}} \circ b_B
      = w_A \circ b_A \circ (\sigma^{*}C \lhd \Spec_L f)_{\mathrm{left}}
      = c_A \circ (\sigma^{*}C \lhd \Spec_L f)_{\mathrm{left}} . \qedhere
  \]
\end{proof}

\begin{lemma}[The transition is flat, affine and surjective]
  \label{lem:baseFieldTransition_flat_affine_surjective}
  \dcref{custom}
  \lean{AlgebraicGeometry.flat_specMap_algHom, AlgebraicGeometry.surjective_specMap_algHom,
    AlgebraicGeometry.flat_baseFieldTransition,
    AlgebraicGeometry.isAffineHom_baseFieldTransition,
    AlgebraicGeometry.surjective_baseFieldTransition}
  \uses{thm:isPullback_baseFieldTransition}
  \leanok
  Let \(K_1\) be a field and \(K_2\) a nonzero commutative ring. Then
  \(\Spec\varphi : \Spec K_2 \to \Spec K_1\) is flat and surjective, and it is affine
  since both schemes are affine. Consequently, for field extensions \(K_1, K_2\) of
  \(k\), the base-field transition \(\pi : C_{K_2} \to C_{K_1}\) is flat, affine and
  surjective.
\end{lemma}
\begin{proof}
  \leanok
  Flatness of \(\Spec\varphi\) is flatness of \(K_2\) as a \(K_1\)-module, and every
  module over a field is flat (free). Surjectivity: \(\Spec K_1\) is a one-point space
  and \(\Spec K_2\) is nonempty (as \(K_2 \neq 0\)), so the unique point is hit. Each of
  the three properties --- flatness, affineness, surjectivity --- is stable under base
  change of schemes, and the transition square
  (\ref{thm:isPullback_baseFieldTransition}) exhibits \(\pi\) as a base change of
  \(\Spec\varphi\); so all three transport to \(\pi\).
\end{proof}

\begin{lemma}[Generic points under surjective morphisms]
  \label{lem:genericPoint_eq_of_surjective}
  \dcref{custom}
  \lean{AlgebraicGeometry.genericPoint_eq_of_surjective}
  \leanok
  A surjective morphism \(f : X \to Y\) of schemes with irreducible underlying spaces
  maps the generic point to the generic point: \(f(\eta_X) = \eta_Y\).
\end{lemma}
\begin{proof}
  \leanok
  Choose \(z \in X\) with \(f(z) = \eta_Y\). The generic point \(\eta_X\) specialises to
  every point of \(X\), in particular \(\eta_X \rightsquigarrow z\); applying the
  continuous map \(f\) gives \(f(\eta_X) \rightsquigarrow \eta_Y\). Conversely
  \(\eta_Y\) specialises to every point of \(Y\), so
  \(\eta_Y \rightsquigarrow f(\eta_X)\). The underlying space of a scheme is sober,
  hence \(T_0\), so the two specialisations force \(f(\eta_X) = \eta_Y\).
\end{proof}

\begin{lemma}[The transition matches generic points]
  \label{lem:baseFieldTransition_genericPoint}
  \dcref{custom}
  \lean{AlgebraicGeometry.baseFieldTransition_genericPoint}
  \uses{lem:genericPoint_eq_of_surjective, lem:baseFieldTransition_flat_affine_surjective,
    thm:baseChangeCurve_integral}
  \leanok
  Let \(C\) be smooth of relative dimension one and geometrically irreducible over the
  field \(k\), and let \(K_1, K_2\) be field extensions of \(k\). Then the base-field
  transition satisfies \(\pi(\eta_{C_{K_2}}) = \eta_{C_{K_1}}\).
\end{lemma}
\begin{proof}
  \leanok
  Both base-changed curves are integral (\ref{thm:baseChangeCurve_integral}), in
  particular irreducible, and \(\pi\) is surjective
  (\ref{lem:baseFieldTransition_flat_affine_surjective}); apply
  \ref{lem:genericPoint_eq_of_surjective}.
\end{proof}

\begin{definition}[Pullback of rational functions]
  \label{def:functionFieldMap}
  \dcref{custom}
  \lean{AlgebraicGeometry.Scheme.Hom.functionFieldMap,
    AlgebraicGeometry.Scheme.Hom.functionFieldMapUnits,
    AlgebraicGeometry.Scheme.Hom.coe_functionFieldMapUnits}
  \leanok
  Let \(f : X \to Y\) be a morphism of schemes with irreducible underlying spaces such
  that \(f(\eta_X) = \eta_Y\). The \emph{function-field map}
  \[
    f^{\sharp} : K(Y) \longrightarrow K(X)
  \]
  is the stalk map of \(f\) at \(\eta_X\), preceded by the canonical identification of
  the stalk of \(Y\) at \(f(\eta_X)\) with the stalk at \(\eta_Y\). It induces a group
  homomorphism \(K(Y)^{\times} \to K(X)^{\times}\) on units. For the base-field
  transition (with \ref{lem:baseFieldTransition_genericPoint} matching the generic
  points) this is the pullback of rational functions
  \(\pi^{\sharp} : K(C_{K_1}) \to K(C_{K_2})\).
\end{definition}

\begin{lemma}[Germ naturality and injectivity of the function-field map]
  \label{lem:functionFieldMap_germ}
  \dcref{custom}
  \lean{AlgebraicGeometry.Scheme.Hom.functionFieldMap_germ,
    AlgebraicGeometry.Scheme.Hom.functionFieldMap_injective}
  \uses{def:functionFieldMap}
  \leanok
  In the setting of \ref{def:functionFieldMap}: for an open \(U \subseteq Y\) containing
  \(\eta_Y\) and a section \(s \in \Gamma(Y, U)\),
  \[
    f^{\sharp}\bigl(\mathrm{germ}_{\eta_Y}(s)\bigr)
      \;=\; \mathrm{germ}_{\eta_X}\bigl(f^{\sharp}s\bigr)
      \qquad \text{in } K(X),
  \]
  where \(f^{\sharp}s \in \Gamma(X, f^{-1}U)\) is the pullback of the section (note
  \(\eta_X \in f^{-1}U\) since \(f(\eta_X) = \eta_Y \in U\)). If moreover \(X\) and
  \(Y\) are integral, the function-field map is injective.
\end{lemma}
\begin{proof}
  \leanok
  Germ naturality: the stalk identification carries \(\mathrm{germ}_{\eta_Y}(s)\) to the
  germ of \(s\) at the point \(f(\eta_X) \in U\), and the stalk map of \(f\) at a point
  sends the germ of a section to the germ of its pullback --- the defining compatibility
  of the stalk map with germs. Injectivity: on integral schemes both function fields are
  fields, and a ring homomorphism out of a field is injective.
\end{proof}

\begin{theorem}[Base change of sections along the transition, pushout form]
  \label{thm:isPushout_transitionSections}
  \dcref{custom}
  \lean{AlgebraicGeometry.Over.isPushout_transitionSections_of_isAffineOpen}
  \uses{thm:isPullback_baseFieldTransition}
  \leanok
  Let \(K_1, K_2\) be field extensions of \(k\), \(\varphi : K_1 \to K_2\) a
  \(k\)-algebra homomorphism, and \(V \subseteq C_{K_1}\) an \emph{affine} open. The
  square of section rings
  \[
    \begin{array}{ccc}
      \Gamma(\Spec K_1, \top) & \xrightarrow{\ \mathrm{pr}_2^{\sharp}\ } &
        \Gamma(C_{K_1}, V) \\
      \downarrow (\Spec\varphi)^{\sharp} & & \downarrow \pi^{\sharp} \\
      \Gamma(\Spec K_2, \top) & \xrightarrow{\ \mathrm{pr}_2^{\sharp}\ } &
        \Gamma\bigl(C_{K_2},\, \pi^{-1}V\bigr)
    \end{array}
  \]
  is a pushout of commutative rings. No flatness enters in the affine case.
\end{theorem}
\begin{proof}
  \leanok
  Restrict the transition square (\ref{thm:isPullback_baseFieldTransition}) over the
  affine opens: \(\pi^{-1}V\) is the fibre product \(V \times_{\Spec K_1} \Spec K_2\) of
  the open subscheme \(V\), a fibre product of affine schemes, hence affine with
  \(\Gamma(\pi^{-1}V) = \Gamma(V) \otimes_{\Gamma(\Spec K_1)} \Gamma(\Spec K_2)\):
  \(\Spec\) converts pushouts of commutative rings into fibre products of affine
  schemes. The four maps of the displayed square are the section-level maps of that
  restricted square, so the square is the tensor-product pushout.
\end{proof}

\begin{definition}[The second-projection algebra structure on sections]
  \label{def:baseChangeSectionsAlgebra}
  \dcref{custom}
  \lean{AlgebraicGeometry.baseChangeBundle, AlgebraicGeometry.baseChangeSectionsAlgebra,
    AlgebraicGeometry.ofHom_algebraMap_baseChangeSections}
  \uses{thm:baseChangeCurve_bundle}
  \leanok
  Let \(K/k\) be a field extension and \(V \subseteq C_K\) an open. The base-changed
  curve, bundled over \(\Spec K\) by its second projection
  (\ref{thm:baseChangeCurve_bundle}), endows \(\Gamma(C_K, V)\) with a \(K\)-algebra
  structure: the structure map is the pullback of global sections of \(\Spec K\) along
  \(\mathrm{pr}_2\), followed by restriction to \(V\),
  \[
    K \;=\; \Gamma(\Spec K, \top) \xrightarrow{\ \mathrm{pr}_2^{\sharp}\ }
      \Gamma(C_K, V).
  \]
  All \(K\)-linear statements about sections of \(C_K\) below refer to this structure.
\end{definition}

\begin{lemma}[Restriction of sections is \(K\)-algebraic]
  \label{lem:baseChangeResAlgHom}
  \dcref{custom}
  \lean{AlgebraicGeometry.ofHom_algebraMap_baseChangeSections_comp_res,
    AlgebraicGeometry.baseChangeResAlgHom, AlgebraicGeometry.baseChangeResAlgHom_apply}
  \uses{def:baseChangeSectionsAlgebra}
  \leanok
  For opens \(W \le V\) of \(C_K\), the structure maps of
  \ref{def:baseChangeSectionsAlgebra} intertwine restriction --- restricting the
  pullback of a constant is the pullback of the constant --- so restriction
  \(\Gamma(C_K, V) \to \Gamma(C_K, W)\) is a homomorphism of \(K\)-algebras.
\end{lemma}
\begin{proof}
  \leanok
  Both maps \(K \to \Gamma(C_K, W)\) --- structure map of \(W\), and structure map of
  \(V\) followed by restriction --- are the pullback of global sections of \(\Spec K\)
  along \(\mathrm{pr}_2\) restricted to \(W\), since pullback of sections commutes with
  restriction. A ring homomorphism commuting with the structure maps is a map of
  \(K\)-algebras.
\end{proof}

\begin{theorem}[The algebra-corner pushout square]
  \label{thm:isPushout_algebraMap_transitionSections}
  \dcref{custom}
  \lean{AlgebraicGeometry.Over.isPushout_algebraMap_transitionSections}
  \uses{thm:isPushout_transitionSections, def:baseChangeSectionsAlgebra}
  \leanok
  In the setting of \ref{thm:isPushout_transitionSections}, the square
  \[
    \begin{array}{ccc}
      K_1 & \longrightarrow & \Gamma(C_{K_1}, V) \\
      \downarrow \varphi & & \downarrow \pi^{\sharp} \\
      K_2 & \longrightarrow & \Gamma\bigl(C_{K_2},\, \pi^{-1}V\bigr)
    \end{array}
  \]
  --- top and bottom the structure maps of the second-projection algebra structures
  (\ref{def:baseChangeSectionsAlgebra}), left the homomorphism \(\varphi\) itself --- is
  a pushout of commutative rings. It is stated for an arbitrary \(k\)-algebra
  homomorphism \(\varphi\); no \(K_1\)-algebra structure on \(K_2\) is chosen.
\end{theorem}
\begin{proof}
  \leanok
  Identify the corners \(\Gamma(\Spec K_i, \top) = K_i\) of the pushout square of
  \ref{thm:isPushout_transitionSections} by the canonical isomorphisms. Under these, the
  left edge \((\Spec\varphi)^{\sharp}\) becomes \(\varphi\) (naturality of the
  identification in the ring), and the two horizontal edges become the structure maps of
  \ref{def:baseChangeSectionsAlgebra}. A square isomorphic to a pushout square is a
  pushout square.
\end{proof}

\begin{definition}[Sections base change along the base-field transition]
  \label{def:transitionSectionsBaseChange}
  \dcref{custom}
  \lean{AlgebraicGeometry.Over.transitionSectionsBaseChangeIso,
    AlgebraicGeometry.Over.transitionSectionsBaseChange}
  \uses{thm:isPushout_algebraMap_transitionSections, def:baseChangeSectionsAlgebra}
  \leanok
  Let \(K_1 \subseteq K_2\) be a tower of field extensions of \(k\) (that is, a
  \(K_1\)-algebra structure on \(K_2\) compatible with the \(k\)-structures), with
  \(\varphi\) the inclusion, and let \(V \subseteq C_{K_1}\) be an affine open. There is
  an isomorphism of rings
  \[
    \Gamma(C_{K_1}, V) \otimes_{K_1} K_2 \;\cong\;
      \Gamma\bigl(C_{K_2},\, \pi^{-1}V\bigr),
  \]
  the \(K_1\)-module structure on \(\Gamma(C_{K_1}, V)\) being the second-projection one
  (\ref{def:baseChangeSectionsAlgebra}).
\end{definition}
\begin{proof}
  \leanok
  The tensor product \(\Gamma(C_{K_1}, V) \otimes_{K_1} K_2\) with its two canonical
  maps is a pushout of the same span \(K_1 \to \Gamma(C_{K_1}, V)\), \(K_1 \to K_2\)
  --- the tensor product of commutative algebras is their pushout --- as the square of
  \ref{thm:isPushout_algebraMap_transitionSections}. Two pushouts of one span are
  canonically isomorphic under the corners.
\end{proof}

\begin{lemma}[The base change on pure tensors]
  \label{lem:transitionSectionsBaseChange_tmul}
  \dcref{custom}
  \lean{AlgebraicGeometry.Over.transitionSectionsBaseChange_tmul_one,
    AlgebraicGeometry.Over.transitionSectionsBaseChange_one_tmul,
    AlgebraicGeometry.Over.transitionSectionsBaseChange_tmul}
  \uses{def:transitionSectionsBaseChange}
  \leanok
  The isomorphism of \ref{def:transitionSectionsBaseChange} sends
  \[
    s \otimes 1 \longmapsto \pi^{\sharp}(s), \qquad
    1 \otimes a \longmapsto \mathrm{pr}_2^{\sharp}(a), \qquad
    s \otimes a \longmapsto \pi^{\sharp}(s) \cdot \mathrm{pr}_2^{\sharp}(a),
  \]
  for \(s \in \Gamma(C_{K_1}, V)\) and \(a \in K_2\).
\end{lemma}
\begin{proof}
  \leanok
  The comparison isomorphism of two pushouts of one span commutes with the two
  coprojections. On the tensor-product side the coprojections are \(s \mapsto s \otimes
  1\) and \(a \mapsto 1 \otimes a\); on the sections side they are \(\pi^{\sharp}\) and
  \(\mathrm{pr}_2^{\sharp}\) (\ref{thm:isPushout_algebraMap_transitionSections}). This
  gives the first two rules, and the third follows since \(s \otimes a = (s \otimes 1)
  \cdot (1 \otimes a)\) and the comparison is a ring homomorphism.
\end{proof}

\begin{lemma}[Naturality of the base change in the open]
  \label{lem:transitionSectionsBaseChange_naturality}
  \dcref{custom}
  \lean{AlgebraicGeometry.Over.transitionSectionsBaseChange_naturality}
  \uses{def:transitionSectionsBaseChange, lem:transitionSectionsBaseChange_tmul,
    lem:baseChangeResAlgHom}
  \leanok
  For affine opens \(W \le V\) of \(C_{K_1}\), restricting the base-changed section is
  base-changing the restriction: the square
  \[
    \begin{array}{ccc}
      \Gamma(C_{K_1}, V) \otimes_{K_1} K_2 & \cong &
        \Gamma\bigl(C_{K_2},\, \pi^{-1}V\bigr) \\
      \downarrow \mathrm{res} \otimes \mathrm{id} & & \downarrow \mathrm{res} \\
      \Gamma(C_{K_1}, W) \otimes_{K_1} K_2 & \cong &
        \Gamma\bigl(C_{K_2},\, \pi^{-1}W\bigr)
    \end{array}
  \]
  commutes, the left vertical being the (\(K_1\)-algebraic,
  \ref{lem:baseChangeResAlgHom}) restriction tensored with the identity of \(K_2\).
\end{lemma}
\begin{proof}
  \leanok
  Both composites are additive, so it suffices to compare them on a pure tensor
  \(s \otimes a\). By \ref{lem:transitionSectionsBaseChange_tmul} the top-right composite
  gives \(\bigl(\pi^{\sharp}(s) \cdot \mathrm{pr}_2^{\sharp}(a)\bigr)|_{\pi^{-1}W} =
  \pi^{\sharp}(s)|_{\pi^{-1}W} \cdot \mathrm{pr}_2^{\sharp}(a)|_{\pi^{-1}W}\)
  (restriction is a ring homomorphism), while the left-bottom composite gives
  \(\pi^{\sharp}(s|_W) \cdot \mathrm{pr}_2^{\sharp}(a)\) read over \(\pi^{-1}W\). These
  agree because pullback of sections commutes with restriction on both factors.
\end{proof}

\section{Two-lattice pairs and their cohomology}
\label{sec:two_lattice_pairs}

The remaining sections of this chapter build the module-level core of the two-term
pushforward engine. All schemes disappear: a \emph{two-lattice pair} axiomatises, in pure
commutative algebra, exactly the data that an affine two-cover extracts from a
quasi-coherent sheaf --- the two chart section modules, the overlap module, the two chart
coordinates, and the two restriction maps, each exhibiting the overlap as a localization
of its chart (the module counterpart of denominator clearing and annihilation on affine
opens, \ref{lem:denominator_clearing} and \ref{lem:section_annihilation}). The \v{C}ech
complex of the two-cover (\ref{def:TwoCover_diff}, \ref{def:TwoCover_H1Cok}) becomes a
two-term complex of modules; this section studies its cohomology in degrees zero and one
and its functoriality, the next section proves the six-term exact sequence of a short
exact sequence of pairs, and the following two sections compute the line-bundle models on
the two-chart projective line and deduce the coherence theorems: finiteness of \(H^1\)
over an \emph{arbitrary} commutative base ring, and of \(H^0\) over a Noetherian one. The
fibrewise Nakayama arguments and the split-rigidity theorems of
Chapter~\ref{chap:Algebra} (Sections~\ref{sec:nakayama_vanishing}
and~\ref{sec:split_rigidity}) consume the finiteness produced here.

Throughout, \(R\) is a commutative ring.

\begin{definition}[Two-lattice pair]
  \label{def:twoLatticePair}
  \dcref{custom}
  \lean{TwoLatticePair}
  \leanok
  A \emph{two-lattice pair} \(P\) over \(R\) consists of:
  \begin{itemize}
    \item \(R\)-modules \(M_0\), \(M_1\) (the \emph{chart lattices}) and \(N\) (the
      \emph{overlap module});
    \item \(R\)-linear endomorphisms \(t_0\) of \(M_0\) and \(t_1\) of \(M_1\) (the
      \emph{coordinate actions}) and an \emph{invertible} \(R\)-linear endomorphism
      \(u\) of \(N\);
    \item \(R\)-linear maps \(\iota_0 : M_0 \to N\) and \(\iota_1 : M_1 \to N\) (the
      \emph{restrictions to the overlap}) intertwining the actions:
      \(\iota_0(t_0 x) = u(\iota_0 x)\) for all \(x \in M_0\) and
      \(\iota_1(t_1 y) = u^{-1}(\iota_1 y)\) for all \(y \in M_1\);
  \end{itemize}
  subject to the localization axioms:
  \begin{enumerate}
    \item \emph{denominator clearing}: every \(n \in N\) admits an integer \(m \ge 0\)
      and \(x \in M_0\) with \(u^m n = \iota_0 x\), and an integer \(m' \ge 0\) and
      \(y \in M_1\) with \(u^{-m'} n = \iota_1 y\);
    \item \emph{annihilation}: if \(\iota_0 x = 0\) then \(t_0^m x = 0\) for some
      \(m \ge 0\), and if \(\iota_1 y = 0\) then \(t_1^m y = 0\) for some \(m \ge 0\).
  \end{enumerate}
  Regarding \(M_0\) and \(N\) as modules over the polynomial ring \(R[t]\), with \(t\)
  acting as \(t_0\) on \(M_0\) and as \(u\) on \(N\), the three conditions ---
  invertibility of \(u\) on \(N\), clearing, annihilation --- say precisely that
  \(\iota_0\) exhibits \(N\) as the localization of \(M_0\) at the powers of \(t\);
  likewise \(\iota_1\) exhibits \(N\) as the localization of \(M_1\) at the powers of
  \(t\), for the structures through \(t_1\) and \(u^{-1}\). In the geometric application
  \(M_0, M_1\) are the sections of a quasi-coherent sheaf on the two charts of an affine
  two-cover, \(N\) its sections on the overlap, and \(t_0, t_1\) the two chart
  coordinates, mutually inverse on the overlap.
\end{definition}

\begin{definition}[The two-term complex of a pair]
  \label{def:twoLattice_diff}
  \dcref{custom}
  \lean{TwoLatticePair.diff}
  \uses{def:twoLatticePair}
  \leanok
  The \emph{difference map} of a two-lattice pair \(P\) is the \(R\)-linear map
  \[
    \delta_P : M_0 \times M_1 \longrightarrow N, \qquad
    \delta_P(x, y) \;=\; \iota_0 x - \iota_1 y .
  \]
  The two-term complex \(M_0 \times M_1 \xrightarrow{\ \delta_P\ } N\) is the module
  avatar of the two-cover \v{C}ech complex.
\end{definition}

\begin{definition}[\(H^0\) of a pair]
  \label{def:twoLattice_H0}
  \dcref{custom}
  \lean{TwoLatticePair.H0}
  \uses{def:twoLattice_diff}
  \leanok
  The degree-zero cohomology of the pair \(P\) is the kernel of the difference map,
  \[
    H^0(P) \;=\; \ker \delta_P
    \;=\; \{\, (x, y) \in M_0 \times M_1 : \iota_0 x = \iota_1 y \,\},
  \]
  an \(R\)-submodule of \(M_0 \times M_1\): the pairs of chart sections agreeing on the
  overlap.
\end{definition}

\begin{definition}[\(H^1\) of a pair]
  \label{def:twoLattice_H1}
  \dcref{custom}
  \lean{TwoLatticePair.H1, TwoLatticePair.imageLattice, TwoLatticePair.h1Mk}
  \uses{def:twoLatticePair}
  \leanok
  The \emph{image lattice} of the pair \(P\) is the \(R\)-submodule
  \(Q(P) = \iota_0(M_0) + \iota_1(M_1) \subseteq N\), and the degree-one cohomology of
  \(P\) is the quotient
  \[
    H^1(P) \;=\; N \big/ Q(P),
  \]
  with quotient projection \(N \twoheadrightarrow H^1(P)\), \(n \mapsto [n]\).
\end{definition}

\begin{lemma}[\(H^1\) is the cokernel of the difference map]
  \label{lem:twoLattice_range_diff}
  \dcref{custom}
  \lean{TwoLatticePair.range_diff_eq}
  \uses{def:twoLattice_diff, def:twoLattice_H1}
  \leanok
  The image of \(\delta_P\) is the image lattice: \(\delta_P(M_0 \times M_1) = Q(P)\).
  In particular \(H^1(P)\) is the cokernel of \(\delta_P\), so \(H^0(P)\) and \(H^1(P)\)
  are the two cohomology modules of the two-term complex of \ref{def:twoLattice_diff}.
\end{lemma}
\begin{proof}
  \leanok
  Each value \(\delta_P(x,y) = \iota_0 x - \iota_1 y\) is a difference of an element of
  \(\iota_0(M_0)\) and an element of \(\iota_1(M_1)\), hence lies in \(Q(P)\).
  Conversely \(\iota_0 x = \delta_P(x, 0)\) and \(\iota_1 y = \delta_P(0, -y)\), so both
  summands of \(Q(P)\) lie in the image, and the image is a submodule.
\end{proof}

\begin{definition}[Map of two-lattice pairs]
  \label{def:twoLattice_hom}
  \dcref{custom}
  \lean{TwoLatticePair.Hom}
  \uses{def:twoLatticePair}
  \leanok
  A \emph{map of two-lattice pairs} \(f : P \to P'\) is a triple of \(R\)-linear maps
  \(f_0 : M_0 \to M_0'\), \(f_1 : M_1 \to M_1'\), \(f_N : N \to N'\) commuting with the
  coordinate actions and the restrictions:
  \[
    f_0 \circ t_0 = t_0' \circ f_0, \quad
    f_1 \circ t_1 = t_1' \circ f_1, \quad
    f_N \circ u = u' \circ f_N, \quad
    f_N \circ \iota_0 = \iota_0' \circ f_0, \quad
    f_N \circ \iota_1 = \iota_1' \circ f_1 .
  \]
  Compatibility with the \emph{inverse} overlap action is automatic: applying the
  \(u\)-compatibility to \(u^{-1} n\) gives \(f_N(n) = u'\bigl(f_N(u^{-1} n)\bigr)\),
  whence \(f_N(u^{-1} n) = u'^{-1}(f_N n)\).
\end{definition}

\begin{lemma}[Naturality of the difference map]
  \label{lem:twoLattice_diff_natural}
  \dcref{custom}
  \lean{TwoLatticePair.Hom.diff_comm}
  \uses{def:twoLattice_hom, def:twoLattice_diff}
  \leanok
  For a map of pairs \(f : P \to P'\) and \((x, y) \in M_0 \times M_1\),
  \[
    \delta_{P'}\bigl(f_0 x, f_1 y\bigr) \;=\; f_N\bigl(\delta_P(x, y)\bigr) .
  \]
\end{lemma}
\begin{proof}
  \leanok
  Expand both sides:
  \(\delta_{P'}(f_0 x, f_1 y) = \iota_0'(f_0 x) - \iota_1'(f_1 y)
  = f_N(\iota_0 x) - f_N(\iota_1 y) = f_N(\iota_0 x - \iota_1 y)\), using the two
  restriction compatibilities and linearity of \(f_N\).
\end{proof}

\begin{definition}[The induced map on \(H^0\)]
  \label{def:twoLattice_h0Map}
  \dcref{custom}
  \lean{TwoLatticePair.Hom.h0Map}
  \uses{def:twoLattice_hom, def:twoLattice_H0}
  \leanok
  A map of pairs \(f : P \to P'\) induces an \(R\)-linear map
  \(H^0(f) : H^0(P) \to H^0(P')\), \((x, y) \mapsto (f_0 x, f_1 y)\).
\end{definition}
\begin{proof}
  \leanok
  \uses{lem:twoLattice_diff_natural}
  If \(\delta_P(x,y) = 0\) then
  \(\delta_{P'}(f_0 x, f_1 y) = f_N(\delta_P(x,y)) = 0\) by naturality
  (\ref{lem:twoLattice_diff_natural}), so the restriction of \(f_0 \times f_1\) to
  \(H^0(P)\) lands in \(H^0(P')\); it is \(R\)-linear as a restriction of an \(R\)-linear
  map.
\end{proof}

\begin{definition}[The induced map on \(H^1\)]
  \label{def:twoLattice_h1Map}
  \dcref{custom}
  \lean{TwoLatticePair.Hom.h1Map}
  \uses{def:twoLattice_hom, def:twoLattice_H1}
  \leanok
  A map of pairs \(f : P \to P'\) induces an \(R\)-linear map
  \(H^1(f) : H^1(P) \to H^1(P')\), \([n] \mapsto [f_N n]\).
\end{definition}
\begin{proof}
  \leanok
  \(f_N\) carries the image lattice into the image lattice:
  \(f_N(\iota_0 x) = \iota_0'(f_0 x)\) and \(f_N(\iota_1 y) = \iota_1'(f_1 y)\), so
  \(f_N(Q(P)) \subseteq Q(P')\) and \(f_N\) descends to the quotients.
\end{proof}

\begin{lemma}[Right exactness of \(H^1\) in the pair]
  \label{lem:twoLattice_h1Map_surjective}
  \dcref{custom}
  \lean{TwoLatticePair.Hom.h1Map_surjective}
  \uses{def:twoLattice_h1Map}
  \leanok
  If a map of pairs \(f : P \to P'\) is surjective on overlap modules, then
  \(H^1(f) : H^1(P) \to H^1(P')\) is surjective.
\end{lemma}
\begin{proof}
  \leanok
  A class in \(H^1(P')\) is \([n']\) for some \(n' \in N'\); write \(n' = f_N(n)\) by
  surjectivity; then \(H^1(f)[n] = [f_N n] = [n']\).
\end{proof}

\section{The six-term exact sequence of two-lattice pairs}
\label{sec:two_lattice_sixterm}

A short exact sequence of two-lattice pairs induces the entire long exact cohomology
sequence of the two-term complexes; since a two-term complex has no cohomology beyond
degree one, the sequence has six terms and closes on both ends:
\[
  0 \to H^0(P) \to H^0(P') \to H^0(P'')
    \xrightarrow{\ \partial\ } H^1(P) \to H^1(P') \to H^1(P'') \to 0 .
\]
The connecting map and the six exactness statements are proved by explicit element
chases on the two-term complexes.

\begin{definition}[Short exact sequence of pairs]
  \label{def:twoLattice_shortExact}
  \dcref{custom}
  \lean{TwoLatticePair.IsShortExact}
  \uses{def:twoLattice_hom}
  \leanok
  A pair of maps of two-lattice pairs \(\varphi : P \to P'\) and \(\psi : P' \to P''\)
  is a \emph{short exact sequence} \(0 \to P \to P' \to P'' \to 0\) if it is
  componentwise short exact: on each of the three carriers (the two chart lattices and
  the overlap module), \(\varphi\) is injective, \(\psi\) is surjective, and the image
  of \(\varphi\) is the kernel of \(\psi\).
\end{definition}

\begin{definition}[The connecting homomorphism]
  \label{def:twoLattice_delta}
  \dcref{custom}
  \lean{TwoLatticePair.IsShortExact.delta, TwoLatticePair.IsShortExact.delta_apply}
  \uses{def:twoLattice_shortExact, def:twoLattice_H0, def:twoLattice_H1}
  \leanok
  For a short exact sequence \(0 \to P \xrightarrow{\varphi} P' \xrightarrow{\psi}
  P'' \to 0\) of two-lattice pairs, the \emph{connecting homomorphism}
  \[
    \partial : H^0(P'') \longrightarrow H^1(P)
  \]
  is the \(R\)-linear map characterized by: whenever \(z = (z_0, z_1) \in H^0(P'')\),
  \(e_0 \in M_0'\) and \(e_1 \in M_1'\) satisfy \(\psi_0 e_0 = z_0\) and
  \(\psi_1 e_1 = z_1\), and \(n \in N\) satisfies
  \(\varphi_N n = \delta_{P'}(e_0, e_1)\), then \(\partial z = [n]\).
\end{definition}
\begin{proof}
  \leanok
  \uses{lem:twoLattice_diff_natural, lem:twoLattice_range_diff}
  \emph{Existence of the chase data.} Given \(z \in H^0(P'')\), componentwise
  surjectivity of \(\psi\) provides lifts \(e_0, e_1\). Then
  \(\psi_N\bigl(\delta_{P'}(e_0,e_1)\bigr) = \delta_{P''}(\psi_0 e_0, \psi_1 e_1)
  = \delta_{P''}(z) = 0\) by naturality (\ref{lem:twoLattice_diff_natural}) and
  \(z \in H^0(P'')\), so \(\delta_{P'}(e_0,e_1)\) lies in \(\ker \psi_N =
  \varphi_N(N)\), and \(n\) exists; it is unique because \(\varphi_N\) is injective.

  \emph{Independence of the lifts.} If \(e_0', e_1'\) is a second componentwise lift of
  \(z\), with \(\varphi_N n' = \delta_{P'}(e_0', e_1')\), then
  \(\psi_0(e_0' - e_0) = 0\) and \(\psi_1(e_1' - e_1) = 0\), so by componentwise
  exactness \(e_0' - e_0 = \varphi_0 a\) and \(e_1' - e_1 = \varphi_1 b\) for some
  \(a \in M_0\), \(b \in M_1\). Then
  \[
    \varphi_N(n' - n) \;=\; \delta_{P'}(\varphi_0 a, \varphi_1 b)
    \;=\; \varphi_N\bigl(\delta_P(a, b)\bigr)
  \]
  by naturality, and injectivity of \(\varphi_N\) gives
  \(n' - n = \delta_P(a,b) \in Q(P)\) (\ref{lem:twoLattice_range_diff}); hence
  \([n'] = [n]\) in \(H^1(P)\).

  \emph{Linearity.} The chase data may be chosen additively --- lifts of a sum as the
  sum of lifts, and similarly for scalar multiples --- and \(n\) depends linearly on
  the chosen data; since the class \([n]\) is independent of the choices, \(\partial\)
  is \(R\)-linear.
\end{proof}

\begin{theorem}[Six-term exactness, left end]
  \label{thm:twoLattice_h0Map_injective}
  \dcref{custom}
  \lean{TwoLatticePair.IsShortExact.h0Map_injective}
  \uses{def:twoLattice_shortExact, def:twoLattice_h0Map}
  \leanok
  For a short exact sequence \(0 \to P \xrightarrow{\varphi} P' \xrightarrow{\psi}
  P'' \to 0\), the map \(H^0(\varphi) : H^0(P) \to H^0(P')\) is injective.
\end{theorem}
\begin{proof}
  \leanok
  \(H^0(\varphi)\) acts componentwise as \((\varphi_0, \varphi_1)\), and both
  \(\varphi_0\) and \(\varphi_1\) are injective.
\end{proof}

\begin{theorem}[Six-term exactness at \(H^0(P')\)]
  \label{thm:twoLattice_exact_h0_middle}
  \dcref{custom}
  \lean{TwoLatticePair.IsShortExact.range_h0Map_eq_ker_h0Map}
  \uses{def:twoLattice_shortExact, def:twoLattice_h0Map}
  \leanok
  In the situation of \ref{thm:twoLattice_h0Map_injective}, the image of
  \(H^0(\varphi)\) equals the kernel of \(H^0(\psi)\).
\end{theorem}
\begin{proof}
  \leanok
  \uses{lem:twoLattice_diff_natural}
  \(\subseteq\): \(\psi \circ \varphi\) vanishes componentwise (image of \(\varphi\)
  inside kernel of \(\psi\)), so \(H^0(\psi) \circ H^0(\varphi) = 0\).
  \(\supseteq\): let \((z_0, z_1) \in H^0(P')\) with \(\psi_0 z_0 = 0\) and
  \(\psi_1 z_1 = 0\). By componentwise exactness \(z_0 = \varphi_0 a\) and
  \(z_1 = \varphi_1 b\). The pair \((a, b)\) lies in \(H^0(P)\): by naturality
  (\ref{lem:twoLattice_diff_natural})
  \(\varphi_N\bigl(\delta_P(a,b)\bigr) = \delta_{P'}(z_0, z_1) = 0\), and
  \(\varphi_N\) is injective. Then \(H^0(\varphi)(a, b) = (z_0, z_1)\).
\end{proof}

\begin{theorem}[Six-term exactness at \(H^0(P'')\)]
  \label{thm:twoLattice_exact_h0_delta}
  \dcref{custom}
  \lean{TwoLatticePair.IsShortExact.range_h0Map_eq_ker_delta}
  \uses{def:twoLattice_delta, def:twoLattice_h0Map}
  \leanok
  In the situation of \ref{thm:twoLattice_h0Map_injective}, the image of
  \(H^0(\psi)\) equals the kernel of the connecting homomorphism \(\partial\).
\end{theorem}
\begin{proof}
  \leanok
  \uses{lem:twoLattice_range_diff, lem:twoLattice_diff_natural}
  \(\subseteq\): for \(z = H^0(\psi)(w)\) with \(w \in H^0(P')\), take the components
  of \(w\) themselves as the lifts in the characterization of \(\partial\)
  (\ref{def:twoLattice_delta}); their difference is \(\delta_{P'}(w) = 0 =
  \varphi_N(0)\), so \(\partial z = [0] = 0\).
  \(\supseteq\): let \(z \in H^0(P'')\) with \(\partial z = 0\); choose lifts
  \(e_0, e_1\) and \(n\) with \(\varphi_N n = \delta_{P'}(e_0, e_1)\), so that
  \(\partial z = [n]\). Vanishing of \([n]\) means \(n \in Q(P) =
  \delta_P(M_0 \times M_1)\) (\ref{lem:twoLattice_range_diff}); write
  \(n = \delta_P(a, b)\). Then \(w := (e_0 - \varphi_0 a,\, e_1 - \varphi_1 b)\)
  satisfies
  \(\delta_{P'}(w) = \delta_{P'}(e_0,e_1) - \varphi_N(\delta_P(a,b)) = 0\)
  by naturality, so \(w \in H^0(P')\); and since \(\psi\circ\varphi = 0\)
  componentwise, \(H^0(\psi)(w) = (\psi_0 e_0, \psi_1 e_1) = z\).
\end{proof}

\begin{theorem}[Six-term exactness at \(H^1(P)\)]
  \label{thm:twoLattice_exact_delta_h1}
  \dcref{custom}
  \lean{TwoLatticePair.IsShortExact.range_delta_eq_ker_h1Map}
  \uses{def:twoLattice_delta, def:twoLattice_h1Map}
  \leanok
  In the situation of \ref{thm:twoLattice_h0Map_injective}, the image of the
  connecting homomorphism \(\partial\) equals the kernel of \(H^1(\varphi)\).
\end{theorem}
\begin{proof}
  \leanok
  \uses{lem:twoLattice_range_diff, lem:twoLattice_diff_natural}
  \(\subseteq\): if \(\partial z = [n]\) with \(\varphi_N n = \delta_{P'}(e_0, e_1)\),
  then \(H^1(\varphi)[n] = [\varphi_N n] = [\delta_{P'}(e_0,e_1)] = 0\), the class of
  an element of the image lattice (\ref{lem:twoLattice_range_diff}).
  \(\supseteq\): let \([n] \in H^1(P)\) with \(H^1(\varphi)[n] = 0\), i.e.\
  \(\varphi_N n \in Q(P') = \delta_{P'}(M_0' \times M_1')\); write
  \(\varphi_N n = \delta_{P'}(e_0, e_1)\). Then
  \(z := (\psi_0 e_0, \psi_1 e_1)\) lies in \(H^0(P'')\):
  \(\delta_{P''}(z) = \psi_N\bigl(\delta_{P'}(e_0,e_1)\bigr) = \psi_N(\varphi_N n) = 0\)
  by naturality and \(\psi\circ\varphi = 0\). By the characterization of \(\partial\)
  (\ref{def:twoLattice_delta}), \(\partial z = [n]\).
\end{proof}

\begin{theorem}[Six-term exactness at \(H^1(P')\)]
  \label{thm:twoLattice_exact_h1_middle}
  \dcref{custom}
  \lean{TwoLatticePair.IsShortExact.range_h1Map_eq_ker_h1Map}
  \uses{def:twoLattice_shortExact, def:twoLattice_h1Map}
  \leanok
  In the situation of \ref{thm:twoLattice_h0Map_injective}, the image of
  \(H^1(\varphi)\) equals the kernel of \(H^1(\psi)\).
\end{theorem}
\begin{proof}
  \leanok
  \uses{lem:twoLattice_range_diff, lem:twoLattice_diff_natural}
  \(\subseteq\): \(H^1(\psi)\bigl(H^1(\varphi)[n]\bigr) = [\psi_N(\varphi_N n)] = 0\).
  \(\supseteq\): let \([n'] \in H^1(P')\) with \([\psi_N n'] = 0\), i.e.\
  \(\psi_N n' \in Q(P'') = \delta_{P''}(M_0'' \times M_1'')\); write
  \(\psi_N n' = \delta_{P''}(z_0, z_1)\) and lift \(z_0 = \psi_0 e_0\),
  \(z_1 = \psi_1 e_1\). Then
  \(\psi_N\bigl(n' - \delta_{P'}(e_0, e_1)\bigr) =
  \psi_N n' - \delta_{P''}(z_0, z_1) = 0\) by naturality, so by exactness on the
  overlap \(n' - \delta_{P'}(e_0,e_1) = \varphi_N w\) for some \(w \in N\). Since
  \([\delta_{P'}(e_0,e_1)] = 0\) in \(H^1(P')\) (\ref{lem:twoLattice_range_diff}),
  \([n'] = [\varphi_N w] = H^1(\varphi)[w]\).
\end{proof}

\begin{theorem}[Six-term exactness, right end]
  \label{thm:twoLattice_h1Map_surjective_ses}
  \dcref{custom}
  \lean{TwoLatticePair.IsShortExact.h1Map_surjective}
  \uses{def:twoLattice_shortExact, def:twoLattice_h1Map}
  \leanok
  In the situation of \ref{thm:twoLattice_h0Map_injective}, the map
  \(H^1(\psi) : H^1(P') \to H^1(P'')\) is surjective.
\end{theorem}
\begin{proof}
  \leanok
  \uses{lem:twoLattice_h1Map_surjective}
  \(\psi\) is surjective on overlap modules; apply the right exactness of \(H^1\) in
  the pair (\ref{lem:twoLattice_h1Map_surjective}).
\end{proof}

\section{The line-bundle models and the finite surjection}
\label{sec:two_lattice_models}

The coherence theorems are proved by comparing an arbitrary pair with \emph{models}: the
two-lattice pairs of the direct sums \(\bigoplus_{i} \mathcal O(m_i)\) of line bundles on
the two-chart projective line over \(R\), whose cohomology is the classical monomial
window. Say a two-lattice pair \(P\) has \emph{finite lattices} if \(M_0\) is a finitely
generated module over the polynomial ring \(R[t]\), with \(t\) acting as \(t_0\), and
likewise \(M_1\) over \(R[t]\) through \(t_1\). Every pair with finite lattices receives
a componentwise surjection from a model --- the two-cover form of the finite free
approximation used by Mumford in the coherence dévissage on abelian varieties, here
terminating in one step because generation is measured over \(R[t]\) rather than over
\(R\).

\begin{definition}[The family line-bundle model]
  \label{def:twoLattice_model}
  \dcref{custom}
  \lean{TwoLatticePair.model}
  \uses{def:twoLatticePair}
  \leanok
  Let \(I\) be a finite index set and \(m : I \to \Z\) a family of twists. The
  \emph{model} \(E(m)\) is the two-lattice pair over \(R\) with chart lattices
  \(M_0 = M_1 = (R[t])^I\), overlap module \(N = (R[t, t^{-1}])^I\), coordinate actions
  \(t_0 = t_1 =\) componentwise multiplication by \(t\), invertible overlap action
  \(u =\) componentwise multiplication by \(t\), and restrictions
  \[
    \iota_0(p)_i \;=\; p_i(t), \qquad
    \iota_1(q)_i \;=\; t^{m_i} \cdot q_i\bigl(t^{-1}\bigr) ,
  \]
  the componentwise inclusion of polynomials into Laurent polynomials on the chart-0
  side, and the twisted substitution \(t \mapsto t^{-1}\) on the chart-1 side. It is
  the pair of the sheaf \(\bigoplus_{i \in I} \mathcal O(m_i)\) on the two standard
  charts of the projective line over \(R\).
\end{definition}
\begin{proof}
  \leanok
  The axioms of a two-lattice pair hold. \emph{Equivariance:} on the chart-0 side,
  multiplying a polynomial by \(t\) multiplies its Laurent image by \(t\); on the
  chart-1 side, \(t^{m_i} \cdot (t q)(t^{-1}) = t^{m_i} \cdot t^{-1} \cdot q(t^{-1})\),
  so \(\iota_1(t_1 q) = u^{-1}(\iota_1 q)\). \emph{Denominator clearing:} a Laurent
  vector \(n\) has finitely many components, each with exponents bounded below, so a
  single power \(u^M\) (componentwise multiplication by \(t^M\)) makes every component a
  polynomial and lands \(n\) in \(\iota_0\bigl((R[t])^I\bigr)\). Symmetrically, the
  exponents of each \(n_i\) are bounded above, say by \(D_i\); for any
  \(M' \ge \max_i (D_i - m_i)\) with \(M' \ge 0\), every exponent of
  \(t^{-M' - m_i}\, n_i\) is nonpositive, so there is a polynomial \(q_i\) with
  \(q_i(t^{-1}) = t^{-M' - m_i}\, n_i\), and then
  \(\bigl(u^{-M'} n\bigr)_i = t^{-M'} n_i = t^{m_i}\, q_i(t^{-1})\), i.e.\
  \(u^{-M'} n = \iota_1(q)\). \emph{Annihilation:} both
  \(\iota_0\) and \(\iota_1\) are injective --- the inclusion of polynomials is
  injective, the substitution \(t \mapsto t^{-1}\) is an automorphism of the Laurent
  ring, and multiplication by the unit \(t^{m_i}\) is injective --- so a section
  vanishing on the overlap is already zero.
\end{proof}

\begin{theorem}[The monomial window: \(H^1\) of a model over any ring]
  \label{thm:moduleFinite_h1_model}
  \lean{TwoLatticePair.moduleFinite_h1_model}
  \uses{def:twoLattice_model, def:twoLattice_H1}
  \leanok
  \dcref{custom}
  Let \(R\) be \emph{any} commutative ring, \(I\) finite and \(m : I \to \Z\). Then
  \(H^1\bigl(E(m)\bigr)\) is a finite \(R\)-module. (It is spanned by the classes of
  the monomials \(t^j\) in the components \(i\) with \(m_i < j < 0\) --- the classical
  negative-monomial window of the projective line.)
\end{theorem}
\begin{proof}
  \leanok
  \uses{cor:two_lattice_loc_pow}
  Apply the two-lattice finiteness ladder (\ref{cor:two_lattice_loc_pow}) with
  \(k = R\), the Laurent module \(N = (R[t,t^{-1}])^I\) --- finitely generated over
  \(R[t,t^{-1}]\), being free on \(I\) --- and the \(R\)-submodules
  \(Q_0 = \iota_0\bigl((R[t])^I\bigr)\), \(Q_1 = \iota_1\bigl((R[t])^I\bigr)\). \(Q_0\)
  is stable under multiplication by \(t\) and \(Q_1\) under multiplication by
  \(t^{-1}\), by the equivariance of \(\iota_0\) and \(\iota_1\)
  (\ref{def:twoLattice_model}); and every element of \(N\) is absorbed into \(Q_0\) by
  a large power of \(t\) and into \(Q_1\) by a large power of \(t^{-1}\), by the
  denominator-clearing axioms verified in \ref{def:twoLattice_model}. The ladder
  concludes that \(N/(Q_0 + Q_1) = H^1\bigl(E(m)\bigr)\) is a finite \(R\)-module.
\end{proof}

\begin{theorem}[The degree window: \(H^0\) of a model over a Noetherian ring]
  \label{thm:moduleFinite_h0_model}
  \lean{TwoLatticePair.moduleFinite_h0_model}
  \uses{def:twoLattice_model, def:twoLattice_H0}
  \leanok
  \dcref{custom}
  Let \(R\) be a \emph{Noetherian} commutative ring, \(I\) finite and \(m : I \to \Z\).
  Then \(H^0\bigl(E(m)\bigr)\) is a finite \(R\)-module.
\end{theorem}
\begin{proof}
  \leanok
  A class in \(H^0(E(m))\) is a pair \((p, q)\) of polynomial vectors with
  \(p_i(t) = t^{m_i} q_i(t^{-1})\) in \(R[t,t^{-1}]\) for every \(i\). Comparing
  coefficients: the right side has no exponent above \(m_i\) (the exponents of
  \(q_i(t^{-1})\) are \(\le 0\)), so \(\deg p_i \le m_i\) --- in particular \(p_i = 0\)
  when \(m_i < 0\). Moreover the first component determines the class: \(\iota_1\) is
  injective (\ref{def:twoLattice_model}) and \(\iota_1 q = \iota_0 p\). Hence
  \(H^0(E(m))\) embeds \(R\)-linearly into
  \(\bigoplus_{i} \{\, p \in R[t] : \deg p \le m_i \,\}\), a finitely generated free
  \(R\)-module (of rank \(\sum_i \max(m_i + 1, 0)\)). A submodule of a finitely
  generated module over the Noetherian ring \(R\) is finitely generated.
\end{proof}

\begin{definition}[The Laurent action on the overlap module]
  \label{def:twoLattice_laurentToEnd}
  \dcref{custom}
  \lean{TwoLatticePair.laurentToEnd}
  \uses{def:twoLatticePair}
  \leanok
  For a two-lattice pair \(P\), the \(R\)-algebra homomorphism
  \[
    R[t, t^{-1}] \longrightarrow \operatorname{End}_R(N), \qquad t \longmapsto u ,
  \]
  extending the assignment \(t^j \mapsto u^j\) (\(j \in \Z\), negative powers through
  the inverse of \(u\)) by \(R\)-linearity. It exists because \(u\) is invertible, so
  \(j \mapsto u^j\) is a homomorphism from \(\Z\) to the units of
  \(\operatorname{End}_R(N)\), and the
  Laurent ring is the monoid algebra of \(\Z\) over \(R\). We write \(g(u)\) for the
  action of \(g \in R[t,t^{-1}]\) on \(N\).
\end{definition}

\begin{lemma}[The Laurent action intertwines the chart actions]
  \label{lem:twoLattice_laurent_intertwine}
  \dcref{custom}
  \lean{TwoLatticePair.laurentToEnd_toLaurent_ι₀,
    TwoLatticePair.laurentToEnd_invert_toLaurent_ι₁}
  \uses{def:twoLattice_laurentToEnd}
  \leanok
  For a polynomial \(g \in R[t]\), an element \(x \in M_0\) and an element
  \(y \in M_1\),
  \[
    g(u)\bigl(\iota_0 x\bigr) = \iota_0\bigl(g(t_0)\, x\bigr)
    \qquad\text{and}\qquad
    g\bigl(u^{-1}\bigr)\bigl(\iota_1 y\bigr) = \iota_1\bigl(g(t_1)\, y\bigr) ,
  \]
  where \(g(t_0)\) and \(g(t_1)\) denote the evaluations of \(g\) at the coordinate
  actions.
\end{lemma}
\begin{proof}
  \leanok
  Both sides are additive in \(g\), so it suffices to treat a monomial
  \(g = c\,t^k\). Iterating the equivariance \(\iota_0(t_0 x) = u(\iota_0 x)\) gives
  \(\iota_0(t_0^k x) = u^k(\iota_0 x)\), and multiplying by \(c\) gives the chart-0
  identity; iterating \(\iota_1(t_1 y) = u^{-1}(\iota_1 y)\) gives the chart-1
  identity with \(u^{-k}\).
\end{proof}

\begin{definition}[The model map determined by matched generators]
  \label{def:twoLattice_modelHom}
  \dcref{custom}
  \lean{TwoLatticePair.modelHom}
  \uses{def:twoLattice_model, def:twoLattice_hom, def:twoLattice_laurentToEnd}
  \leanok
  Let \(P\) be a two-lattice pair, \(I\) a finite index set, \(m : I \to \Z\), and let
  \(a : I \to M_0\) and \(b : I \to M_1\) be families whose overlap images match up to
  the twist:
  \[
    u^{m_i}\bigl(\iota_0 a_i\bigr) \;=\; \iota_1 b_i \qquad (i \in I).
  \]
  The \emph{model map} \(f : E(m) \to P\) is the map of pairs with components
  \[
    f_0(p) = \sum_{i \in I} p_i(t_0)\, a_i, \qquad
    f_1(q) = \sum_{i \in I} q_i(t_1)\, b_i, \qquad
    f_N(n) = \sum_{i \in I} n_i(u)\, \bigl(\iota_0 a_i\bigr) ,
  \]
  the polynomial actions on the generators on the charts and the Laurent action
  (\ref{def:twoLattice_laurentToEnd}) on their overlap images.
\end{definition}
\begin{proof}
  \leanok
  \uses{lem:twoLattice_laurent_intertwine}
  The five compatibilities of a map of pairs (\ref{def:twoLattice_hom}) hold term by
  term. Compatibility with the coordinate actions: multiplying \(p_i\) (resp.\ \(q_i\),
  \(n_i\)) by \(t\) multiplies \(p_i(t_0)\) (resp.\ \(q_i(t_1)\), \(n_i(u)\)) by
  \(t_0\) (resp.\ \(t_1\), \(u\)), which is exactly the model action carried to \(P\).
  Compatibility with \(\iota_0\): for a polynomial vector \(p\),
  \(f_N(\iota_0 p) = \sum_i p_i(u)(\iota_0 a_i) = \sum_i \iota_0\bigl(p_i(t_0)a_i\bigr)
  = \iota_0(f_0 p)\) by the chart-0 intertwining
  (\ref{lem:twoLattice_laurent_intertwine}). Compatibility with \(\iota_1\): for a
  polynomial vector \(q\), the \(i\)-th term of \(f_N(\iota_1 q)\) is
  \(\bigl(t^{m_i} q_i(t^{-1})\bigr)(u)\,(\iota_0 a_i) =
  q_i(u^{-1})\Bigl(u^{m_i}(\iota_0 a_i)\Bigr) = q_i(u^{-1})(\iota_1 b_i)
  = \iota_1\bigl(q_i(t_1) b_i\bigr)\), using the matching hypothesis and the chart-1
  intertwining; summing gives \(f_N(\iota_1 q) = \iota_1(f_1 q)\).
\end{proof}

\begin{lemma}[Models have finite lattices]
  \label{lem:moduleFinite_aeval_model}
  \dcref{custom}
  \lean{TwoLatticePair.moduleFinite_aeval_model_t₀, TwoLatticePair.moduleFinite_aeval_model_t₁}
  \uses{def:twoLattice_model}
  \leanok
  The chart lattices of a model \(E(m)\) are finite: \((R[t])^I\), with \(t\) acting by
  componentwise multiplication, is a finitely generated \(R[t]\)-module.
\end{lemma}
\begin{proof}
  \leanok
  The \(R[t]\)-module structure through the coordinate action coincides with the
  standard one on the free module \((R[t])^I\), which is generated by the \(|I|\)
  standard basis vectors.
\end{proof}

\begin{theorem}[The finite model surjection]
  \label{thm:exists_hom_model_surjective}
  \dcref{custom}
  \lean{TwoLatticePair.exists_hom_model_surjective}
  \uses{def:twoLattice_modelHom}
  \leanok
  Let \(P\) be a two-lattice pair with finite lattices over a commutative ring \(R\).
  Then there exist a finite index set \(I\), twists \(m : I \to \Z\), and a map of
  pairs \(f : E(m) \to P\) surjective on the chart-0 lattice, on the chart-1 lattice,
  \emph{and} on the overlap module.
\end{theorem}
\begin{proof}
  \leanok
  \uses{def:twoLatticePair, lem:twoLattice_laurent_intertwine}
  Choose finite generating families \((a_j)_{j \in J_0}\) of \(M_0\) over \(R[t]\)
  (acting through \(t_0\)) and \((b_j)_{j \in J_1}\) of \(M_1\) over \(R[t]\) (through
  \(t_1\)). Match each generator across the overlap by denominator clearing
  (\ref{def:twoLatticePair}): for \(j \in J_0\) there are \(k_j \ge 0\) and
  \(y_j \in M_1\) with \(u^{-k_j}(\iota_0 a_j) = \iota_1 y_j\); set \(m_j = -k_j\), so
  \(u^{m_j}(\iota_0 a_j) = \iota_1 y_j\). For \(j \in J_1\) there are \(k_j \ge 0\) and
  \(x_j \in M_0\) with \(u^{k_j}(\iota_1 b_j) = \iota_0 x_j\); set \(m_j = -k_j\), so
  \(u^{m_j}(\iota_0 x_j) = u^{-k_j} u^{k_j} (\iota_1 b_j) = \iota_1 b_j\). Take
  \(I = J_0 \sqcup J_1\), with generator pairs \((a_j, y_j)\) for \(j \in J_0\) and
  \((x_j, b_j)\) for \(j \in J_1\), and let \(f : E(m) \to P\) be the model map of
  \ref{def:twoLattice_modelHom}.

  \emph{Chart surjectivity.} Any \(x \in M_0\) is an \(R[t]\)-combination
  \(x = \sum_{j \in J_0} g_j(t_0)\, a_j\); the vector with components \(g_j\) on
  \(J_0\) and \(0\) on \(J_1\) maps to \(x\) under \(f_0\). Symmetrically for \(M_1\).

  \emph{Overlap surjectivity.} Given \(n \in N\), clear denominators:
  \(u^k n = \iota_0 x\) for some \(k \ge 0\), \(x \in M_0\); pick \(p\) with
  \(f_0(p) = x\). Since \(f\) commutes with the overlap actions and their inverses and
  with \(\iota_0\),
  \[
    f_N\Bigl(u_E^{-k}\bigl(\iota_0^{E} p\bigr)\Bigr)
    \;=\; u^{-k}\bigl(\iota_0(f_0 p)\bigr)
    \;=\; u^{-k}\bigl(u^{k} n\bigr) \;=\; n ,
  \]
  where \(u_E\) and \(\iota_0^E\) denote the overlap action and restriction of the
  model.
\end{proof}

\section{The coherence theorems for two-lattice pairs}
\label{sec:two_lattice_coherence}

The two finiteness keystones of the engine: for a pair with finite lattices, \(H^1\) is
a finite \(R\)-module over \emph{every} commutative ring \(R\), and \(H^0\) is a finite
\(R\)-module over every \emph{Noetherian} \(R\). Noetherianity enters only through
\(H^0\); every downstream consequence that only needs \(H^1\) --- vanishing propagation
and the openness of the vanishing locus (Section~\ref{sec:nakayama_vanishing}) --- holds
relative to an arbitrary base ring.

\begin{definition}[The kernel pair of a map of pairs]
  \label{def:twoLattice_kernelPair}
  \dcref{custom}
  \lean{TwoLatticePair.Hom.kernelPair, TwoLatticePair.Hom.kernelPairIncl}
  \uses{def:twoLattice_hom}
  \leanok
  For a map of two-lattice pairs \(f : P \to P'\), the \emph{kernel pair} \(K(f)\) is
  the two-lattice pair with carriers \(\ker f_0 \subseteq M_0\),
  \(\ker f_1 \subseteq M_1\), \(\ker f_N \subseteq N\), the restricted coordinate and
  overlap actions, and the restricted maps \(\iota_0, \iota_1\); the componentwise
  inclusions form a map of pairs \(K(f) \to P\).
\end{definition}
\begin{proof}
  \leanok
  \uses{def:twoLatticePair}
  The kernels are stable under the actions and the restrictions because \(f\) commutes
  with them (for the inverse overlap action, by the automatic compatibility of
  \ref{def:twoLattice_hom}), so all the structure restricts; the annihilation axioms
  restrict verbatim. Denominator clearing must be re-proved: let
  \(n \in \ker f_N\). Clearing in \(P\) gives \(k \ge 0\) and \(x \in M_0\) with
  \(u^k n = \iota_0 x\). Then
  \(\iota_0'\bigl(f_0 x\bigr) = f_N(u^k n) = u'^k\bigl(f_N n\bigr) = 0\), so by the
  annihilation axiom of \(P'\) there is \(j \ge 0\) with \(t_0'^j (f_0 x) = 0\), i.e.\
  \(f_0\bigl(t_0^j x\bigr) = 0\): the element \(t_0^j x\) lies in \(\ker f_0\), and
  \[
    u^{j + k}\, n \;=\; u^j\bigl(\iota_0 x\bigr) \;=\; \iota_0\bigl(t_0^j x\bigr)
  \]
  exhibits clearing inside the kernel pair. The chart-1 side is symmetric. The
  componentwise inclusions commute with all the structure by construction.
\end{proof}

\begin{lemma}[The kernel pair sits in a short exact sequence]
  \label{lem:twoLattice_kernelPair_shortExact}
  \dcref{custom}
  \lean{TwoLatticePair.Hom.isShortExact_kernelPairIncl}
  \uses{def:twoLattice_kernelPair, def:twoLattice_shortExact}
  \leanok
  If a map of pairs \(f : P \to P'\) is surjective on all three carriers, then
  \[
    0 \longrightarrow K(f) \longrightarrow P \xrightarrow{\ f\ } P' \longrightarrow 0
  \]
  is a short exact sequence of two-lattice pairs.
\end{lemma}
\begin{proof}
  \leanok
  On each carrier, the inclusion of the kernel is injective, its image is exactly the
  kernel of the corresponding component of \(f\), and that component is surjective by
  hypothesis.
\end{proof}

\begin{lemma}[Kernel lattices are finite over a Noetherian ring]
  \label{lem:moduleFinite_aeval_kernelPair}
  \dcref{custom}
  \lean{TwoLatticePair.Hom.moduleFinite_aeval_kernelPair_t₀,
    TwoLatticePair.Hom.moduleFinite_aeval_kernelPair_t₁}
  \uses{def:twoLattice_kernelPair}
  \leanok
  Let \(R\) be Noetherian and let \(f : P \to P'\) be a map of pairs where \(P\) has
  finite lattices. Then the kernel pair \(K(f)\) has finite lattices.
\end{lemma}
\begin{proof}
  \leanok
  \(\ker f_0\) is an \(R[t]\)-submodule of \(M_0\) for the structure through \(t_0\)
  (it is stable under \(t_0\)). By Hilbert's basis theorem \(R[t]\) is Noetherian, so
  the finitely generated \(R[t]\)-module \(M_0\) is a Noetherian module, and every
  submodule --- in particular \(\ker f_0\) --- is finitely generated. Symmetrically
  for \(\ker f_1\).
\end{proof}

\begin{theorem}[COH-1: coherence of \(H^1\) over any base ring]
  \label{thm:twoLattice_coh1}
  \dcref{custom}
  \lean{TwoLatticePair.moduleFinite_h1}
  \uses{def:twoLatticePair, def:twoLattice_H1}
  \leanok
  Let \(R\) be \emph{any} commutative ring --- no Noetherian hypothesis --- and let
  \(P\) be a two-lattice pair with finite lattices. Then \(H^1(P)\) is a finite
  \(R\)-module.
\end{theorem}
\begin{proof}
  \leanok
  \uses{thm:exists_hom_model_surjective, thm:moduleFinite_h1_model,
    lem:twoLattice_h1Map_surjective}
  Choose a finite model surjection \(f : E(m) \to P\)
  (\ref{thm:exists_hom_model_surjective}). It is surjective on overlap modules, so
  \(H^1(f) : H^1(E(m)) \to H^1(P)\) is surjective
  (\ref{lem:twoLattice_h1Map_surjective}). The source is a finite \(R\)-module --- the
  monomial window (\ref{thm:moduleFinite_h1_model}) --- and a quotient of a finite
  module is finite.
\end{proof}

\begin{theorem}[COH-0: coherence of \(H^0\) over a Noetherian base ring]
  \label{thm:twoLattice_coh0}
  \dcref{custom}
  \lean{TwoLatticePair.moduleFinite_h0}
  \uses{def:twoLatticePair, def:twoLattice_H0}
  \leanok
  Let \(R\) be a \emph{Noetherian} commutative ring and let \(P\) be a two-lattice
  pair with finite lattices. Then \(H^0(P)\) is a finite \(R\)-module.
\end{theorem}
\begin{proof}
  \leanok
  \uses{thm:exists_hom_model_surjective, def:twoLattice_kernelPair,
    lem:twoLattice_kernelPair_shortExact, lem:moduleFinite_aeval_kernelPair,
    lem:moduleFinite_aeval_model, def:twoLattice_delta, thm:twoLattice_exact_h0_delta,
    thm:twoLattice_coh1, thm:moduleFinite_h0_model}
  Choose a finite model surjection \(f : E(m) \to P\)
  (\ref{thm:exists_hom_model_surjective}), surjective on all three carriers, and form
  the short exact sequence \(0 \to K(f) \to E(m) \to P \to 0\) of
  \ref{lem:twoLattice_kernelPair_shortExact}. Its connecting homomorphism
  \(\partial : H^0(P) \to H^1(K(f))\) (\ref{def:twoLattice_delta}) has:
  \begin{itemize}
    \item kernel equal to the image of \(H^0(E(m)) \to H^0(P)\)
      (\ref{thm:twoLattice_exact_h0_delta}), a quotient of the finite \(R\)-module
      \(H^0(E(m))\) of the degree window (\ref{thm:moduleFinite_h0_model}, using that
      \(R\) is Noetherian), hence finite;
    \item image a submodule of \(H^1(K(f))\). The model has finite lattices
      (\ref{lem:moduleFinite_aeval_model}), so the kernel pair has finite lattices
      over the Noetherian \(R\) (\ref{lem:moduleFinite_aeval_kernelPair}), and COH-1
      applied to \(K(f)\) (\ref{thm:twoLattice_coh1}) makes \(H^1(K(f))\) a finite
      \(R\)-module; over the Noetherian \(R\) it is a Noetherian module, so the
      submodule \(\partial\bigl(H^0(P)\bigr)\) is finite.
  \end{itemize}
  Thus \(H^0(P)\) is an extension of the finite module \(\partial(H^0(P))\) by the
  finite module \(\ker \partial\), hence finite.
\end{proof}

\section{The rigid engine on a packaged two-cover}
\label{sec:rigid_engine_assembly}

This section assembles the module-level rigid engine --- the coherence theorems of
two-lattice pairs (\ref{thm:twoLattice_coh1}, \ref{thm:twoLattice_coh0}), the fibrewise
Nakayama propagation (Section~\ref{sec:nakayama_vanishing}) and the split rigidity
(Section~\ref{sec:split_rigidity}) --- at the level of sheaves. For a sheaf \(F\) of
\(R\)-modules on the small Zariski site of a scheme \(X\), packaged quasi-coherently
(\ref{def:qcohOn}) on both charts of an affine two-cover \(U_0 \cup U_1 = X\), a small
amount of coordinate data produces a two-lattice pair whose cohomology \emph{is} the
cohomology of the site with coefficients in \(F\); firing the engine on the pair yields
the curve-lite form of Kleiman's theorem on the module \(\mathcal Q\): vanishing of the
degree-one fibres forces \(H^1 = 0\) and makes \(H^0\) finite projective, with formation
of \(H^0\) commuting with all coefficients.

The degree-one carrier is the landed two-cover cokernel
(\ref{thm:twoCoverH1LinearEquiv}); the missing degree-zero carrier --- the kernel twin of
the cokernel bridge --- is built first.

\begin{theorem}[The Mayer--Vietoris degree-zero kernel bridge]
  \label{thm:sectionsKerEquiv}
  \dcref{custom}
  \lean{CategoryTheory.GrothendieckTopology.MayerVietorisSquare.sectionsToKerModuleDiff,
    CategoryTheory.GrothendieckTopology.MayerVietorisSquare.sectionsToKerModuleDiff_coe,
    CategoryTheory.GrothendieckTopology.MayerVietorisSquare.sectionsKerEquiv,
    CategoryTheory.GrothendieckTopology.MayerVietorisSquare.sectionsKerEquiv_apply_coe}
  \uses{def:moduleDiff, lem:mv_exact_degree0}
  \leanok
  For a Mayer--Vietoris square \(S\) and a sheaf \(F\) of \(R\)-modules, the pair of
  restrictions, corestricted to the kernel of the restriction-difference map
  (\ref{def:moduleDiff}), is an \(R\)-linear equivalence
  \[
    F(X_4) \;\cong\; \ker\bigl(\mathrm{moduleDiff} : F(X_2) \times F(X_3) \to
      F(X_1)\bigr),
    \qquad s \longmapsto \bigl(s|_{X_2},\, s|_{X_3}\bigr).
  \]
\end{theorem}
\begin{proof}
  \leanok
  The map lands in the kernel and is injective and surjective onto it by degree-zero
  exactness (\ref{lem:mv_exact_degree0}): injectivity is the separation half of the
  sheaf condition of the square, surjectivity onto the kernel is the gluing half.
  Linearity is linearity of the two restrictions.
\end{proof}

\begin{theorem}[The two-cover \(H^0\) kernel bridge]
  \label{thm:twoCoverH0LinearEquiv}
  \dcref{custom}
  \lean{AlgebraicGeometry.Scheme.twoCoverH0LinearEquiv}
  \uses{thm:sectionsKerEquiv, def:twoCoverSquare, def:HModule_linearEquivHModule',
    def:HModule'_linearEquiv0}
  \leanok
  Let \(X\) be a scheme, \(U_0 \cup U_1 = X\) a two-cover by opens, and \(F\) \emph{any}
  sheaf of \(R\)-modules on the small Zariski site. Then degree-zero cohomology of the
  site is the kernel of the two-cover restriction-difference map: an \(R\)-linear
  equivalence
  \[
    H^0(X, F) \;\cong\;
      \ker\Bigl(F(U_0) \times F(U_1) \longrightarrow F(U_0 \cap U_1)\Bigr).
  \]
  No affineness is needed: degree-zero exactness is the sheaf condition alone.
\end{theorem}
\begin{proof}
  \leanok
  Identify \(H^0(X, F)\) with \(H'^0(\top, F) = F(\top)\)
  (\ref{def:HModule_linearEquivHModule'}, \ref{def:HModule'_linearEquiv0}), then apply
  the kernel bridge \ref{thm:sectionsKerEquiv} to the two-cover square
  (\ref{def:twoCoverSquare}), whose restriction-difference map is exactly the displayed
  one.
\end{proof}

\begin{definition}[The two-cover pair data of a packaged sheaf]
  \label{def:twoCoverPairData}
  \dcref{custom}
  \lean{AlgebraicGeometry.Scheme.TwoCoverPairData}
  \uses{def:qcohOn}
  \leanok
  Let \(F\) be a sheaf of \(R\)-modules on the small Zariski site of \(X\), packaged
  quasi-coherently (\ref{def:qcohOn}) on two opens \(U_0, U_1\). \emph{Two-cover pair
  data} for \(F\) consists of chart coordinates \(g_0 \in \Gamma(X, U_0)\),
  \(g_1 \in \Gamma(X, U_1)\) such that:
  \begin{itemize}
    \item the overlap is the basic open of either coordinate:
      \(U_0 \cap U_1 = D(g_0) = D(g_1)\);
    \item the packaged actions of \(g_0\) and of \(g_1\) commute with the \(R\)-scalars
      on sections of every open \(W \le U_i\);
    \item on the overlap the two packaged actions are mutually inverse (the relation
      \(g_0\, g_1 = 1\), expressed on the actions: acting by \(g_0\) and then by
      \(g_1\), in either order, is the identity of \(F(U_0 \cap U_1)\)).
  \end{itemize}
\end{definition}

\begin{definition}[The two-lattice pair of a packaged sheaf]
  \label{def:twoCoverPairData_pair}
  \dcref{custom}
  \lean{AlgebraicGeometry.Scheme.TwoCoverPairData.end₀,
    AlgebraicGeometry.Scheme.TwoCoverPairData.end₁,
    AlgebraicGeometry.Scheme.TwoCoverPairData.end₀_apply,
    AlgebraicGeometry.Scheme.TwoCoverPairData.end₁_apply,
    AlgebraicGeometry.Scheme.TwoCoverPairData.end₀_pow_apply,
    AlgebraicGeometry.Scheme.TwoCoverPairData.end₁_pow_apply,
    AlgebraicGeometry.Scheme.TwoCoverPairData.pair,
    AlgebraicGeometry.Scheme.TwoCoverPairData.pair_diff,
    AlgebraicGeometry.Scheme.TwoCoverPairData.pair_t₀,
    AlgebraicGeometry.Scheme.TwoCoverPairData.pair_t₁}
  \uses{def:twoCoverPairData, def:twoLatticePair, def:qcohOn_module,
    def:qcohOn_secRes_linear, thm:qcohOn_isLocalizedModule, def:twoCoverSquare,
    def:moduleDiff}
  \leanok
  Let \(U_0, U_1\) be \emph{affine} opens and \(d\) two-cover pair data for a packaged
  sheaf \(F\). Then the data
  \[
    M_0 = F(U_0), \quad M_1 = F(U_1), \quad N = F(U_0 \cap U_1),
  \]
  with coordinate actions the packaged actions of \(g_0, g_1\) (made \(R\)-linear
  endomorphisms by the \(R\)-linearity field of \(d\)), overlap automorphism the action
  of \(g_0\) on \(N\) (invertible with inverse the action of \(g_1\), by the
  mutual-inverse field), and restrictions \(\iota_0, \iota_1\) the section restrictions
  to the overlap, is a two-lattice pair (\ref{def:twoLatticePair}) over \(R\). Its
  difference map is the Mayer--Vietoris restriction-difference map of the two-cover
  square with coefficients in \(F\), on the nose.
\end{definition}
\begin{proof}
  \leanok
  The intertwining laws \(\iota_i(g_i \cdot x) = g_i \cdot \iota_i(x)\) are the
  (P1) naturality of the packaged action under restriction. The localization axioms are
  the (P2) axioms of the packaging, transported to the overlap along the identification
  \(U_0 \cap U_1 = U_i \cap D(g_i)\): denominator clearing for the pair asks that every
  \(n \in N\) admit \(m \ge 0\) and \(x \in M_i\) with \(g_i^m \cdot n =
  \iota_i(x)\), which is (P2) denominator clearing on the affine open \(U_i\); and
  annihilation for the pair asks that \(\iota_i(x) = 0\) force \(g_i^m \cdot x = 0\),
  which is (P2) defect annihilation --- exactly the three-axiom match by which the
  packaging exhibits the overlap sections as the localization of the chart sections
  (\ref{thm:qcohOn_isLocalizedModule}). Powers of the action are the actions of powers
  (by (P1) associativity), which is used implicitly throughout. The difference map
  \((x, y) \mapsto \iota_0 x - \iota_1 y\) is by construction the
  restriction-difference map of \ref{def:moduleDiff} for the two-cover square
  (\ref{def:twoCoverSquare}).
\end{proof}

\begin{definition}[The cohomology carriers of the pair]
  \label{def:twoCoverPair_h0h1}
  \dcref{custom}
  \lean{AlgebraicGeometry.Scheme.TwoCoverPairData.h0Equiv,
    AlgebraicGeometry.Scheme.TwoCoverPairData.h1Equiv}
  \uses{def:twoCoverPairData_pair, thm:twoCoverH0LinearEquiv,
    thm:twoCoverH1LinearEquiv, thm:hModule1_vanishing_qcoh, def:twoLattice_H0,
    def:twoLattice_H1, lem:twoLattice_range_diff}
  \leanok
  Let \(U_0, U_1\) be affine opens with \(U_0 \cup U_1 = X\) and \(d\) two-cover pair
  data for the packaged sheaf \(F\), with pair \(P\)
  (\ref{def:twoCoverPairData_pair}). There are \(R\)-linear equivalences
  \[
    H^0(X, F) \;\cong\; H^0(P), \qquad H^1(X, F) \;\cong\; H^1(P) :
  \]
  the cohomology of the site with coefficients in \(F\) is the cohomology of the
  two-term complex of the pair.
\end{definition}
\begin{proof}
  \leanok
  Degree zero: the kernel bridge (\ref{thm:twoCoverH0LinearEquiv}) lands in the kernel
  of the restriction-difference map, which is the kernel of the pair's difference map
  since the two maps are equal (\ref{def:twoCoverPairData_pair}), i.e.\ \(H^0(P)\)
  (\ref{def:twoLattice_H0}). Degree one: the packaging on the two affine charts gives
  the twisted affine vanishing \(H'^1(U_i, F) = 0\)
  (\ref{thm:hModule1_vanishing_qcoh}), so the two-cover computation
  (\ref{thm:twoCoverH1LinearEquiv}) identifies \(H^1(X, F)\) with the cokernel of the
  restriction-difference map; and the cokernel of the pair's difference map is
  \(H^1(P)\) (\ref{lem:twoLattice_range_diff}, \ref{def:twoLattice_H1}).
\end{proof}

\begin{lemma}[Vanishing of \(H^1\) of the pair makes the differential surjective]
  \label{lem:twoCoverPair_surjective_diff}
  \dcref{custom}
  \lean{AlgebraicGeometry.Scheme.TwoCoverPairData.surjective_diff}
  \uses{def:twoCoverPairData_pair, lem:twoLattice_range_diff,
    lem:rigid_surjective_of_coker_zero}
  \leanok
  If \(H^1\) of the pair of \ref{def:twoCoverPairData_pair} is trivial (a
  subsingleton), then the difference map \(\delta : F(U_0) \times F(U_1) \to
  F(U_0 \cap U_1)\) is surjective.
\end{lemma}
\begin{proof}
  \leanok
  \(H^1\) of the pair is the cokernel of \(\delta\) (\ref{lem:twoLattice_range_diff}),
  and a map with trivial cokernel is surjective
  (\ref{lem:rigid_surjective_of_coker_zero}).
\end{proof}

\begin{theorem}[The rigid engine, sheaf level]
  \label{thm:twoCoverPair_rigidEngine}
  \lean{AlgebraicGeometry.Scheme.TwoCoverPairData.rigidEngine}
  \uses{def:twoCoverPair_h0h1, def:twoCoverPairData_pair}
  \leanok
  \dcref{custom}
  Let \(R\) be Noetherian, \(U_0, U_1\) affine opens covering \(X\), and \(d\)
  two-cover pair data for the packaged sheaf \(F\) whose pair \(P\) has finite lattices
  (the chart sections \(F(U_i)\) are finite over \(R[t]\) through the coordinate
  actions). Assume the section modules are tame: \(F(U_0) \times F(U_1)\) is a flat
  \(R\)-module and \(F(U_0 \cap U_1)\) is a projective \(R\)-module (in the intended
  applications all three are free). If every residue-field fibre of \(H^1\) of the pair
  vanishes,
  \[
    H^1(P) \otimes_R \kappa(\mathfrak p) = 0 \quad \text{for every prime }
      \mathfrak p \subseteq R,
  \]
  then
  \[
    H^1(X, F) = 0, \qquad \text{and } H^0(X, F) \text{ is a finite projective }
      R\text{-module}.
  \]
  This is the curve-lite form of Kleiman's criterion for the module \(\mathcal Q\):
  vanishing of the degree-one cohomology of the fibres ((v) of the source) forces the
  representing module to be free/projective ((i)) --- here delivered on a two-chart
  cover with the two-term \v Cech complex playing the role of the pushforward. The
  fibre hypothesis is deliberately consumed in this \emph{complex} form (fibres of
  \(H^1\) of the pair); its sheaf-level discharge --- identifying the fibre of the
  complex with the cohomology of a fibre sheaf --- is a seam of a later campaign and is
  not part of this statement.
\end{theorem}
\begin{proof}
  \leanok
  \uses{thm:twoLattice_coh1, thm:twoLattice_coh0, thm:fibrewise_vanishing_propagates,
    lem:twoCoverPair_surjective_diff, thm:rigid_projective_ker}
  \(H^1(P)\) is a finite \(R\)-module by COH-1 (\ref{thm:twoLattice_coh1}), so the
  fibrewise vanishing hypothesis propagates to \(H^1(P) = 0\) by Nakayama
  (\ref{thm:fibrewise_vanishing_propagates}), and along the degree-one carrier
  (\ref{def:twoCoverPair_h0h1}) this gives \(H^1(X, F) = 0\). The differential is then
  surjective (\ref{lem:twoCoverPair_surjective_diff}). \(H^0(P) = \ker\delta\) is
  finite over the Noetherian \(R\) by COH-0 (\ref{thm:twoLattice_coh0}), and the split
  rigidity of a surjective two-term complex with flat domain and projective codomain
  makes the finite kernel projective (\ref{thm:rigid_projective_ker}). Transport both
  along the degree-zero carrier (\ref{def:twoCoverPair_h0h1}).
\end{proof}

\begin{theorem}[The openness export]
  \label{thm:twoCoverPair_isOpen_vanishing}
  \alignmentcheck{TO CHECK: the cited Kleiman argument uses a Noetherian base, whereas this
  export claims the same openness over an arbitrary commutative ring via finite-module
  coherence; verify the source-to-blueprint hypothesis change.}
  \lean{AlgebraicGeometry.Scheme.TwoCoverPairData.rigidEngine_isOpen_vanishing}
  \uses{def:twoCoverPairData_pair, thm:twoLattice_coh1, thm:isOpen_fibre_vanishing}
  \leanok
  \dcref{custom}
  In the setting of \ref{def:twoCoverPairData_pair}, with finite lattices but \emph{no}
  Noetherian hypothesis on \(R\), the fibrewise-vanishing locus
  \[
    \bigl\{\, \mathfrak p \in \Spec R :
      H^1(P) \otimes_R \kappa(\mathfrak p) = 0 \,\bigr\}
  \]
  is open in \(\Spec R\). (In the source the analogous openness is obtained from
  Serre's finiteness theorem over a Noetherian base; here it falls out of the coherence
  of \(H^1\) over an arbitrary base ring, Noetherian-free.)
\end{theorem}
\begin{proof}
  \leanok
  \(H^1(P)\) is a finite \(R\)-module over any commutative \(R\)
  (\ref{thm:twoLattice_coh1}); for a finite module the fibrewise-vanishing locus is the
  complement of the support, which is closed (\ref{thm:isOpen_fibre_vanishing}).
\end{proof}

\begin{definition}[The module-coefficient clause]
  \label{def:twoCoverPair_h0TensorEquiv}
  \lean{AlgebraicGeometry.Scheme.TwoCoverPairData.h0TensorEquiv}
  \uses{def:twoCoverPair_h0h1, lem:twoCoverPair_surjective_diff, thm:rigid_ker_tensor}
  \leanok
  \dcref{custom}
  In the setting of \ref{def:twoCoverPair_h0h1}, assume \(H^1(P)\) is trivial and the
  overlap module \(F(U_0 \cap U_1)\) is flat over \(R\). Then for \emph{every}
  \(R\)-module \(P'\) there is an \(R\)-linear equivalence
  \[
    H^0(X, F) \otimes_R P' \;\cong\;
      \ker\bigl(\delta \otimes_R P'\bigr) :
  \]
  formation of \(H^0\) commutes with all module coefficients. This is the two-chart
  form of the isomorphism (eq.\ Q of the source)
  \(\mathrm{Hom}(\mathcal Q, \mathcal N) \cong f_*(\mathcal F \otimes f^*\mathcal N)\)
  in its right-exact regime: on the vanishing locus the functor
  \(P' \mapsto \ker(\delta \otimes P')\) is represented by \(H^0\).
\end{definition}
\begin{proof}
  \leanok
  The differential is surjective (\ref{lem:twoCoverPair_surjective_diff}) with flat
  codomain, so the kernel-tensor rigidity (\ref{thm:rigid_ker_tensor}) gives
  \(\ker(\delta) \otimes P' \cong \ker(\delta \otimes P')\); compose with the
  degree-zero carrier (\ref{def:twoCoverPair_h0h1}) tensored with \(P'\).
\end{proof}

\begin{definition}[The ring-map clause, abstract half]
  \label{def:twoCoverPair_h0BaseChangeEquiv}
  \dcref{custom}
  \lean{AlgebraicGeometry.Scheme.TwoCoverPairData.h0BaseChangeEquiv}
  \uses{def:twoCoverPair_h0h1, lem:twoCoverPair_surjective_diff, thm:rigid_ker_baseChange}
  \leanok
  In the setting of \ref{def:twoCoverPair_h0TensorEquiv}, for every commutative
  \(R\)-algebra \(R'\) there is an \(R'\)-\emph{linear} equivalence
  \[
    R' \otimes_R H^0(X, F) \;\cong\; \ker\bigl(\delta \otimes_R R'\bigr),
  \]
  the module-coefficient clause upgraded to carry the \(R'\)-structure. Reaching
  \(H^0\) of a base-changed \emph{sheaf} on the nose is a separate step: it threads
  this equivalence through term identifications of the base-changed complex
  (\ref{def:relTwistH0BaseChange} below, for the relative twisted sheaf).
\end{definition}
\begin{proof}
  \leanok
  As in \ref{def:twoCoverPair_h0TensorEquiv}, with the \(R'\)-linear kernel base-change
  form of the rigidity (\ref{thm:rigid_ker_baseChange}) in place of the module form.
\end{proof}

\section{The pair data of the twisted sheaf}
\label{sec:twist_pair_data}

The Layer-2 discharge of the pair data for the cocycle-glued twisted sheaf \(F_g\) of
Section~\ref{sec:twisted_sheaf}, generic in the base ring \(k\): every field of
\ref{def:twoCoverPairData} reduces, through the chart trivialisations, to section
algebra in the structure sheaf.

\begin{definition}[Multiplication by a section as an endomorphism]
  \label{def:mulSectionEnd}
  \dcref{custom}
  \lean{AlgebraicGeometry.Scheme.mulSectionEnd, AlgebraicGeometry.Scheme.mulSectionEnd_apply}
  \uses{def:moduleKSheaf}
  \leanok
  Let \(X\) be a scheme over \(\Spec k\) and \(U\) an open. Multiplication by a fixed
  section \(c \in \Gamma(X, U)\) is a \(k\)-linear endomorphism of \(\Gamma(X, U)\)
  (for the \(k\)-structure through the structure morphism, \ref{def:moduleKSheaf}):
  the comparison endomorphism through which the twisted pair's coordinate actions are
  identified with honest section multiplication.
\end{definition}

\begin{lemma}[The packaged action of the twisted sheaf, and its \(k\)-linearity]
  \label{lem:twistSheaf_qsmul_eq}
  \dcref{custom}
  \lean{AlgebraicGeometry.twistSheaf_qsmul₀_eq, AlgebraicGeometry.twistSheaf_qsmul₁_eq,
    AlgebraicGeometry.twistSheaf_smul_qsmul₀, AlgebraicGeometry.twistSheaf_smul_qsmul₁}
  \uses{thm:twistSheaf_qcohOn, def:twistTriv, def:qcohOn_transport}
  \leanok
  Fix opens \(V_0, V_1\) and a unit cocycle \(g\), with twisted sheaf \(F_g\). For
  \(W \le V_i\), the packaged action (\ref{thm:twistSheaf_qcohOn}) of
  \(r \in \Gamma(X, V_i)\) on \(F_g(W)\) is
  trivialise--multiply--untrivialise:
  \[
    r \cdot m \;=\; \mathrm{triv}_i^{-1}\bigl(r|_W \cdot \mathrm{triv}_i(m)\bigr),
  \]
  and it commutes with the \(k\)-scalars.
\end{lemma}
\begin{proof}
  \leanok
  The packaging of \(F_g\) on the chart \(V_i\) is by definition the transport
  (\ref{def:qcohOn_transport}) of the structure-sheaf packaging along the chart
  trivialisation (\ref{def:twistTriv}), whose action is exactly the displayed conjugate;
  the first claim is definitional. For \(k\)-linearity: multiplication of sections
  commutes with the \(k\)-scalars (both are multiplications in a commutative ring), and
  the trivialisations are \(k\)-linear, so the conjugate commutes with the scalars too.
\end{proof}

\begin{lemma}[The two trivialisations differ by the cocycle on the overlap]
  \label{lem:twistTriv_inf_eq}
  \dcref{custom}
  \lean{AlgebraicGeometry.twistTriv₀_inf_eq, AlgebraicGeometry.twistTriv₁_inf_eq}
  \uses{def:twistTriv, def:twistSections}
  \leanok
  On the overlap sections \(F_g(V_0 \cap V_1)\),
  \[
    \mathrm{triv}_0(p) \;=\; g \cdot \mathrm{triv}_1(p), \qquad
    \mathrm{triv}_1(p) \;=\; g^{-1} \cdot \mathrm{triv}_0(p).
  \]
\end{lemma}
\begin{proof}
  \leanok
  A twisted section \(p = (s_0, s_1)\) over \(V_0 \cap V_1\) satisfies the matching
  relation \(s_0| = g \cdot s_1|\) on \((V_0 \cap V_1) \cap V_0 \cap V_1 = V_0 \cap
  V_1\) (\ref{def:twistSections}); the two trivialisations are the two components
  (\ref{def:twistTriv}), so the relation reads \(\mathrm{triv}_0(p) = g \cdot
  \mathrm{triv}_1(p)\). Multiply by \(g^{-1}\) for the second form.
\end{proof}

\begin{lemma}[The mutual-inverse law of the coordinate actions]
  \label{lem:twistSheaf_qsmul_inverse}
  \dcref{custom}
  \lean{AlgebraicGeometry.twistSheaf_qsmul₀_qsmul₁, AlgebraicGeometry.twistSheaf_qsmul₁_qsmul₀}
  \uses{lem:twistSheaf_qsmul_eq, lem:twistTriv_inf_eq}
  \leanok
  Let \(r_0 \in \Gamma(X, V_0)\), \(r_1 \in \Gamma(X, V_1)\) satisfy
  \(r_0|_{\cap} \cdot r_1|_{\cap} = 1\) on the overlap. Then the packaged actions of
  \(r_0\) and \(r_1\) on \(F_g(V_0 \cap V_1)\) are mutually inverse.
\end{lemma}
\begin{proof}
  \leanok
  Compute the composite on \(m\) through the chart-0 trivialisation. Acting by \(r_1\)
  is \(\mathrm{triv}_1^{-1}(r_1| \cdot \mathrm{triv}_1(m))\)
  (\ref{lem:twistSheaf_qsmul_eq}); by \ref{lem:twistTriv_inf_eq} its chart-0
  trivialisation is \(g \cdot r_1| \cdot g^{-1} \cdot \mathrm{triv}_0(m)
  = r_1| \cdot \mathrm{triv}_0(m)\) (the ring is commutative). Acting then by \(r_0\)
  multiplies by \(r_0|\), giving \(r_0|\, r_1| \cdot \mathrm{triv}_0(m) =
  \mathrm{triv}_0(m)\); untrivialising returns \(m\). The other order is symmetric,
  through the chart-1 trivialisation.
\end{proof}

\begin{definition}[The pair data of the twisted sheaf]
  \label{def:twistPairData}
  \dcref{custom}
  \lean{AlgebraicGeometry.twistPairData}
  \uses{def:twoCoverPairData, thm:twistSheaf_qcohOn, lem:twistSheaf_qsmul_eq,
    lem:twistSheaf_qsmul_inverse}
  \leanok
  Let \(r_0 \in \Gamma(X, V_0)\) and \(r_1 \in \Gamma(X, V_1)\) satisfy
  \(V_0 \cap V_1 = D(r_0) = D(r_1)\) and \(r_0|_{\cap} \cdot r_1|_{\cap} = 1\). Then
  \((r_0, r_1)\) is two-cover pair data (\ref{def:twoCoverPairData}) for the twisted
  sheaf \(F_g\) on \((V_0, V_1)\): the basic-open identifications are the hypotheses,
  \(k\)-linearity of the packaged actions is \ref{lem:twistSheaf_qsmul_eq}, and the
  mutual-inverse law is \ref{lem:twistSheaf_qsmul_inverse}.
\end{definition}

\begin{theorem}[Chart-lattice finiteness of the twisted pair]
  \label{thm:moduleFinite_aeval_twistPair}
  \dcref{custom}
  \lean{AlgebraicGeometry.moduleFinite_aeval'_twistPair_t₀,
    AlgebraicGeometry.moduleFinite_aeval'_twistPair_t₁}
  \uses{def:twistPairData, def:twoCoverPairData_pair, def:mulSectionEnd, def:twistTriv,
    lem:aevalCongr}
  \leanok
  In the setting of \ref{def:twistPairData} with \(V_0, V_1\) affine: if the chart
  sections \(\Gamma(X, V_i)\) are finite over \(k[t]\) with \(t\) acting as
  multiplication by \(r_i\) (\ref{def:mulSectionEnd}), then the pair of the twisted
  sheaf (\ref{def:twoCoverPairData_pair}) has finite lattices.
\end{theorem}
\begin{proof}
  \leanok
  The chart trivialisation \(\mathrm{triv}_i : F_g(V_i) \cong \Gamma(X, V_i)\)
  (\ref{def:twistTriv}) intertwines the pair's coordinate action --- which is
  trivialise--multiply--untrivialise (\ref{lem:twistSheaf_qsmul_eq}) --- with
  multiplication by \(r_i\). Chart-lattice finiteness transports along an equivariant
  equivalence (\ref{lem:aevalCongr}).
\end{proof}

\begin{theorem}[Freeness of the twisted section modules]
  \label{thm:free_twistSheafSections}
  \dcref{custom}
  \lean{AlgebraicGeometry.free_twistSheafSections₀, AlgebraicGeometry.free_twistSheafSections₁}
  \uses{def:twistTriv}
  \leanok
  For \(W \le V_i\), if \(\Gamma(X, W)\) is a free \(k\)-module then so is
  \(F_g(W)\): transport a basis along the chart trivialisation
  (\ref{def:twistTriv}).
\end{theorem}
\begin{proof}
  \leanok
  A module equivalent to a free module is free.
\end{proof}

\section{The pinned relative cover: coordinates, finiteness and freeness}
\label{sec:relative_cover_coordinates}

The geometric half of the engine's Layer-2 discharge: the pinned affine two-cover of the
curve --- the chart preimages of a finite morphism \(\pi\) to \(\PP^1\) --- carries
polynomial chart algebras through which the finiteness of \(\pi\) becomes chart-lattice
finiteness, first over the field \(k\) and then, by base change, over every test ring
\(R\). Throughout, \(k\) is a field and \(\pi : Y \to \PP^1\) an affine morphism (finite
where stated).

\begin{definition}[The pinned affine two-cover]
  \label{def:fiberTwoCover}
  \dcref{custom}
  \lean{AlgebraicGeometry.fiberTwoCover, AlgebraicGeometry.fiberTwoCover_V₀,
    AlgebraicGeometry.fiberTwoCover_V₁}
  \uses{def:AffineTwoCover, lem:preimage_chart_affine, lem:preimage_chart_sup,
    def:fiberCover, lem:P1_chartOpen_inf}
  \leanok
  For an affine morphism \(\pi : Y \to \PP^1\), the chart preimages
  \(V_0 = \pi^{-1}D_+(X_0)\), \(V_1 = \pi^{-1}D_+(X_1)\) form an \emph{affine two-cover}
  of \(Y\) (\ref{def:AffineTwoCover}): both charts are affine
  (\ref{lem:preimage_chart_affine}), they cover \(Y\)
  (\ref{lem:preimage_chart_sup}), and the overlap
  \(V_0 \cap V_1 = \pi^{-1}D_+(X_0 X_1)\) (\ref{lem:P1_chartOpen_inf}) is affine, being
  the preimage of an affine open under an affine morphism. This is the cover of
  \ref{def:fiberCover}, upgraded with the affineness of the overlap; the whole engine
  runs on it and on its base change.
\end{definition}

\begin{theorem}[The overlap is the basic open of the chart-1 coordinate]
  \label{thm:preimage_inf_eq_basicOpen_fiberCoord1}
  \dcref{custom}
  \lean{AlgebraicGeometry.preimage_inf_eq_basicOpen_fiberCoord₁}
  \uses{def:fiberCover, def:fiberCoordUnit, thm:P1_basicOpen_chartCoord}
  \leanok
  \(V_0 \cap V_1 = D(t_1)\), where \(t_1 = \pi^{\sharp}(X_0/X_1) \in \Gamma(Y, V_1)\)
  is the pulled-back chart-1 coordinate (\ref{def:fiberCoordUnit}) --- the chart-1 twin
  of the identification \(V_0 \cap V_1 = D(t_0)\) of \ref{def:fiberCover}.
\end{theorem}
\begin{proof}
  \leanok
  On \(\PP^1\) the basic open of the chart-1 coordinate \(X_0/X_1\) in \(D_+(X_1)\) is
  the chart overlap (\ref{thm:P1_basicOpen_chartCoord}); basic opens commute with
  pullback of sections, so applying \(\pi^{-1}\) gives
  \(V_1 \cap V_0 = D\bigl(\pi^{\sharp}(X_0/X_1)\bigr)\).
\end{proof}

\begin{definition}[The chart polynomial algebras]
  \label{def:fiberPolyHom}
  \dcref{custom}
  \lean{AlgebraicGeometry.fiberPolyHom₀, AlgebraicGeometry.fiberPolyHom₁,
    AlgebraicGeometry.fiberPolyHom₀_X, AlgebraicGeometry.fiberPolyHom₁_X}
  \uses{def:P1_chartSectionsEquiv0, def:P1_chartSectionsEquiv1, def:fiberCover,
    def:fiberCoordUnit}
  \leanok
  The \emph{chart-\(i\) polynomial algebra} is the ring homomorphism
  \[
    k[t] \;\cong\; \Gamma\bigl(\PP^1, D_+(X_i)\bigr)
      \xrightarrow{\ \pi^{\sharp}\ } \Gamma(Y, V_i),
  \]
  the chart identification (\ref{def:P1_chartSectionsEquiv0},
  \ref{def:P1_chartSectionsEquiv1}) followed by pullback along \(\pi\). It sends \(t\)
  to the pulled-back chart coordinate: \(t \mapsto t_0 = \pi^{\sharp}(X_1/X_0)\) on
  chart \(0\) (\ref{def:fiberCover}) and \(t \mapsto t_1 = \pi^{\sharp}(X_0/X_1)\) on
  chart \(1\) (\ref{def:fiberCoordUnit}).
\end{definition}
\begin{proof}
  \leanok
  Only the values at \(t\) need comment: the chart identification matches \(t\) with
  the chart coordinate \(X_{1-i}/X_i \in \Gamma(\PP^1, D_+(X_i))\), whose pullback is
  \(t_i\) by definition.
\end{proof}

\begin{theorem}[E-i: the chart polynomial algebras are finite ring maps]
  \label{thm:fiberPolyHom_finite}
  \dcref{custom}
  \lean{AlgebraicGeometry.fiberPolyHom₀_finite, AlgebraicGeometry.fiberPolyHom₁_finite}
  \uses{def:fiberPolyHom, lem:finite_app_chart}
  \leanok
  If \(\pi\) is finite, both chart polynomial algebras
  \(k[t] \to \Gamma(Y, V_i)\) are finite ring homomorphisms.
\end{theorem}
\begin{proof}
  \leanok
  For a finite morphism, the induced map on sections over an affine chart of the target
  is a finite ring map (\ref{lem:finite_app_chart}); precomposing with the ring
  isomorphism \(k[t] \cong \Gamma(\PP^1, D_+(X_i))\) preserves finiteness.
\end{proof}

\begin{theorem}[Scalar compatibility of the chart polynomial algebras]
  \label{thm:fiberPolyHom_algebraMap}
  \dcref{custom}
  \lean{AlgebraicGeometry.fiberPolyHom₀_algebraMap, AlgebraicGeometry.fiberPolyHom₁_algebraMap}
  \uses{def:fiberPolyHom, lem:P1_structureMap_appTop_awayToSection,
    lem:appLE_overAlgebraMap, def:overAlgebraMap}
  \leanok
  Let \(Y = C\) be the challenge curve and suppose \(\pi\) factors the structure
  morphism, \(\pi \circ (\text{structure map of } \PP^1) = (C \to \Spec k)\). Then the
  chart polynomial algebras carry constants to the structure algebra map
  (\ref{def:overAlgebraMap}): for \(c \in k\),
  \[
    \text{fiberPolyHom}_i(c) \;=\; \text{the image of } c \text{ under }
      k \to \Gamma(C, V_i).
  \]
\end{theorem}
\begin{proof}
  \leanok
  The chart identification carries the constant \(c \in k[t]\) to the restriction to
  \(D_+(X_i)\) of the structure algebra map of \(\PP^1\)
  (\ref{lem:P1_structureMap_appTop_awayToSection}). Pullback along \(\pi\) carries the
  structure algebra map of \(\PP^1\) to that of \(C\), because \(\pi\) commutes with
  the two structure morphisms (\ref{lem:appLE_overAlgebraMap}, applied to the
  factorization hypothesis).
\end{proof}

\begin{theorem}[E-i in chart-lattice form, field level]
  \label{thm:moduleFinite_aeval_fiberCoord}
  \dcref{custom}
  \lean{AlgebraicGeometry.moduleFinite_aeval'_mulLeft_fiberCoord₀,
    AlgebraicGeometry.moduleFinite_aeval'_mulLeft_fiberCoord₁}
  \uses{thm:fiberPolyHom_finite, thm:fiberPolyHom_algebraMap, lem:moduleFinite_aeval_ringHom}
  \leanok
  In the setting of \ref{thm:fiberPolyHom_algebraMap} with \(\pi\) finite: the chart
  sections \(\Gamma(C, V_i)\) are finite over \(k[t]\), with \(t\) acting as
  multiplication by the pulled-back chart coordinate \(t_i\).
\end{theorem}
\begin{proof}
  \leanok
  Apply the finite-ring-map criterion (\ref{lem:moduleFinite_aeval_ringHom}) to the
  chart polynomial algebra: it is finite (\ref{thm:fiberPolyHom_finite}), sends \(t\)
  to \(t_i\) (\ref{def:fiberPolyHom}), and sends constants to the structure algebra
  map, which is the \(k\)-scalar action on sections
  (\ref{thm:fiberPolyHom_algebraMap}).
\end{proof}

\begin{definition}[The relative chart coordinates]
  \label{def:relFiberCoord}
  \dcref{custom}
  \lean{AlgebraicGeometry.relFiberCoord₀, AlgebraicGeometry.relFiberCoord₁,
    AlgebraicGeometry.relPullbackSection_mul}
  \uses{def:relSectionsBaseChange, def:relCover, def:fiberTwoCover, def:fiberCover,
    def:fiberCoordUnit}
  \leanok
  Let \(R\) be a commutative \(k\)-algebra and consider the base-changed pinned cover
  \(V_i^R = \mathrm{pr}_1^{-1} V_i\) of the relative curve \(C_R\)
  (\ref{def:relCover}, \ref{def:fiberTwoCover}). The \emph{relative chart coordinates}
  are the pullbacks along the first projection,
  \[
    t_0^R \;=\; \mathrm{pr}_1^{\sharp}(t_0) \in \Gamma(C_R, V_0^R), \qquad
    t_1^R \;=\; \mathrm{pr}_1^{\sharp}(t_1) \in \Gamma(C_R, V_1^R)
  \]
  (\ref{def:relSectionsBaseChange}); pullback of sections is multiplicative.
\end{definition}

\begin{theorem}[The relative overlap identifications]
  \label{thm:relCover_inf_eq_basicOpen}
  \dcref{custom}
  \lean{AlgebraicGeometry.relCover_fiberTwoCover_inf_eq_basicOpen₀,
    AlgebraicGeometry.relCover_fiberTwoCover_inf_eq_basicOpen₁}
  \uses{def:relFiberCoord, lem:relCover_inf, def:fiberCover,
    thm:preimage_inf_eq_basicOpen_fiberCoord1}
  \leanok
  On the base-changed pinned cover,
  \[
    V_0^R \cap V_1^R \;=\; D(t_0^R) \;=\; D(t_1^R).
  \]
\end{theorem}
\begin{proof}
  \leanok
  The overlap of the base-changed cover is the preimage of the overlap
  (\ref{lem:relCover_inf}); the overlap downstairs is \(D(t_0) = D(t_1)\)
  (\ref{def:fiberCover}, \ref{thm:preimage_inf_eq_basicOpen_fiberCoord1}); and basic
  opens commute with pullback of sections along \(\mathrm{pr}_1\).
\end{proof}

\begin{theorem}[The relative coordinates are mutually inverse on the overlap]
  \label{thm:relFiberCoord_mul}
  \dcref{custom}
  \lean{AlgebraicGeometry.relFiberCoord_mul}
  \uses{def:relFiberCoord, def:fiberCoordUnit, lem:relCover_inf}
  \leanok
  \(t_0^R|_{\cap} \cdot t_1^R|_{\cap} = 1\) in
  \(\Gamma\bigl(C_R,\ V_0^R \cap V_1^R\bigr)\).
\end{theorem}
\begin{proof}
  \leanok
  The identity \(t_0|_{\cap} \cdot t_1|_{\cap} = 1\) holds downstairs
  (\ref{def:fiberCoordUnit}); pull it back along \(\mathrm{pr}_1^{\sharp}\) over the
  overlap, a ring homomorphism commuting with restriction, using
  \(\mathrm{pr}_1^{-1}(V_0 \cap V_1) = V_0^R \cap V_1^R\) (\ref{lem:relCover_inf}).
\end{proof}

\begin{theorem}[The relative E-i: chart-lattice finiteness over the test ring]
  \label{thm:moduleFinite_aeval_relFiberCoord}
  \dcref{custom}
  \lean{AlgebraicGeometry.moduleFinite_aeval'_mulSectionEnd_relFiberCoord₀,
    AlgebraicGeometry.moduleFinite_aeval'_mulSectionEnd_relFiberCoord₁}
  \uses{thm:moduleFinite_aeval_fiberCoord, thm:moduleFinite_aeval_baseChange,
    lem:aevalCongr, def:relSectionsBaseChange, def:relFiberCoord, def:mulSectionEnd}
  \leanok
  In the setting of \ref{thm:moduleFinite_aeval_fiberCoord}, for every commutative
  \(k\)-algebra \(R\): the relative chart sections \(\Gamma(C_R, V_i^R)\) are finite
  over \(R[t]\), with \(t\) acting as multiplication by the relative coordinate
  \(t_i^R\) (\ref{def:mulSectionEnd}).
\end{theorem}
\begin{proof}
  \leanok
  The field-level chart lattice is finite over \(k[t]\)
  (\ref{thm:moduleFinite_aeval_fiberCoord}); its base change to \(R\) is finite over
  \(R[t]\) (\ref{thm:moduleFinite_aeval_baseChange}). The sections base change
  \(R \otimes_k \Gamma(C, V_i) \cong \Gamma(C_R, V_i^R)\)
  (\ref{def:relSectionsBaseChange}) intertwines the base-changed multiplication
  endomorphism \(\mathrm{id} \otimes (t_i \cdot)\) with multiplication by \(t_i^R\):
  on pure tensors, \(a \otimes t_i s \mapsto \mathrm{pr}_2^{\sharp}(a)\,
  \mathrm{pr}_1^{\sharp}(t_i s) = t_i^R \cdot \bigl(\mathrm{pr}_2^{\sharp}(a)\,
  \mathrm{pr}_1^{\sharp}(s)\bigr)\), since \(\mathrm{pr}_1^{\sharp}\) is multiplicative
  (\ref{def:relFiberCoord}). Chart-lattice finiteness transports along the equivariant
  equivalence (\ref{lem:aevalCongr}).
\end{proof}

\begin{theorem}[Freeness of the relative section modules]
  \label{thm:free_relSections}
  \dcref{custom}
  \lean{AlgebraicGeometry.free_relSections}
  \uses{def:relSectionsBaseChange}
  \leanok
  For every quasi-compact quasi-separated open \(V \subseteq C\) and every commutative
  \(k\)-algebra \(R\), the \(R\)-module \(\Gamma(C_R, \mathrm{pr}_1^{-1}V)\) is free:
  a \(k\)-basis of \(\Gamma(C, V)\) becomes an \(R\)-basis through the sections base
  change (\ref{def:relSectionsBaseChange}). This feeds the flatness and projectivity
  hypotheses of the engine: all three section modules of the pinned relative cover are
  free.
\end{theorem}
\begin{proof}
  \leanok
  \(\Gamma(C, V)\) is a free \(k\)-module (\(k\) is a field), so
  \(R \otimes_k \Gamma(C, V)\) is a free \(R\)-module, and freeness transports along
  the \(R\)-linear equivalence of \ref{def:relSectionsBaseChange}.
\end{proof}

\section{The rigid engine on the relative twisted sheaf}
\label{sec:rel_twist_rigid_engine}

The keystones: the engine of Section~\ref{sec:rigid_engine_assembly} fired on the
relative twisted sheaf \(F_g\) on the base change of the pinned cover, with every
hypothesis discharged from the preceding two sections --- and the base-change clause
threaded, through the term identifications of Section~\ref{sec:relative_h1_base_change},
onto \(H^0\) of the base-changed sheaf on the nose.

One staging note, recorded here because it is the honest state of the statements: the
engine's vanishing hypothesis is consumed in \emph{complex form} --- vanishing of the
residue-field fibres of \(H^1\) of the two-lattice pair of \(F_g\). Its sheaf-level
discharge (identifying the glued fibre sheaf with the divisor sheaf of a witness and
citing the fibrewise large-twist vanishing of Chapter~\ref{chap:RiemannRoch}), and the
rank export \(\dim_{\kappa(\mathfrak p)} H^0 \otimes \kappa(\mathfrak p) = \deg +
\chi\), are deliberately staged behind a later seam and are \emph{not} claimed by the
statements below; consumers supply the fibre hypothesis.

Throughout: \(k\) a field, \(C\) the challenge curve, \(\pi : C \to \PP^1\) finite,
\(R\) a commutative \(k\)-algebra, and \(g\) a unit cocycle on the overlap of the
base-changed pinned cover \((V_0^R, V_1^R)\) of \(C_R\), with relative twisted sheaf
\(F_g\) (\ref{def:relTwistSheaf}).

\begin{definition}[The pair data of the relative twisted sheaf]
  \label{def:relTwistPairData}
  \dcref{custom}
  \lean{AlgebraicGeometry.relTwistSheaf.instQcohOn₀,
    AlgebraicGeometry.relTwistSheaf.instQcohOn₁, AlgebraicGeometry.relTwistPairData}
  \uses{def:relTwistSheaf, def:twistPairData, thm:twistSheaf_qcohOn, def:relFiberCoord,
    thm:relCover_inf_eq_basicOpen, thm:relFiberCoord_mul, def:fiberTwoCover}
  \leanok
  The relative twisted sheaf is packaged quasi-coherently on both base-changed charts
  (\ref{thm:twistSheaf_qcohOn}, read on the cover of \(C_R\)), and the relative chart
  coordinates \((t_0^R, t_1^R)\) satisfy the two hypotheses of \ref{def:twistPairData}
  on the base-changed pinned cover: the overlap-as-basic-open identifications
  (\ref{thm:relCover_inf_eq_basicOpen}) and the mutual-inverse relation
  (\ref{thm:relFiberCoord_mul}). This yields the two-cover pair data of \(F_g\), and
  with it the two-lattice pair of the relative twisted sheaf that the engine consumes.
\end{definition}

\begin{theorem}[Chart-lattice finiteness of the relative twisted pair]
  \label{thm:moduleFinite_aeval_relTwistPair}
  \dcref{custom}
  \lean{AlgebraicGeometry.moduleFinite_aeval'_relTwistPair_t₀,
    AlgebraicGeometry.moduleFinite_aeval'_relTwistPair_t₁}
  \uses{def:relTwistPairData, thm:moduleFinite_aeval_twistPair,
    thm:moduleFinite_aeval_relFiberCoord}
  \leanok
  If \(R\) acts through the structure factorization of \(\pi\) (the hypothesis of
  \ref{thm:fiberPolyHom_algebraMap}), the pair of the relative twisted sheaf
  (\ref{def:relTwistPairData}) has finite lattices over \(R[t]\).
\end{theorem}
\begin{proof}
  \leanok
  The relative E-i gives finiteness of the structure-sheaf chart lattices at the
  relative coordinates (\ref{thm:moduleFinite_aeval_relFiberCoord}); the twisted pair's
  lattices are trivialisation-transports of these
  (\ref{thm:moduleFinite_aeval_twistPair}).
\end{proof}

\begin{theorem}[Freeness of the twisted section modules on the pinned relative cover]
  \label{thm:free_relTwistSections}
  \dcref{custom}
  \lean{AlgebraicGeometry.free_relTwistSections₀, AlgebraicGeometry.free_relTwistSections₁,
    AlgebraicGeometry.free_relTwistSectionsOverlap}
  \uses{thm:free_relSections, thm:free_twistSheafSections, def:fiberTwoCover,
    def:relTwistPairData}
  \leanok
  The three section modules of the pair of the relative twisted sheaf --- the two
  charts and the overlap --- are free \(R\)-modules.
\end{theorem}
\begin{proof}
  \leanok
  The charts and the overlap of the pinned cover are affine
  (\ref{def:fiberTwoCover}), in particular quasi-compact and quasi-separated, so the
  structure-sheaf sections of their base changes are free
  (\ref{thm:free_relSections}); the twisted sections are trivialisation-equivalent to
  them (\ref{thm:free_twistSheafSections}).
\end{proof}

\begin{theorem}[The rigid engine on the relative twisted sheaf]
  \label{thm:relTwistRigidEngine}
  \lean{AlgebraicGeometry.relTwistRigidEngine}
  \uses{thm:twoCoverPair_rigidEngine, def:relTwistPairData,
    thm:moduleFinite_aeval_relTwistPair, thm:free_relTwistSections, lem:relCover_charts}
  \leanok
  \dcref{custom}
  Let \(R\) be Noetherian. If every residue-field fibre of \(H^1\) of the pair of
  \(F_g\) vanishes,
  \[
    H^1(\text{pair}) \otimes_R \kappa(\mathfrak p) = 0 \quad \text{for every prime }
      \mathfrak p \subseteq R,
  \]
  then
  \[
    H^1\bigl(C_R, F_g\bigr) = 0, \qquad \text{and } H^0\bigl(C_R, F_g\bigr)
      \text{ is a finite projective } R\text{-module}.
  \]
  Every other hypothesis of the abstract engine is discharged: this is the curve-lite
  Kleiman criterion, (v)\,\(\Rightarrow\)\,(i), on the relative twisted sheaf. The
  fibre hypothesis remains in complex form (see the staging note at the head of this
  section).
\end{theorem}
\begin{proof}
  \leanok
  The base-changed pinned cover is an affine two-cover of \(C_R\)
  (\ref{lem:relCover_charts}); the pair data is \ref{def:relTwistPairData}; the chart
  lattices are finite (\ref{thm:moduleFinite_aeval_relTwistPair}); the section modules
  are free (\ref{thm:free_relTwistSections}), so the product of the chart modules is
  flat and the overlap module is projective. Fire the sheaf-level engine
  (\ref{thm:twoCoverPair_rigidEngine}).
\end{proof}

\begin{lemma}[Fibrewise vanishing propagates to the pair]
  \label{lem:relTwist_subsingleton_pairH1}
  \dcref{custom}
  \lean{AlgebraicGeometry.relTwist_subsingleton_pairH1}
  \uses{thm:moduleFinite_aeval_relTwistPair, thm:twoLattice_coh1,
    thm:fibrewise_vanishing_propagates}
  \leanok
  With no Noetherian hypothesis on \(R\): if every residue-field fibre of \(H^1\) of
  the pair of \(F_g\) vanishes, then \(H^1\) of the pair vanishes. (Consumers holding
  the fibrewise hypothesis use this to feed the two clauses below, which are stated on
  the vanishing locus.)
\end{lemma}
\begin{proof}
  \leanok
  \(H^1\) of the pair is a finite \(R\)-module by COH-1
  (\ref{thm:twoLattice_coh1}, with the finite lattices of
  \ref{thm:moduleFinite_aeval_relTwistPair}); fibrewise vanishing of a finite module
  propagates by Nakayama (\ref{thm:fibrewise_vanishing_propagates}).
\end{proof}

\begin{theorem}[The openness export on the relative twisted sheaf]
  \label{thm:relTwistRigidEngine_isOpen}
  \lean{AlgebraicGeometry.relTwistRigidEngine_isOpen_vanishing}
  \uses{thm:twoCoverPair_isOpen_vanishing, thm:moduleFinite_aeval_relTwistPair}
  \leanok
  \dcref{custom}
  With no Noetherian hypothesis on \(R\), the fibrewise-vanishing locus
  \(\{\mathfrak p \in \Spec R : H^1(\text{pair of } F_g) \otimes_R
  \kappa(\mathfrak p) = 0\}\) is open in \(\Spec R\) --- the strata over which the
  engine fires.
\end{theorem}
\begin{proof}
  \leanok
  The openness export of the abstract engine
  (\ref{thm:twoCoverPair_isOpen_vanishing}), with the finite lattices of
  \ref{thm:moduleFinite_aeval_relTwistPair}.
\end{proof}

\begin{definition}[The module-coefficient clause on the relative twisted sheaf]
  \label{def:relTwistH0TensorEquiv}
  \lean{AlgebraicGeometry.relTwistH0TensorEquiv}
  \uses{def:twoCoverPair_h0TensorEquiv, thm:free_relTwistSections, def:relTwistPairData}
  \leanok
  \dcref{custom}
  On the vanishing locus (\(H^1\) of the pair of \(F_g\) trivial), for every
  \(R\)-module \(P\) there is an \(R\)-linear equivalence
  \[
    H^0\bigl(C_R, F_g\bigr) \otimes_R P \;\cong\;
      \ker\bigl(\delta_{F_g} \otimes_R P\bigr),
  \]
  \(\delta_{F_g}\) the twisted two-cover differential: formation of the twisted
  \(H^0\) commutes with all module coefficients (the curve-lite eq.\ Q).
\end{definition}
\begin{proof}
  \leanok
  The abstract module-coefficient clause (\ref{def:twoCoverPair_h0TensorEquiv}), the
  flatness of the overlap module coming from its freeness
  (\ref{thm:free_relTwistSections}).
\end{proof}

\begin{definition}[The ring-map clause on the relative twisted sheaf, abstract half]
  \label{def:relTwistH0BaseChangeEquiv}
  \dcref{custom}
  \lean{AlgebraicGeometry.relTwistH0BaseChangeEquiv}
  \uses{def:twoCoverPair_h0BaseChangeEquiv, thm:free_relTwistSections, def:relTwistPairData}
  \leanok
  On the vanishing locus, for every commutative \(R\)-algebra \(R'\) there is an
  \(R'\)-linear equivalence
  \[
    R' \otimes_R H^0\bigl(C_R, F_g\bigr) \;\cong\;
      \ker\bigl(\delta_{F_g} \otimes_R R'\bigr).
  \]
\end{definition}
\begin{proof}
  \leanok
  The abstract ring-map clause (\ref{def:twoCoverPair_h0BaseChangeEquiv}), with
  flatness of the overlap module from freeness (\ref{thm:free_relTwistSections}).
\end{proof}

\begin{lemma}[The twisted differential in trivialised coordinates]
  \label{lem:twistTriv_pairDiff}
  \dcref{custom}
  \lean{AlgebraicGeometry.twistTriv₀_pairDiff}
  \uses{def:twistTriv, lem:twistTriv_res, lem:twistTriv_inf_eq}
  \leanok
  For any twisted sheaf \(F_g\) on opens \(V_0, V_1\): under the chart-0
  trivialisation of the overlap, the twisted restriction-difference map
  \((p_0, p_1) \mapsto p_0|_{\cap} - p_1|_{\cap}\) becomes
  \[
    (s_0, s_1) \longmapsto s_0|_{\cap} - g \cdot s_1|_{\cap},
    \qquad s_i = \mathrm{triv}_i(p_i) :
  \]
  the defect-map shape of \ref{def:twistSections}, with the cocycle carrying the
  chart-1 term across.
\end{lemma}
\begin{proof}
  \leanok
  The trivialisations commute with restriction (\ref{lem:twistTriv_res}), so
  \(\mathrm{triv}_0(p_0|_{\cap}) = s_0|_{\cap}\); and on the overlap the chart-0
  trivialisation of \(p_1|_{\cap}\) is \(g\) times its chart-1 trivialisation
  (\ref{lem:twistTriv_inf_eq}), which is \(s_1|_{\cap}\), again by
  \ref{lem:twistTriv_res}. Subtract.
\end{proof}

\begin{definition}[CBC-1 for the twisted complex, assembled]
  \label{def:relTwistDomBaseChange}
  \dcref{custom}
  \lean{AlgebraicGeometry.relTwistDomBaseChange,
    AlgebraicGeometry.twistTriv₀_relTwistTermBaseChange₀,
    AlgebraicGeometry.twistTriv₁_relTwistTermBaseChange₁,
    AlgebraicGeometry.twistTriv₀_relTwistOverlapBaseChange}
  \uses{def:relTwistTermBaseChange, thm:relTermBaseChange, def:relSectionsMap,
    def:twistTriv, lem:relDiffBaseChange}
  \leanok
  Let \(R \to R'\) be a homomorphism of commutative \(k\)-algebras, \(D\) an affine
  two-chart cover of \(C\), \(g\) a unit cocycle on the base-changed cover of \(C_R\)
  and \(g' = \rho^{\sharp}(g)\) its base change. The product of the two chart term
  equivalences of \ref{def:relTwistTermBaseChange} is an \(R'\)-linear equivalence
  \[
    R' \otimes_R \bigl(F_g(V_0^R) \times F_g(V_1^R)\bigr) \;\cong\;
      F_{g'}(V_0^{R'}) \times F_{g'}(V_1^{R'}),
  \]
  the degree-zero term of the twisted complex base-changed. In trivialised coordinates
  each term equivalence computes, on a pure tensor \(r' \otimes p\), as
  \(r' \cdot \rho^{\sharp}\bigl(\mathrm{triv}(p)\bigr)\) --- the structure-sheaf term
  base change of \ref{thm:relTermBaseChange} (at the overlap, the codomain equivalence
  of \ref{lem:relDiffBaseChange}) read through the trivialisations.
\end{definition}
\begin{proof}
  \leanok
  Each twisted term equivalence is by construction
  (\ref{def:relTwistTermBaseChange}) the conjugate of the corresponding
  structure-sheaf term equivalence by the chart trivialisations over \(R\) and
  \(R'\); the displayed formula on pure tensors is the formula of
  \ref{thm:relTermBaseChange} for the structure sheaf, conjugated. The product of two
  \(R'\)-linear equivalences is one.
\end{proof}

\begin{theorem}[The twisted \(\delta\)-naturality square]
  \label{thm:relTwistDiffBaseChange}
  \dcref{custom}
  \lean{AlgebraicGeometry.relTwistDiff, AlgebraicGeometry.relTwistDiff_apply,
    AlgebraicGeometry.relTwistPair_diff, AlgebraicGeometry.relTwistDiffBaseChange,
    AlgebraicGeometry.relTwistPairDiffBaseChange}
  \uses{lem:twistTriv_pairDiff, def:relTwistDomBaseChange, def:relCocycleBaseChange,
    lem:relSectionsMap_props, def:relTwistTermBaseChange, def:relTwistSheaf}
  \leanok
  In the setting of \ref{def:relTwistDomBaseChange}, write \(\delta_g\) for the
  twisted restriction-difference map over \(R\) and \(\delta_{g'}\) for the one over
  \(R'\). Under the CBC-1 term identifications --- the domain equivalence of
  \ref{def:relTwistDomBaseChange} and the overlap equivalence of
  \ref{def:relTwistTermBaseChange} --- the base change of the twisted differential is
  carried to the twisted differential of the base-changed cocycle:
  \[
    \mathrm{overlap}\circ\bigl(\delta_g \otimes_R R'\bigr)
      \;=\; \delta_{g'} \circ \mathrm{dom} .
  \]
  This is the cocycle-conjugated analogue of the untwisted \(\delta\)-naturality
  (\ref{lem:relDiffBaseChange}); on the pinned cover, \(\delta_g\) is on the nose the
  differential of the pair of the relative twisted sheaf, so the square also reads on
  the pairs.
\end{theorem}
\begin{proof}
  \leanok
  Both composites are additive, so compare them on \(r' \otimes (p_0, p_1)\). Read
  everything in chart-0 trivialised coordinates. By \ref{lem:twistTriv_pairDiff} the
  twisted differential is \((s_0, s_1) \mapsto s_0|_{\cap} - g \cdot s_1|_{\cap}\), so
  along the top and right the tensor goes to
  \[
    r' \cdot \rho^{\sharp}\bigl(s_0|_{\cap}\bigr)
      - r' \cdot \rho^{\sharp}(g)\,\rho^{\sharp}\bigl(s_1|_{\cap}\bigr),
  \]
  by \ref{def:relTwistDomBaseChange} and multiplicativity of \(\rho^{\sharp}\). Along
  the left and bottom it goes to
  \(\bigl(r' \cdot \rho^{\sharp}s_0\bigr)\big|_{\cap} - g' \cdot \bigl(r' \cdot
  \rho^{\sharp}s_1\bigr)\big|_{\cap}\), where \(g' = \rho^{\sharp}(g)\)
  (\ref{def:relCocycleBaseChange}). The two agree because \(\rho^{\sharp}\) commutes
  with restriction (\ref{lem:relSectionsMap_props}), the \(R'\)-action commutes with
  restriction, and the \(g'\)-factor commutes with the \(R'\)-scalar (commutative
  ring). The pair form on the pinned cover is the same identity, the pair's
  differential being definitionally the twisted restriction-difference map.
\end{proof}

\begin{theorem}[\(H^1\) of the twisted sheaf vanishes after base change]
  \label{thm:relTwist_h1_baseChange}
  \dcref{custom}
  \lean{AlgebraicGeometry.relTwist_surjective_diff_baseChange,
    AlgebraicGeometry.relTwist_subsingleton_h1_baseChange}
  \uses{thm:relTwistDiffBaseChange, lem:twoCoverPair_surjective_diff, lem:kerCongr,
    def:twoCoverPair_h0h1, lem:twoLattice_range_diff, def:relTwistPairData}
  \leanok
  On the vanishing locus (\(H^1\) of the pair of \(F_g\) over \(R\) trivial), for
  every homomorphism \(R \to R'\) of commutative \(k\)-algebras: the twisted
  differential over \(R'\) at the base-changed cocycle \(g'\) is surjective, and
  \[
    H^1\bigl(C_{R'},\, F_{g'}\bigr) \;=\; 0 .
  \]
\end{theorem}
\begin{proof}
  \leanok
  Over \(R\) the twisted differential is surjective
  (\ref{lem:twoCoverPair_surjective_diff}); tensoring with \(R'\) preserves
  surjectivity (right exactness), and surjectivity transports across the
  \(\delta\)-naturality square (\ref{thm:relTwistDiffBaseChange}) by
  \ref{lem:kerCongr}. So \(\delta_{g'}\) is surjective; hence \(H^1\) of the pair of
  \(F_{g'}\) --- its cokernel (\ref{lem:twoLattice_range_diff}) --- is trivial, and
  the degree-one carrier over \(R'\) (\ref{def:twoCoverPair_h0h1}, for the pair data
  of \(F_{g'}\), \ref{def:relTwistPairData}) transports the vanishing to
  \(H^1(C_{R'}, F_{g'})\).
\end{proof}

\begin{definition}[\(H^0\) base change on the nose]
  \label{def:relTwistH0BaseChange}
  \lean{AlgebraicGeometry.relTwistH0BaseChange}
  \uses{def:relTwistH0BaseChangeEquiv, lem:kerCongr, thm:relTwistDiffBaseChange,
    def:twoCoverPair_h0h1, def:relTwistPairData}
  \leanok
  \dcref{custom}
  On the vanishing locus, for every homomorphism \(R \to R'\) of commutative
  \(k\)-algebras there is an \(R'\)-linear equivalence
  \[
    R' \otimes_R H^0\bigl(C_R, F_g\bigr) \;\cong\;
      H^0\bigl(C_{R'},\, F_{g'}\bigr), \qquad g' = \rho^{\sharp}(g) :
  \]
  formation of the twisted \(H^0\) commutes on the nose with every change of the test
  ring --- the curve-lite form of Kleiman's clause (iii) at ring maps.
\end{definition}
\begin{proof}
  \leanok
  Compose three equivalences: the abstract ring-map clause
  \(R' \otimes_R H^0(C_R, F_g) \cong \ker(\delta_g \otimes_R R')\)
  (\ref{def:relTwistH0BaseChangeEquiv}); the kernel transport
  \(\ker(\delta_g \otimes_R R') \cong \ker(\delta_{g'})\) along the twisted
  \(\delta\)-naturality square (\ref{lem:kerCongr},
  \ref{thm:relTwistDiffBaseChange}); and the inverse of the degree-zero carrier of the
  pair of \(F_{g'}\) (\ref{def:twoCoverPair_h0h1}, \ref{def:relTwistPairData}),
  \(\ker(\delta_{g'}) \cong H^0(C_{R'}, F_{g'})\).
\end{proof}

\section{Base-field invariance of \(H^0\), \(H^1\) and the genus}
\label{sec:h1_base_field_invariance}

Fix a field \(k\) and an object \(C \in \Over(\Spec k)\). Section
\ref{sec:relative_h1_base_change} base-changed the two-cover complex along a
homomorphism \(R \to R'\) of test rings; this section runs the same mechanism once more
with the \emph{absolute} curve as source, for an arbitrary commutative \(k\)-algebra
\(R\): the two-term \v{C}ech complex of an affine two-cover of \(C\) base-changes
termwise by the free identification \ref{def:relSectionsBaseChange}, and \(H^1\), being
its cokernel, commutes with any base change by right-exactness. Specialised to a field
extension \(R = K\), this yields \(H^i(C, \struct C) \otimes_k K \cong
H^i(C_K, \struct{C_K})\) for \(i = 0, 1\) and the base-field invariance of the genus.
No flatness, semicontinuity or exchange argument is used.

\begin{lemma}[Pullback of sections commutes with restriction]
  \label{lem:relPullbackSection_resHom}
  \dcref{custom}
  \lean{AlgebraicGeometry.relPullbackSection_resHom}
  \uses{def:relSectionsBaseChange}
  \leanok
  Let \(R\) be a commutative \(k\)-algebra and \(W \leq V\) opens of \(C\). For a
  section \(s \in \Gamma(C, V)\), restricting on \(C\) and then pulling back along the
  first projection \(\mathrm{pr}_1 : C_R \to C\) agrees with pulling back and then
  restricting on the relative curve:
  \[
    \mathrm{pr}_1^{\sharp}\bigl(s|_W\bigr)
      = \bigl(\mathrm{pr}_1^{\sharp} s\bigr)\big|_{\mathrm{pr}_1^{-1} W}.
  \]
\end{lemma}
\begin{proof}
  \leanok
  Both sides are the image of \(s\) under the map of sections induced by
  \(\mathrm{pr}_1\) from \(V\) to \(\mathrm{pr}_1^{-1}W \subseteq \mathrm{pr}_1^{-1}V\):
  pullback along a morphism of schemes is compatible with restriction on source and
  target.
\end{proof}

\begin{definition}[Termwise base change of the two-cover complex of the curve]
  \label{def:curveDiffTermsBaseChange}
  \dcref{custom}
  \lean{AlgebraicGeometry.curveDiffDomBaseChange, AlgebraicGeometry.curveDiffCodBaseChange,
    AlgebraicGeometry.curveDiffCodBaseChange_tmul}
  \uses{def:relSectionsBaseChange, def:relCover, lem:relCover_inf, def:AffineTwoCover}
  \leanok
  Let \(D = (V_0, V_1)\) be an affine two-cover of \(C\) and \(R\) a commutative
  \(k\)-algebra. The two terms of the two-cover complex of \(C\) base-change freely by
  \ref{def:relSectionsBaseChange}: an \(R\)-linear equivalence
  \[
    R \otimes_k \bigl(\Gamma(C, V_0) \times \Gamma(C, V_1)\bigr)
    \;\cong\;
    \Gamma(C_R, V_0^R) \times \Gamma(C_R, V_1^R)
  \]
  on the degree-zero term (product of the two chart identifications; the charts are
  affine, hence qcqs), and
  \[
    R \otimes_k \Gamma(C, V_0 \cap V_1) \;\cong\; \Gamma(C_R, V_0^R \cap V_1^R)
  \]
  on the overlap term (the overlap of affines is qcqs, and
  \(V_0^R \cap V_1^R = \mathrm{pr}_1^{-1}(V_0 \cap V_1)\) by \ref{lem:relCover_inf}),
  where \(V_i^R = \mathrm{pr}_1^{-1} V_i\). On a pure tensor the overlap identification
  is \(a \otimes s \mapsto a \cdot \mathrm{pr}_1^\sharp(s)\).
\end{definition}

\begin{theorem}[The base-changed difference map is the relative difference map]
  \label{thm:curveDiffBaseChange}
  \dcref{custom}
  \lean{AlgebraicGeometry.curveDiffBaseChange, AlgebraicGeometry.curveDiffBaseChange_range}
  \uses{def:curveDiffTermsBaseChange, def:TwoCover_diff, lem:relPullbackSection_resHom}
  \leanok
  Under the two term equivalences of \ref{def:curveDiffTermsBaseChange}, the base
  change \(\mathrm{id}_R \otimes \mathrm{diff}\) of the two-cover
  restriction-difference map of \(C\) (\ref{def:TwoCover_diff}) is carried to the
  restriction-difference map of the relative curve \(C_R\) with its base-changed cover.
  In particular the overlap-term equivalence maps
  \(\mathrm{range}(\mathrm{id}_R \otimes \mathrm{diff})\) onto
  \(\mathrm{range}(\mathrm{diff}_R)\).
\end{theorem}
\begin{proof}
  \leanok
  Both sides are \(R\)-linear, so it suffices to check on a pure tensor
  \(a \otimes (s_0, s_1)\). The left leg gives
  \(a \cdot \mathrm{pr}_1^\sharp\bigl(s_0|_\cap - s_1|_\cap\bigr)\). The right leg
  first sends \(a \otimes (s_0, s_1)\) to
  \((a \cdot \mathrm{pr}_1^\sharp s_0,\, a \cdot \mathrm{pr}_1^\sharp s_1)\), then takes
  the difference of restrictions on \(C_R\). By
  \ref{lem:relPullbackSection_resHom}, \(\mathrm{pr}_1^\sharp(s_i|_\cap) =
  (\mathrm{pr}_1^\sharp s_i)|_{\cap_R}\); restriction on \(C_R\) is \(R\)-linear (the
  \(R\)-action is pullback of constants along \(\mathrm{pr}_2\), itself compatible with
  restriction), so the right leg gives
  \(a \cdot \bigl((\mathrm{pr}_1^\sharp s_0)|_{\cap_R} -
  (\mathrm{pr}_1^\sharp s_1)|_{\cap_R}\bigr)\), the same element. The range statement
  follows since both term equivalences are surjective.
\end{proof}

\begin{theorem}[Base change of the two-cover \(H^1\) cokernel of the curve]
  \label{thm:curveH1CokBaseChange}
  \dcref{custom}
  \lean{AlgebraicGeometry.curveH1CokBaseChange, AlgebraicGeometry.curveH1CokBaseChange_tmul_mk}
  \uses{thm:quot_range_baseChange, thm:curveDiffBaseChange, def:TwoCover_H1Cok}
  \leanok
  For every commutative \(k\)-algebra \(R\) there is an \(R\)-linear equivalence
  \[
    R \otimes_k \mathrm{H1Cok}(C; V_0, V_1) \;\cong\;
      \mathrm{H1Cok}(C_R; V_0^R, V_1^R),
  \]
  sending \(a \otimes [s]\) to \([a \cdot \mathrm{pr}_1^\sharp(s)]\). No flatness of
  \(R\) is needed.
\end{theorem}
\begin{proof}
  \leanok
  \(\mathrm{H1Cok}\) is the cokernel of \(\mathrm{diff}\) (\ref{def:TwoCover_H1Cok}),
  and base change commutes with cokernels (\ref{thm:quot_range_baseChange}):
  \(R \otimes_k \mathrm{H1Cok}(C) \cong (R \otimes_k \Gamma(C, V_0 \cap V_1)) /
  \mathrm{range}(\mathrm{id}_R \otimes \mathrm{diff})\). Now transport along the
  overlap-term equivalence of \ref{def:curveDiffTermsBaseChange}, which by
  \ref{thm:curveDiffBaseChange} carries the range of
  \(\mathrm{id}_R \otimes \mathrm{diff}\) onto the range of the relative difference
  map; the induced map of quotients is the displayed equivalence, and its effect on the
  class of a pure tensor is read off the overlap identification.
\end{proof}

\begin{theorem}[Cohomology commutes with base change in degree one, from the absolute
  curve]
  \label{thm:curveH1BaseChange}
  \dcref{custom}
  \lean{AlgebraicGeometry.curveH1BaseChange}
  \uses{thm:TwoCover_h1CokEquiv, thm:curveH1CokBaseChange, def:relTwoCoverH1,
    def:moduleKSheaf, def:HModule}
  \leanok
  Let \(D\) be an affine two-cover of \(C\) and \(R\) a commutative \(k\)-algebra.
  There is an \(R\)-linear equivalence
  \[
    R \otimes_k H^1\bigl(C, \struct{C}\bigr) \;\cong\;
      H^1\bigl(C_R, \struct{C_R}\bigr)
  \]
  on the cohomology carriers.
\end{theorem}
\begin{proof}
  \leanok
  Compose: the base change of the absolute cokernel bridge
  \(H^1(C, \struct C) \cong \mathrm{H1Cok}(C)\) (\ref{thm:TwoCover_h1CokEquiv},
  tensored with \(R\)); the cokernel base change
  \(R \otimes_k \mathrm{H1Cok}(C) \cong \mathrm{H1Cok}(C_R)\)
  (\ref{thm:curveH1CokBaseChange}); and the inverse of the relative cokernel bridge
  \(H^1(C_R, \struct{C_R}) \cong \mathrm{H1Cok}(C_R)\) (\ref{def:relTwoCoverH1}).
\end{proof}

\begin{definition}[A pinned affine two-cover of the challenge curve]
  \label{def:curveCover}
  \dcref{custom}
  \lean{AlgebraicGeometry.curveCover}
  \uses{thm:AffineTwoCover_nonempty_of_curve}
  \leanok
  For the challenge curve \(C\) (proper, smooth of relative dimension one,
  geometrically irreducible over \(k\)), fix once and for all an affine two-chart cover
  of \(C\), which exists by \ref{thm:AffineTwoCover_nonempty_of_curve}; all base-change
  equivalences below are built on this one cover.
\end{definition}

\begin{theorem}[Cohomology commutes with base change in degree zero]
  \label{thm:curveH0BaseChange}
  \dcref{custom}
  \lean{AlgebraicGeometry.curveH0BaseChange}
  \uses{def:moduleKSheafHZero, def:relSectionsBaseChange}
  \leanok
  Let \(C\) be the challenge curve and \(R\) a commutative \(k\)-algebra. There is an
  \(R\)-linear equivalence
  \[
    R \otimes_k H^0\bigl(C, \struct C\bigr) \;\cong\; H^0\bigl(C_R, \struct{C_R}\bigr).
  \]
\end{theorem}
\begin{proof}
  \leanok
  Both carriers are global sections (\ref{def:moduleKSheafHZero}, over \(k\) and over
  \(R\)). The total space \(\top \subseteq C\) is quasi-compact (\(C\) is proper over a
  field) and quasi-separated, so the free sections base change
  \ref{def:relSectionsBaseChange} applies at \(V = \top\):
  \(R \otimes_k \Gamma(C, \top) \cong \Gamma(C_R, \top)\). Conjugating by the two
  degree-zero identifications gives the claim.
\end{proof}

\begin{theorem}[\(H^1\) base-field invariance]
  \label{thm:h1BaseFieldEquiv}
  \lean{AlgebraicGeometry.h1BaseFieldEquiv}
  \uses{thm:curveH1BaseChange, def:curveCover, thm:baseChangeCurve_bundle}
  \leanok
  \dcref{custom}
  Let \(C\) be the challenge curve over \(k\) and \(K/k\) a field extension. Then
  \[
    H^1\bigl(C, \struct C\bigr) \otimes_k K \;\cong\; H^1\bigl(C_K, \struct{C_K}\bigr)
  \]
  as \(K\)-vector spaces, where \(C_K\) is the base-changed curve regarded over
  \(\Spec K\) (\ref{thm:baseChangeCurve_bundle}).
\end{theorem}
\begin{proof}
  \leanok
  Specialise \ref{thm:curveH1BaseChange} to the algebra \(R = K\) and the pinned cover
  of \ref{def:curveCover}; the relative curve \(C_K\) over \(\Spec K\) is the
  base-changed curve of \ref{thm:baseChangeCurve_bundle}.
\end{proof}

\begin{theorem}[\(H^0\) base-field invariance]
  \label{thm:h0BaseFieldEquiv}
  \dcref{custom}
  \lean{AlgebraicGeometry.h0BaseFieldEquiv}
  \uses{thm:curveH0BaseChange, thm:baseChangeCurve_bundle}
  \leanok
  In the setting of \ref{thm:h1BaseFieldEquiv},
  \(H^0(C, \struct C) \otimes_k K \cong H^0(C_K, \struct{C_K})\) as \(K\)-vector
  spaces.
\end{theorem}
\begin{proof}
  \leanok
  Specialise \ref{thm:curveH0BaseChange} to \(R = K\).
\end{proof}

\begin{theorem}[The scalar identities of base-field invariance]
  \label{thm:finrank_h1_baseField}
  \dcref{custom}
  \lean{AlgebraicGeometry.finrank_h1_baseField, AlgebraicGeometry.finrank_h0_baseField,
    AlgebraicGeometry.finrank_h1_baseField_eq_genus}
  \uses{thm:h1BaseFieldEquiv, thm:h0BaseFieldEquiv, def:genus}
  \leanok
  For the challenge curve \(C\) over \(k\) and any field extension \(K/k\),
  \[
    \dim_K H^1\bigl(C_K, \struct{C_K}\bigr) = \dim_k H^1\bigl(C, \struct C\bigr)
      = \genus(C),
    \qquad
    \dim_K H^0\bigl(C_K, \struct{C_K}\bigr) = \dim_k H^0\bigl(C, \struct C\bigr).
  \]
\end{theorem}
\begin{proof}
  \leanok
  Extension of scalars multiplies no dimensions:
  \(\dim_K(V \otimes_k K) = \dim_k V\) for any \(k\)-vector space \(V\). Apply this
  across the equivalences of \ref{thm:h1BaseFieldEquiv} and \ref{thm:h0BaseFieldEquiv};
  the \(\genus(C)\) identification is Definition~\ref{def:genus}.
\end{proof}

\begin{theorem}[Base-field invariance of the genus]
  \label{thm:genus_baseField}
  \dcref{custom}
  \lean{AlgebraicGeometry.genus_baseField}
  \uses{thm:finrank_h1_baseField, thm:baseChangeCurve_bundle, def:genus}
  \leanok
  Let \(C\) be the challenge curve over \(k\) and \(K/k\) a field extension. The
  base-changed curve \(C_K\), a legal challenge curve over \(K\)
  (\ref{thm:baseChangeCurve_bundle}), satisfies
  \[
    \genus(C_K) = \genus(C).
  \]
\end{theorem}
\begin{proof}
  \leanok
  By Definition~\ref{def:genus} applied over \(K\),
  \(\genus(C_K) = \dim_K H^1(C_K, \struct{C_K})\), which equals \(\genus(C)\) by
  \ref{thm:finrank_h1_baseField}.
\end{proof}
